% Worked solutions for Listas/lista-topic5-nonverifiable-income-termination.md (13 book exercises + originals 14-17).
% Compile: pdflatex twice; check: pdffonts (Type 1 only, uni = yes).
% Every number is asserted in lista-topic5-nonverifiable-income-termination-solutions_check.py.
\documentclass[11pt,a4paper]{article}
\usepackage[utf8]{inputenc}
\usepackage[T1]{fontenc}
\usepackage{lmodern}   % Type 1 vector fonts (replaces PK bitmaps)
\usepackage{cmap}      % ToUnicode CMap -> copyable text and accents
\usepackage[brazilian,english]{babel}
\usepackage{amsmath,amssymb,amsfonts}
\usepackage[a4paper,margin=2.2cm]{geometry}
\usepackage{booktabs}
\usepackage{enumitem}
\usepackage{fancyhdr}
\usepackage{xcolor}
\usepackage{microtype}

\pagestyle{fancy}\fancyhf{}
\lhead{\footnotesize Finan\c{c}as II --- Corporativas (MPE Insper)}
\rhead{\footnotesize Topic 5 list --- Worked solutions}
\cfoot{\footnotesize\thepage}
\renewcommand{\headrulewidth}{0.3pt}
\setlength{\headheight}{14pt}
\setlist{topsep=2pt,itemsep=1pt}

\newcommand{\exsec}[1]{\section*{\color{blue!55!black}#1}\addcontentsline{toc}{section}{#1}}
\newcommand{\notation}{\par\noindent\textbf{Notation.}}
\newcommand{\step}[2]{\par\medskip\noindent\textbf{Step #1 --- #2.}~}
\newcommand{\intuition}{\par\medskip\noindent\textbf{Intuition.}~}
\newcommand{\bookans}{\par\medskip\noindent\textbf{Book answer.}~}
\newcommand{\refbib}{\par\smallskip\noindent\textbf{Ref:}~}
\newcommand{\assume}{\par\smallskip\noindent\textbf{Assumption.}~}
\newcommand{\pt}{\foreignlanguage{brazilian}}
\newcommand{\asked}{\par\medskip\noindent\textbf{What is asked.}~}
\newcommand{\approach}{\par\smallskip\noindent\textbf{Approach.}~}
\newcommand{\answer}[1]{\par\smallskip\noindent\fbox{\textbf{Answer #1.}}~}
\newcommand{\pitfalls}{\par\medskip\noindent\textbf{Common mistakes.}}

\usepackage[hyphens]{url}
%% STATEMENT-PREAMBLE BEGIN (tools/statements_into_solutions.py)
\usepackage{longtable,array,calc}
\providecommand{\tightlist}{\setlength{\itemsep}{0pt}\setlength{\parskip}{0pt}}
\providecommand{\href}[2]{#2}
\providecommand{\real}[1]{#1}
\makeatletter\@ifundefined{c@none}{\newcounter{none}}{}\makeatother  % pandoc's \LTcaptype{none}
\newcommand{\stmtbegin}{\par\medskip\noindent\textbf{Statement.}\par\nopagebreak\begingroup\small\setlength{\tabcolsep}{3pt}\sloppy}
\newcommand{\stmtend}{\par\endgroup\nopagebreak\noindent\rule{\linewidth}{0.4pt}\par\medskip}
\DeclareUnicodeCharacter{2192}{\ensuremath{\rightarrow}}
\DeclareUnicodeCharacter{2212}{\ensuremath{-}}
\DeclareUnicodeCharacter{2022}{\textbullet{}}
\DeclareUnicodeCharacter{2248}{\ensuremath{\approx}}
\DeclareUnicodeCharacter{2032}{\ensuremath{^{\prime}}}
\DeclareUnicodeCharacter{2081}{\ensuremath{_{1}}}
\DeclareUnicodeCharacter{03B8}{\ensuremath{\theta}}
\DeclareUnicodeCharacter{03BB}{\ensuremath{\lambda}}
\DeclareUnicodeCharacter{03C8}{\ensuremath{\psi}}
\DeclareUnicodeCharacter{03B1}{\ensuremath{\alpha}}
\DeclareUnicodeCharacter{25A0}{\rule{0.7em}{0.7em}}
%% STATEMENT-PREAMBLE END
\begin{document}
\begin{center}{\Large\bfseries Exercise list --- Topic 5: Nonverifiable income, strategic default, termination, debt structure}\\[4pt]{\large Worked solutions}\end{center}

\noindent Statements are in \url{Listas/lista-topic5-nonverifiable-income-termination.md} (and its PDF); each exercise is cited by number and printed page and its full text is reproduced right before its solution. Exercises 14--17 are original exercises on the sovereign-debt model (Topic~5 slides 13--29), for which no course book has an exercise; they are marked ``Original exercise''. Short answers: \url{lista-topic5-nonverifiable-income-termination-solutions-short.pdf}. Editions are the local copies: Tirole, \emph{The Theory of Corporate Finance} (2005); Amaro de Matos, \emph{Theoretical Foundations of Corporate Finance} (2001); Berk = Berk, DeMarzo \& Harford, \emph{Fundamentals of Corporate Finance}, 2nd ed. (2011); Grinblatt--Titman, \emph{Financial Markets and Corporate Strategy}, 2nd ed. (2002); Ross = Ross, Westerfield, Jordan \& Lamb, \emph{Fundamentos de Administra\c{c}\~ao Financeira}, 9th ed. (2013).

\tableofcontents

% =============================================================================
\exsec{1. Control passing to creditors --- Grinblatt \& Titman 16.10, p.~593}
%% BEGIN-STATEMENT 1
\stmtbegin
\textbf{Source:} Grinblatt \& Titman, Ch. 16, Exercise 16.10, p.~593 (PDF 606) · Answer in book: no

16.10. In the event of bankruptcy, the control of a firm passes from the equity holders to the debt
holders. Describe differences in the preferences of the equity holders and debt holders and how
decisions following bankruptcy proceedings are likely to change.
\stmtend
%% END-STATEMENT

\asked The exercise has two parts. \emph{(preferences)}: explain how what shareholders want the firm
to do differs from what creditors want, given the claims each one holds. \emph{(decisions)}: explain
which corporate decisions (risk, liquidation, investment, payouts, management) are likely to change once
bankruptcy hands control of the firm to the debt holders.

\approach Each claimant's preferences follow from the shape of its payoff as a function of firm value:
equity is a call option on the firm (convex), debt is a short put plus a riskless bond (concave). Whoever
holds control maximises the expected value of \emph{its own} claim, not total value; this is the
debt holder--equity holder conflict of G\&T \S16.2. In the lecture's language, default is the event that
transfers control to creditors, which is what makes their claim enforceable when income is not
verifiable. A two-state numerical example makes each bias concrete.

\notation
\begin{itemize}
  \item $F>0$: face value of the debt due now; $V_L\geq 0$: value if the firm is liquidated now.
  \item $\tilde V$: random value of the firm if it keeps operating (going-concern value), $E[\tilde V]$ its mean; risk neutrality and a zero interest rate, as in the chapter's examples.
  \item $E(\cdot)$, $D(\cdot)$: payoffs of equity and debt for a given total value $V$.
\end{itemize}

\step{1}{Two claims, two payoff shapes} By absolute priority, at a value $V$ the claims pay
\begin{align}
  D(V) &= \min(V,F), \label{t5e1-debt}\\
  E(V) &= \max(V-F,0). \label{t5e1-equity}
\end{align}
Equity \eqref{t5e1-equity} is a call option on the firm: convex, it gains from upside and is
protected on the downside. Debt \eqref{t5e1-debt} is concave: it gains nothing above $F$ and bears the
losses below $F$. Each group, when in control, picks the action that maximises the expected value
of its own claim, not $E[\tilde V]$ (G\&T Result 16.2, p.~562).

\step{2}{Equity holders' preferences while they control the firm} Near distress, \eqref{t5e1-equity}
makes shareholders (and the managers who represent them) prefer:
(i) \emph{risk} --- asset substitution, even into negative-NPV gambles, because the downside falls on
creditors; (ii) \emph{continuation} over liquidation --- reluctance to liquidate, since liquidation
pays them $\max(V_L-F,0)=0$ when $F>V_L$, like exercising an out-of-the-money option;
(iii) \emph{no new equity-financed investment} --- debt overhang, since the existing debt captures
much of the payoff; (iv) \emph{quick payoffs} --- shortsighted investment (G\&T \S 16.2, pp.~562--578).
Illustration with the book's Eastern Airlines numbers (p.~576): $F=500$; continuing gives
$\tilde V=600$ or $200$ with probability $\tfrac12$ each; liquidating now gives $V_L=480$ (\$ millions).
Using \eqref{t5e1-equity} and \eqref{t5e1-debt},
\begin{align}
  E_{\text{cont}} &= \tfrac12\max(600-500,0)+\tfrac12\max(200-500,0) \label{t5e1-Econt}\\
                  &= \tfrac12(100)+\tfrac12(0) = 50, \nonumber\\
  E_{\text{liq}}  &= \max(480-500,0) = 0, \label{t5e1-Eliq}\\
  D_{\text{cont}} &= \tfrac12\min(600,500)+\tfrac12\min(200,500) \label{t5e1-Dcont}\\
                  &= \tfrac12(500)+\tfrac12(200) = 350, \nonumber\\
  D_{\text{liq}}  &= \min(480,500) = 480. \label{t5e1-Dliq}
\end{align}
Shareholders want to continue ($50>0$), creditors want to liquidate ($480>350$), and the firm is
worth $E[\tilde V]=\tfrac12(600)+\tfrac12(200)=400$ if it continues against $480$ if it liquidates:
equity control destroys 80.

\step{3}{Creditors' preferences once they take control} In bankruptcy (or after a covenant
default) control passes to the debt holders or their representatives (G\&T p.~576). By
\eqref{t5e1-debt} they are \emph{conservative}: they value the downside fully and the upside above
$F$ not at all. They therefore prefer (i) safe strategies and lower risk; (ii) liquidation or asset
sales whenever $V_L$ exceeds the expected value of their claim under continuation, which is
efficient in Step~2's example; (iii) cash collection and cost cutting over long-horizon, risky
investment. The bias can now run the other way: with $F=V_L=500$ and $\tilde V=800$ or $300$
(own numbers), continuing is worth $E[\tilde V]=\tfrac12(800)+\tfrac12(300)=550>500$, but
\begin{align}
  D_{\text{cont}} &= \tfrac12\min(800,500)+\tfrac12\min(300,500) \label{t5e1-Dcont2}\\
                  &= \tfrac12(500)+\tfrac12(300) = 400 < 500 = D_{\text{liq}}, \nonumber
\end{align}
so creditors liquidate a firm worth more alive (the Emruss case, pp.~577--578, has the same logic:
senior creditors liquidate because liquidation pays them in full).

\answer{(preferences)} Equity holders (convex claim $\max(V-F,0)$) favour risk, continuation over
liquidation, no new equity-financed investment and quick payoffs; debt holders (concave claim
$\min(V,F)$) favour safety, liquidation whenever $V_L$ exceeds the expected value of their claim under
continuation, and cash collection. In the example, equity wants to continue ($50>0$), debt wants to
liquidate ($480>350$).

\step{4}{How decisions change after bankruptcy} Relative to equity control: risk is cut, failing
divisions are closed or sold, liquidation becomes more likely (sometimes too likely), dividends stop,
and managers are often replaced. When Chapter~11 converts the debt into equity of the reorganised
firm, the former creditors become residual claimants: debt overhang and the incentive to gamble
disappear, and decisions move back toward maximising total value. Absolute-priority violations
(p.~560) leave shareholders some bargaining power, so the shift of control is only partial.

\answer{(decisions)} After bankruptcy the firm takes less risk, sells or closes failing assets,
liquidates more often (sometimes too often: with $\tilde V=800$ or $300$ and $V_L=F=500$, creditors
liquidate a firm worth $550>500$ alive), stops dividends and often replaces management; once the debt is
converted into equity, decisions move back toward maximising total value.

\intuition Each claimant behaves like the shape of its payoff: the convex claim gambles and delays
liquidation, the concave claim plays safe and liquidates early. Shifting control to creditors in
default is what makes their claim credible --- the same role the termination threat plays in
Topic~5 --- but it replaces one bias by the opposite one unless the claims are restructured.
\pitfalls
\begin{itemize}
  \item Saying that creditors in control maximise firm value: they maximise $E[\min(\tilde V,F)]$ and
        ignore the upside above $F$, so they can liquidate a firm worth more alive.
  \item Treating liquidation as always efficient (or always inefficient): it is efficient only when
        $V_L>E[\tilde V]$; the bias depends on who holds control.
  \item Forgetting absolute-priority violations and debt-for-equity swaps, which change the payoff
        shapes and hence the post-bankruptcy decisions.
\end{itemize}
\refbib{Grinblatt \& Titman Ch.~16, \S 16.1 ``Bankruptcy'', pp.~559--561; \S 16.2 ``Debt Holder--Equity Holder Conflicts'', pp.~561--578 (``The Reluctance to Liquidate Problem'', pp.~575--578); Ch.~18, p.~643 (creditor influence makes managers ``too conservative'').}

% =============================================================================
\exsec{2. Duration of debt: proving (5.11) --- Amaro de Matos 5.2.5 \#3, p.~148}
%% BEGIN-STATEMENT 2
\stmtbegin
\textbf{Source:} Amaro de Matos, Ch. 5, Exercise 5.2.5 \#3, p.~148 (PDF 166) · Answer in book: yes,
p.~268

\begin{enumerate}
\def\labelenumi{\arabic{enumi}.}
\setcounter{enumi}{2}
\tightlist
\item
  Prove equation (5.11).
\end{enumerate}

\paragraph{Context needed to solve it}

\emph{From Amaro de Matos §5.2.3 ``The Choice of the Duration'', pp.~139--141 (PDF 157--159):}

A model developed by Diamond (1991, 1993) illustrates in a very simple way the differences between
short-term and long-term debt financing. The model is based on information asymmetry and credit
rating.

Regaining the notation of Aghion and Bolton, consider an entrepreneur with a two-period project who
needs funds for an initial investment \(K = 1\). The project may be of one of two types: good or
bad. At time \(t = 2\) a good project has a payoff \(X > 0\), whereas a bad project pays that amount
with a probability \(\pi\), and pays zero with probability \(1 - \pi\). Let \(\theta\) denote the
type of project, \(\theta_g\) referring to a good project and \(\theta_b\) referring to a bad
project. The quality of the project is known to the entrepreneur but not to the investor. Given a
project, as of date \(t = 0\) the investor assigns a probability \(f\) that \(\theta = \theta_g\).
This probability is referred to as a \emph{credit rating} and reflects the lender's information
about the project. It follows that, as of time 0, the probability \(q\) that a borrower pays the
loan at date \(t = 2\) is given by

\[q \equiv f + \pi(1 - f) = \pi + f(1 - \pi). \qquad (5.9)\]

At time \(t = 1\) information about the project becomes available to the investor. This new
information is not verifiable, so it makes no sense to write contracts contingent upon that
information, since they cannot be enforced. However, the investors update their beliefs about the
quality of the project and the credit rating will be either upgraded or downgraded.

Under this revision of the credit rating, let \(f^u\) denote the probability that
\(\theta = \theta_g\) given an upgrade and let \(f^d\) denote the probability that
\(\theta = \theta_g\) given a downgrade. In that case, the probability that an entrepreneur who
receives a downgrade repays at time \(t = 2\) is

\[q^d = \pi + f^d(1 - \pi) > f^d. \qquad (5.10)\]

Moreover, it is here assumed that

\[f^d < f < f^u,\]

meaning that the probability that the project is good is greater given an upgrade than given a
downgrade. Furthermore, the simplifying assumption \(f^u = 1\) is also made. Thus, given an upgrade
at date \(t = 1\), it is certain that the project is good. However, it is not certain that all
projects of good quality are upgraded. There is room for one such project to be downgraded. Let
\(e\) denote the probability of such an event.

All projects are also assumed to produce a nonassignable control rent \(C\) if the management has
control at time \(t = 2\). Also, if at time \(t = 1\) the manager defaults on short-term debt,
lenders then have the right to remove the management from control, selling the assets for a
liquidation value \(L < K = 1\). All projects may be liquidated at that point in time, under those
conditions. If the project is liquidated, the entrepreneur is prevented from consumption of his/her
present or future cash flows and control rents. In this case, an optimal contract is a debt contract
enforced by the right to liquidate the project if the debt is not repaid in full.

Just as in the case of the borrowers, lenders are assumed be risk neutral. Moreover, lenders have
alternative investments at \(t = 0\) that provide a return \(R\) at \(t = 2\). The condition
\(L < X/R\) ensures that a good project yields a higher return when not liquidated.

A project may be financed with either long-term or short-term debt. In the context of this model,
long-term debt refers to a contract issued at \(t = 0\) with no refinancing at date \(t = 1\) and
maturing at \(t = 2\). If one such instrument provides \(K = 1\) at time \(t = 0\) with face value
\(\rho \le X\), the lender expects to receive at maturity \(\rho[\pi + f(1 - \pi)] = R^2\). Notice
that if \(\rho > X\), a borrower with credit rating \(f\) will not be able to provide the lender
with the required rate of return \(R^2\). Now, if the project is upgraded at time \(t = 1\), the
market value of the debt contract at that point in time becomes \(V_l^u = \rho/R\). Similarly, if
the project is downgraded, the market value of debt at time \(t = 1\) becomes \(V_l^d = q^d\rho/R\).

Short-term debt, on the other hand, is debt financed at \(t = 0\) with face value \(r_1\) and
maturity at time \(t = 1\). If short-term borrowers are upgraded at \(t = 1\), their debt contract
at that point in time is worth \(V_s^u = r_1\) and they may refinance their project by issuing a new
short-term debt contract with maturity at \(t = 2\) and face value \(r_1 R\) that will be paid with
certainty and guarantees a return of \(R\) per unit invested by the lender. If short-term borrowers
are downgraded, however, in principle they may refinance their debt with a new short-term contract
only if \(q^d X/R > r_1\), since \(X\) is the maximum that a borrower can pay at \(t = 2\). If that
is the case, at \(t = 1\) the market value of the short-term debt issued at \(t = 0\) would be
\(V_s^d = r_1\) if the borrower pays, possibly through the refinancement mechanism. But what happens
if \(q^d X/R < r_1\)? In that case there is still room for a refinancement provided that the
liquidation value \(L\) is less than the expected value to the lender from a refinancement,
\(q^d X/R\). Under a downgrading, the time \(t = 1\) value of the short-term debt is thus

\[V_s^d = \min\left[r_1, \max(L, q^d X/R)\right].\]

The face value of short-term debt is established such as to provide the lenders with an expected
return of \(R\) per period. It then follows that

\[r_1 = \frac{R - [1 - (1 - e)f]V_s^d}{(1 - e)f}. \qquad (5.11)\]
\stmtend
%% END-STATEMENT

\asked Derive the face value $r_1$ of one-period (short-term) debt in Diamond's (1991) rating model,
i.e.\ show that competitive lenders who must earn $R$ per period set $r_1$ as in (5.11).

\approach Short-term debt is priced by a one-period break-even condition: what the lenders expect to
hold at $t=1$, averaged over the upgrade and downgrade states, must equal $R$ per unit lent. The state
probabilities come from the rating model ($f^u=1$: only good projects are upgraded), and the claim's
$t=1$ value in each state is given in the text. This is the lecture's maturity question: short-term
debt is repriced at $t=1$ on the basis of nonverifiable (but observable) news.

\notation
\begin{itemize}
  \item $f\in(0,1)$: date-0 probability that the project is good (credit rating); $f^u=1$, $f^d<f$: posterior after an upgrade / a downgrade.
  \item $e\in(0,1)$: probability that a good project is downgraded at $t=1$; bad projects are always downgraded (since $f^u=1$).
  \item $r_1$: face value of one-period debt due at $t=1$; $V_s^u$, $V_s^d$: its $t=1$ value after an upgrade / a downgrade.
  \item $R>1$: gross one-period return of the lenders' alternative; $L$: liquidation value; $X$: payoff of a successful project; $\pi\in(0,1)$; $q^d=\pi+f^d(1-\pi)$: repayment probability after a downgrade (5.10).
  \item $P_u$, $P_d$: probabilities of an upgrade and a downgrade for a borrower of unknown type; $M\equiv\max(L,q^dX/R)$.
\end{itemize}

\step{1}{Probabilities of upgrade and downgrade} We need the weights of the two $t=1$ states. By the
law of total probability over the two types, with $\Pr(\text{up}\mid\text{good})=1-e$ and
$\Pr(\text{up}\mid\text{bad})=0$ (only good projects are upgraded, since $f^u=1$):
\begin{align}
  P_u &= f\,(1-e)+(1-f)\,(0) \nonumber\\
      &= f(1-e), \label{t5e2-Pu}\\
  P_d &= 1-P_u = 1-(1-e)f. \label{t5e2-Pd}
\end{align}
The downgrade posterior is consistent with these by Bayes' rule, since a downgrade comes from a good
project with probability $fe$ and from a bad one with probability $(1-f)(1)$:
\begin{align}
  f^d = \frac{fe}{fe+(1-f)} = \frac{fe}{P_d}. \label{t5e2-bayes}
\end{align}

\step{2}{Value of the claim at $t=1$} After an upgrade the borrower is known to be good and
refinances $r_1$ with new debt of face value $r_1R$, repaid with certainty, so
\begin{align}
  V_s^u = r_1. \label{t5e2-Vu}
\end{align}
After a downgrade the value is the one stated in the text, $V_s^d=\min[r_1,\max(L,q^dX/R)]=\min(r_1,M)$.

\step{3}{Lenders' break-even} Lenders give $K=1$ at $t=0$ and require an expected value $R$ at
$t=1$ (one period at the per-period return $R$):
\begin{align}
  P_uV_s^u+P_dV_s^d &= R. \label{t5e2-be}
\end{align}
Substituting \eqref{t5e2-Pu}, \eqref{t5e2-Pd} and \eqref{t5e2-Vu} into \eqref{t5e2-be}:
\begin{align}
  \bigl(f(1-e)\bigr)(r_1)+\bigl(1-(1-e)f\bigr)V_s^d &= R, \nonumber\\
  (1-e)f\,r_1 &= R-[1-(1-e)f]V_s^d, \nonumber\\
  r_1 &= \frac{R-[1-(1-e)f]V_s^d}{(1-e)f}, \label{t5e2-r1}
\end{align}
where the second line subtracts $[1-(1-e)f]V_s^d$ from both sides and the third divides by
$(1-e)f>0$. This is (5.11). $\blacksquare$

\step{4}{Solving the fixed point} Since $V_s^d$ depends on $r_1$, \eqref{t5e2-r1} is an equation in
$r_1$; we solve it case by case to see when the debt is risky. Case A, $r_1\leq M$: then $V_s^d=r_1$ and
\eqref{t5e2-r1} gives
\begin{align}
  (1-e)f\,r_1+[1-(1-e)f](r_1) &= R, \nonumber\\
  r_1\bigl[(1-e)f+1-(1-e)f\bigr] &= R, \nonumber\\
  r_1 &= R, \label{t5e2-caseA}
\end{align}
consistent if and only if $R\leq M$: the debt is riskless. Case B, $R>M$: then $V_s^d=M$ and
\begin{align}
  r_1 = \frac{R-[1-(1-e)f](M)}{(1-e)f}. \label{t5e2-caseB}
\end{align}
It is consistent ($r_1>M$) because
\begin{align}
  r_1-M &= \frac{R-[1-(1-e)f]M-(1-e)f\,M}{(1-e)f} \nonumber\\
        &= \frac{R-M}{(1-e)f}>0, \nonumber
\end{align}
and $r_1>R$ because $r_1-R=[1-(1-e)f](R-M)/[(1-e)f]>0$ (subtract $R=(1-e)fR/[(1-e)f]$ from
\eqref{t5e2-caseB}). The premium $r_1-R$ pays for the shortfall $R-M$ that lenders bear in the
downgrade state.

\answer{(5.11)} $r_1=\dfrac{R-[1-(1-e)f]\,V_s^d}{(1-e)f}$, from the break-even
$f(1-e)\,r_1+[1-(1-e)f]\,V_s^d=R$. Since $V_s^d=\min(r_1,M)$, it gives $r_1=R$ (riskless debt) when
$R\le M$ and $r_1=[R-(1-(1-e)f)M]/[(1-e)f]>R$ otherwise.

\intuition Short-term lenders are paid in full whenever the borrower is upgraded and receive at
most the refinancing or liquidation value after a downgrade; the face value is set so that the
probability-weighted average of the two equals the safe return. Because the downgrade state carries
the loss, $r_1$ rises with the probability $1-(1-e)f$ of a downgrade.
\pitfalls
\begin{itemize}
  \item Requiring $R^2$ instead of $R$: short-term debt runs one period, so the lenders' target at
        $t=1$ is $R$ (long-term debt is the one priced against $R^2$).
  \item Writing the upgrade probability as $f$ or $1-e$: it is $f(1-e)$, because bad projects are never
        upgraded ($f^u=1$) and good ones are upgraded only with probability $1-e$.
  \item Treating $V_s^d$ as a constant: it is $\min[r_1,\max(L,q^dX/R)]$, so (5.11) is a fixed point and
        must be checked for consistency.
\end{itemize}
\bookans Matches the book's answer (p.~268), which writes the break-even condition
$P_uV_s^u+P_dV_s^d=R$ with $P_u=f(1-e)$ and solves for $r_1$. Step~4 is an addition.
\refbib{Amaro de Matos Ch.~5, \S 5.2.3 ``The Choice of the Duration'', pp.~139--141; Appendix C.5.2, p.~268.}

% =============================================================================
\exsec{3. Threat of termination: continuation probability --- Amaro de Matos 5.2.5 \#1, p.~148}
%% BEGIN-STATEMENT 3
\stmtbegin
\textbf{Source:} Amaro de Matos, Ch. 5, Exercise 5.2.5 \#1, p.~148 (PDF 166) · Answer in book: yes,
p.~267

\begin{enumerate}
\def\labelenumi{\arabic{enumi}.}
\tightlist
\item
  Show that in the optimal debt contract, the probability \(\beta_\circ\) defined in (5.8) decreases
  as \(q\) increases as stated in the text.
\end{enumerate}

\paragraph{Context needed to solve it}

\emph{From Amaro de Matos §5.2.1 ``Extending the Model'', pp.~134--137 (PDF 152--155):}

As before, all the agents are risk neutral and the risk-free rate is assumed to be zero. The time
schedule refers to a two-period model and is characterized as follows: At time \(t = 0\) an amount
\(K\) is invested in the project. The entrepreneur, as before, is assumed to have no initial wealth.
In this specification of the model, the state of nature is characterized by the cash flow at
\(t = 1\). It may be \(\theta_g = x\) with probability \(q\) and \(\theta_b = 0\) with probability
\(1 - q\). The cash flow is observable for both the investor and the entrepreneur, but not for any
third party, and is thus not verifiable. Therefore, as before, a contract between the investor and
the entrepreneur is incomplete.

The manager is assumed to have the ability to divert cash flow without the control of the investor.
In that case a zero payoff would be declared at \(t = 1\) when the actual payoff was \(x\). If the
declared cash flow at \(t = 1\) is \(x\), the contract states a payment of \(R_x\) to the
creditholder and, with probability \(\beta_x\), the right of the creditholder to liquidate a part of
the assets that is worth \(L_x\). If the cash flow is null, the values are, respectively,
\(R_\circ\), \(\beta_\circ\), and \(L_\circ\).

Notice that taking \(\beta_x = 0\), \(\beta_\circ = 1\), and \(L_x = L_\circ\), being equal to the
value of the assets of the project at time \(t = 1\), this reduces to a usual debt contract with
face value \(R_x < x\) and one-period maturity. The issue in designing a contract is to find the
optimal values for the preceding parameters. Therefore, it is possible in this simple setting to
analyze the conditions under which a debt contract would result as optimal.

In the second period, two things may happen. Either the manager continues the project or the project
is liquidated, in the sense that it is separated from the manager. If the manager's activity
continues, the cash flow at time \(t = 2\) is \(y\). In the case of no liquidation, the project thus
has a net present value of \(qx + y - K\).

The expected return to the firm is then written as

\[\pi_f = q[x - R_x + (1 - \beta_x)y] + (1 - q)[-R_\circ + (1 - \beta_\circ)y]. \qquad (5.3)\]

On the other hand, the investor expects to get

\[\pi_i = q[R_x + \beta_x L_x] + (1 - q)[R_\circ + \beta_\circ L_\circ] - K. \qquad (5.4)\]

Suppose for a moment that at time \(t = 1\) the cash flow is \(x\). The expected return to the firm
would then be \(x - R_x + (1 - \beta_x)y\), decomposed as follows: The cash flow \(x\) goes to the
entrepreneur, who has to pay \(R_x\) to the investor. With probability \(\beta_x\) the project is
liquidated at date \(t = 1\), leading to a continuation value of zero, and with probability
\(1 - \beta_x\) the project continues until time \(t = 2\), generating a cash flow \(y\) at that
point. If, however, the manager hides from the outer investor the cash flow \(x\), then only
\(R_\circ\) will be paid at time \(t = 1\), and liquidation occurs with probability \(\beta_\circ\),
leaving to the firm a possible payoff \(S\), not necessarily zero, as before. This happens because
the liquidation value \(L_\circ\) does not correspond to the true value of the assets in that state
of nature. In that case the expected return to the firm is
\(x - R_\circ + \beta_\circ S + (1 - \beta_\circ)y\).

For the manager to be induced to pay the due \(R_x\) when the cash flow is \(x\), one must therefore
impose

\[x - R_x + (1 - \beta_x)y \geq x - R_\circ + \beta_\circ S + (1 - \beta_\circ)y. \qquad (5.5)\]

Noticing that \(R_\circ = \beta_x = 0\) is optimal in the sense that these values maximize the
profit to the firm and facilitates the verification of (5.5), the preceding inequality reads

\[R_x \leq \beta_\circ(y - S).\]

Among all these possible values for \(R_x\) that ensure a trustworthy behavior from the manager of
the project, any investor would maximize its expected return by choosing the greatest value. Thus,
at the optimum, this restriction is binding. Substituting the optimal values of the parameters in
expression (5.4) for the expected result of the investor, one gets

\[\pi_i = \beta_\circ[q(y - S) + (1 - q)L_\circ] - K. \qquad (5.6)\]

Using the optimal values of the parameters, namely those minimizing the profit to the investor, in
the profit of the firm (5.3), it follows that

\[\pi_f = (qx + y - K) - (1 - q)\beta_\circ(y - L_\circ). \qquad (5.7)\]

This result can be interpreted as the net present value of the project in the absence of liquidation
minus a cost due to contract incompleteness.

For the profit in (5.6) to be positive, ensuring the minimum condition of feasibility of the
contract, one must have

\[1 \geq \beta_\circ \geq \frac{K}{q(y - S) + (1 - q)L_\circ},\]

where the first inequality comes from the fact that \(\beta_\circ\) is a probability and the second
inequality comes from the condition \(\pi_i \geq 0\). Notice that \(\pi_f\) is decreasing in
\(\beta_\circ\) whereas \(\pi_i\) is monotonically increasing in that parameter. Maximizing the
expected return to the firm conditioned to a positive expected return to the investor defines the
optimal value of \(\beta_\circ\) as the lowest possible. Thus

\[\beta_\circ = \frac{K}{q(y - S) + (1 - q)L_\circ}.\]

Both \(S\) and \(L_\circ\) (and thus also \(\beta_\circ\)) are still to be specified. There is no
need to specify \(L_x\), since the main feature of the optimal contract just described is that there
is no liquidation in the good state. Liquidation occurs only, if at all, in the case of a poor
performance. In order to understand what determines \(S\) and \(L_\circ\) some additional
assumptions are required.

It is assumed that the entrepreneur is not able to continue managing the project in the case of
bankruptcy after a zero cash flow. This is the case of a \emph{liquidity default}. In that case, the
creditor may sell the position to a third party who obtains at \(t = 2\) a cash flow \(\alpha y\),
with \(\alpha \in [0, 1]\), meaning that the incoming manager cannot get better results out of the
assets than the manager who left the firm.

In the case of bankruptcy after a nonzero cash flow, the original manager will negotiate the buyback
of the firm with the creditholder, sharing the second-period return \(y\). In that situation,
\(S = y/2\). This is the case of a \emph{strategic default}.

It is also assumed that, in the case of liquidity default, the buyer of the assets incurs an
informational cost \(c \in [0, \bar c]\). This cost is modeled as a uniformly distributed random
variable. Therefore, by a negotiating process similar to the preceding, the remaining value of the
assets, \(\alpha y\), is equally divided between the creditor and the buyer. The buyer then
negotiates only if \(c \leq \alpha y/2\). Thus, assuming that \(\alpha y/2 \leq \bar c\), the
probability that this happens is

\[\Pr\left(c \leq \frac{\alpha y}{2}\right) = \frac{\alpha y}{2\bar c}.\]

Therefore, the expected value of the liquidation payoff in the case of a zero cash flow at \(t = 1\)
is

\[L_\circ = \frac{\alpha^2 y^2}{4\bar c},\]

and the optimal \(\beta_\circ\) reads

\[\beta_\circ = \frac{K}{q\frac{y}{2} + (1 - q)\frac{\alpha^2 y^2}{4\bar c}}. \qquad (5.8)\]

It follows\textsuperscript{3} that the optimal value of \(\beta_\circ\) decreases as \(q\)
increases. Thus, the probability of liquidation of assets decreases as the probability of positive
cash flow increases. Thus, the inefficiency due to contract incompleteness, expressed in (5.7) as
\((1 - q)\beta_\circ(y - L_\circ)\), is lower for higher \(q\), or, in other words, with lower risk
of bankruptcy.

\emph{From Amaro de Matos, Notes to Chapter Five, note 3:}

\begin{enumerate}
\def\labelenumi{\arabic{enumi}.}
\setcounter{enumi}{2}
\tightlist
\item
  This is left as an exercise.
\end{enumerate}
\stmtend
%% END-STATEMENT

\asked Show that the optimal liquidation probability $\beta_\circ$ after a declared zero cash flow,
given by (5.8), falls when the probability $q$ of the good cash flow rises (i.e.\ sign the derivative
$\partial\beta_\circ/\partial q$).

\approach With nonverifiable income, repayment is enforced only by the \emph{threat of termination}:
after a zero report the investor liquidates with probability $\beta_\circ$, set at the lowest level that
lets her break even. The denominator of (5.8) is her expected recovery per unit of $\beta_\circ$, a
$q$-weighted average of the recovery after a \emph{strategic default} ($y/2$) and after a \emph{liquidity
default} ($L_\circ$). Differentiating the ratio and signing $y/2-L_\circ$ answers the question.

\notation
\begin{itemize}
  \item $q\in(0,1)$: probability of the cash flow $x$ at $t=1$; $y>0$: second-period cash flow; $K>0$: investment.
  \item $\beta_\circ\in[0,1]$: probability of liquidation after a declared zero cash flow (the continuation probability is $1-\beta_\circ$).
  \item $S=y/2$: entrepreneur's share in a strategic-default buyback, so the investor gets $y-S=y/2$.
  \item $\alpha\in[0,1]$: efficiency of an outside buyer; $c\sim U[0,\bar c]$: buyer's information cost, with $\alpha y/2\leq\bar c$; $L_\circ=\alpha^2y^2/(4\bar c)$: expected liquidation payoff.
  \item $\Pi(q)$: investor's expected recovery per unit of $\beta_\circ$, the denominator of (5.8).
\end{itemize}

\step{1}{Write $\beta_\circ$ as a ratio} From (5.8),
\begin{align}
  \beta_\circ(q) &= \frac{K}{\Pi(q)}, \label{t5e3-beta}\\
  \Pi(q) &= q\,\frac{y}{2}+(1-q)\frac{\alpha^2y^2}{4\bar c}. \label{t5e3-Pi}
\end{align}

\step{2}{Differentiate} The sign of the comparative static is the sign of a derivative, so we compute
it. Writing \eqref{t5e3-beta} as $\beta_\circ=K\,\Pi(q)^{-1}$ and applying the chain rule ($K$, $y$,
$\alpha$, $\bar c$ do not depend on $q$):
\begin{align}
  \frac{\partial\beta_\circ}{\partial q} &= K\cdot(-1)\,\Pi(q)^{-2}\cdot\Pi'(q) \nonumber\\
  &= -\frac{K\,\Pi'(q)}{\Pi(q)^2}. \label{t5e3-dbeta}
\end{align}
\eqref{t5e3-Pi} is linear in $q$; $\frac{d}{dq}\bigl[q\,\tfrac y2\bigr]=\tfrac y2$ and
$\frac{d}{dq}\bigl[(1-q)\tfrac{\alpha^2y^2}{4\bar c}\bigr]=-\tfrac{\alpha^2y^2}{4\bar c}$, so
\begin{align}
  \Pi'(q) &= \frac{y}{2}-\frac{\alpha^2y^2}{4\bar c}. \label{t5e3-dPi}
\end{align}
Substituting \eqref{t5e3-dPi} and \eqref{t5e3-Pi} into \eqref{t5e3-dbeta}:
\begin{align}
  \frac{\partial\beta_\circ}{\partial q} &= \frac{-K\left(\dfrac{y}{2}-\dfrac{\alpha^2y^2}{4\bar c}\right)}{\left[q\,\dfrac{y}{2}+(1-q)\dfrac{\alpha^2y^2}{4\bar c}\right]^2}. \label{t5e3-final}
\end{align}

\step{3}{Sign} The denominator of \eqref{t5e3-final} is positive and $K>0$, so the sign is the
opposite of that of $\Pi'(q)$. Factor $y/2$ out of \eqref{t5e3-dPi}:
\begin{align}
  \Pi'(q) &= \frac{y}{2}\left[1-\frac{\alpha^2y^2}{4\bar c}\cdot\frac{2}{y}\right] \nonumber\\
          &= \frac{y}{2}\left[1-\frac{\alpha^2y}{2\bar c}\right] \nonumber\\
          &= \frac{y}{2}\left[1-\alpha\cdot\frac{\alpha y}{2\bar c}\right]. \label{t5e3-factor}
\end{align}
The assumption $\alpha y/2\leq\bar c$ gives $\alpha y/(2\bar c)\leq 1$, and $\alpha\leq1$, so the
bracket is $\geq0$; it is $>0$ unless $\alpha=1$ and $y/2=\bar c$. Hence $\Pi'(q)\geq0$ and
\begin{align}
  \frac{\partial\beta_\circ}{\partial q}\leq 0, \label{t5e3-sign}
\end{align}
strictly except in that knife-edge case, where $L_\circ=y/2$ and $\beta_\circ=K/(y/2)$ does not
depend on $q$. $\blacksquare$ (Interior solution assumed: $\Pi(q)\geq K$, so that $\beta_\circ\leq1$.)
The inefficiency $(1-q)\beta_\circ(y-L_\circ)$ in (5.7) is then also decreasing in $q$. By the
product rule, with $L_\circ$ independent of $q$ and $y-L_\circ>0$ (because $L_\circ\le y/2$),
\begin{align}
  \frac{\partial}{\partial q}\Bigl[(1-q)\beta_\circ(y-L_\circ)\Bigr]
  &= (y-L_\circ)\Bigl[-\beta_\circ+(1-q)\frac{\partial\beta_\circ}{\partial q}\Bigr]<0, \label{t5e3-ineff}
\end{align}
since $\beta_\circ>0$ and $\partial\beta_\circ/\partial q\le0$ by \eqref{t5e3-sign}.

\answer{(sign)} $\dfrac{\partial\beta_\circ}{\partial q}
=-\dfrac{K\left(\frac y2-\frac{\alpha^2y^2}{4\bar c}\right)}{\left[q\frac y2+(1-q)\frac{\alpha^2y^2}{4\bar c}\right]^2}\le0$
(strictly negative unless $\alpha=1$ and $y/2=\bar c$); the incompleteness cost $(1-q)\beta_\circ(y-L_\circ)$
also falls with $q$.

\intuition The investor recovers $y/2$ when she threatens a firm that defaulted strategically
(it had $x$ and buys itself back) but only $L_\circ\leq y/2$ when the default is a genuine liquidity
default (an outside buyer has to be found). A higher $q$ puts more weight on the first, richer
recovery, so a smaller liquidation probability suffices to break even --- as in the lecture's
$1-\gamma_0=(I-A)/[pE(\tilde R_2)+(1-p)L]$ (slide 11).
\pitfalls
\begin{itemize}
  \item Letting $L_\circ$ depend on $q$: it is fixed by the buyer's bargaining and information cost,
        $L_\circ=\alpha^2y^2/(4\bar c)$, so only the weights $q$ and $1-q$ move.
  \item Signing the derivative without the assumption $\alpha y/2\le\bar c$: it is what makes
        $\Pr(c\le\alpha y/2)=\alpha y/(2\bar c)$ a probability and gives $L_\circ\le y/2$.
  \item Claiming a strict decrease always: at $\alpha=1$, $y/2=\bar c$ the two recoveries coincide and
        $\beta_\circ$ does not depend on $q$.
\end{itemize}
\bookans Matches the book's answer (p.~267), which gives \eqref{t5e3-final} with ``$\leq0$''. The
book's argument assumes $\alpha y/2<\bar c$ (strict) and writes the step as ``$\alpha^2y/2<\bar c$'';
with the text's weak inequality $\alpha y/2\leq\bar c$ the derivative can be zero in the knife-edge
case of Step~3.
\refbib{Amaro de Matos Ch.~5, \S 5.2.1 ``Extending the Model'', pp.~134--137, eq.~(5.8); Appendix C.5.2, p.~267.}

% =============================================================================
\exsec{4. Entrepreneur control and renegotiation --- Amaro de Matos 5.1.6 \#3, p.~134}
%% BEGIN-STATEMENT 4
\stmtbegin
\textbf{Source:} Amaro de Matos, Ch. 5, Exercise 5.1.6 \#3, p.~134 (PDF 152) · Answer in book: yes,
p.~266

\begin{enumerate}
\def\labelenumi{\arabic{enumi}.}
\setcounter{enumi}{2}
\tightlist
\item
  Under entrepreneur control, show that the renegotiation of the contract whenever \(t_s < \hat t\)
  for \(s = 0, 1\) leads to an increase in the expected payoff of the entrepreneur.
\end{enumerate}

\paragraph{Context needed to solve it}

\emph{From Amaro de Matos §5.1.1 ``The Model'', pp.~127--129 (PDF 145--147):}

The relationship between incomplete financial contracts and the allocation of control in a firm may
be discussed within a simple analytical context, as in the paper of Aghion and Bolton (1992). This
section discusses their model, where a bilateral contract is considered between an entrepreneur
having no initial wealth and an investor who has the resources to invest the amount \(K\) at
\(t = 0\). This is a model with three points in time characterizing the following dynamic setting:

At \(t = 1\), the state of nature \(\theta\) is revealed. This state may be either good
(\(\theta_g\)) with probability \(q\) or bad (\(\theta_b\)) with probability \(1 - q\). Finally, at
\(t = 2\), the result of the investment is obtained, expressed as a rate of return \(r\). After the
realization of the state of nature at \(t = 1\) and before the maturity of the contract at
\(t = 2\), an action \(a\) is undertaken by one of the agents, either the investor or the
entrepreneur. The decision about who takes the action is to be contracted at \(t = 0\). The agent
who chooses the action is said to be the one who owns the project or who retains control.

The action \(a\) and the state of nature \(\theta\) are determining elements of the stochastic
return \(r\). An action is said to be optimal if it maximizes the expected wealth conditional to a
given state of nature. Let \(a_g\) denote the optimal action under the state of nature \(\theta_g\)
and let \(a_b\) denote the optimal action under the state of nature \(\theta_b\).

The model follows the idea of Grossman and Hart (1986) that incomplete contracts are characterized
by the fact that the state of nature is impossible or very costly to describe ex-ante and,
therefore, the ex-ante contract cannot be contingent on the state \(\theta\). Unlike Grossman and
Hart (1986), however, it is assumed here that at \(t = 1\) there is a publicly verifiable signal
\(s \in \{0, 1\}\) that is correlated with the state of nature \(\theta\) and that the contract may
be contingent on \(s\). One then defines the two conditional probabilities

\[\beta_g = \text{Prob}(s = 1 \mid \theta = \theta_g) > 1/2\]
\[\beta_b = \text{Prob}(s = 1 \mid \theta = \theta_b) < 1/2.\]

Note that \(\beta_g = 1\) and \(\beta_b = 0\) imply that the public signal \(s\) is a perfect proxy
for \(\theta\). In other words, one may infer \(\theta\) directly from the observable signal \(s\),
and thus the contract is no longer incomplete. For a contract to be incomplete, either \(\beta_g\)
is at a certain distance from 1 or \(\beta_b\) is away from zero, or both situations occur. One may
then use the sum of these two distances, namely,

\[d = |\beta_g - 1| + \beta_b, \qquad (5.1)\]

as a measure of the degree of incompleteness of the contract.

The entrepreneur (E) and the investor (I) are assumed to be risk neutral in income and to have
utility functions given, respectively, by

\[U_E(r, a) = g(r) + l(a, \theta)\] \[U_I(r, a) = r - g(r),\]

where \(l(a, \theta)\) is part of the total payoff of the entrepreneur that, although depending on
the action undertaken and on the state of nature, is nonmonetary and not observable by third
parties. Therefore, this term is taken as reflecting the \emph{private benefits} for the
entrepreneur.

It is also assumed that the stochastic return at \(t = 2\) may take only two values: \(r = 1\) or
\(r = 0\). As described previously, the stochastic return depends on both the state of nature
\(\theta\) and on the undertaken action \(a\). Define \(y^i_j\) as

\[y^i_j = E\left[r \mid \theta = \theta_i;\ a = a_j\right].\]

In other words, \(y^i_j\) is the expected rate of return given that the state of nature is
\(\theta_i\) and the action undertaken is \(a_j\). Under the same state of nature and action, let
the private benefit to the entrepreneur be denoted by \(l^i_j\). Then, for the project to be
feasible from the point of view of the investor, one must have a positive net present value for the
project,

\[q y^g_g + (1 - q) y^b_b > K.\]

Also, the optimality of \(a_g\) and \(a_b\) is assumed to be expressed as

\[y^g_g + l^g_g > y^g_b + l^g_b\] \[y^b_b + l^b_b > y^b_g + l^b_g.\]

Finally, it is assumed that the entrepreneur is compensated with a positive monetary transfer
schedule \(g(r) = t(s, r)\), which is a function of the first-time-period public signal, revealing
information about the state of nature and also the function of the return of the project at the end
of the second period. Given that \(r\) can take only the values 0 and 1, with no loss of generality,
this implies that the schedule \(t(s, r)\) is composed of a fixed part and a variable part,
proportional to \(r\):

\[t(s, r) = t_s r + t'_s \geq 0, \qquad (5.2)\]

where \(t_s\) and \(t'_s\) are constants. All residual returns go to the investor. Hence, if each
one of the agents could choose the action to be taken maximizing his/her utilitiy, then the
entrepreneur would maximize \(t_s y^i_j + t'_s + l^i_j\) and the investor would maximize
\((1 - t_s) y^i_j - t'_s\). The problem to be discussed within this model is the following: How does
the structure presented induce the agents to settle an ex-ante agreement on who should decide which
action should be taken after the state of nature is revealed? Is there room for renegotiation of the
contract once the state of the world is revealed, but before the action is undertaken?

\emph{From Amaro de Matos §5.1.2 ``Entrepreneur Control'', pp.~129--131 (PDF 147--149):}

In this section, the case where the entrepreneur has full control is studied. This means that the
entrepreneur chooses the action \(a\). The action will therefore be chosen so as to maximize the
entrepreneur's reward and will be given by

\[a^E(\theta_i, s) = \arg\max_{a_j}\left\{t_s y^i_j + t'_s + l^i_j\right\},\]

for \(i = g, b\) and \(j = g, b\). In order to describe the nature of the different contracts that
may arise in this context, we say that a contract is \emph{efficient} if the action chosen is the
optimal action for the realized state of nature. Also, we say that a contract is \emph{feasible} if
the investor expects to obtain at least as much as the amount invested, \(K\).

First, consider the case where the private benefits (\(l\)) satisfy the inequalities
\(l^g_g > l^g_b\) and \(l^b_b > l^b_g\). In that case, it is said that the private benefits (\(l\))
are comonotonic\textsuperscript{1} with the total revenues (\(y + l\)). It is fairly simple to show
that, in this situation, a contract with a schedule \(t(s, r)\) given by a constant \(t\),
independent of the signal perceived and of the return obtained, is efficient. Under the assumption
of comonotonic private benefits,

\[t + l^g_g > t + l^g_b\] \[t + l^b_b > t + l^b_g,\]

thus ensuring that the revelation of a good state (\(i = g\)) implies \(a^E = a_g\) and the
revelation of a bad state (\(i = b\)) implies \(a^E = a_b\). In other words, the optimal choice of
the entrepreneur coincides with the optimal choice for the project and one may say that entrepreneur
control is efficient.

In this case, for the contract also to be feasible, the constant fee schedule \(t \geq 0\) must be
chosen such that \(q y^g_g + (1 - q) y^b_b - t = K\). One such constant always exists given the
assumption \(q y^g_g + (1 - q) y^b_b > K\) and it follows

\textbf{Proposition 5.1.1 (Aghion and Bolton).} \emph{If \(l^g_g > l^g_b\) and \(l^b_b > l^b_g\),
then entrepreneur control is always efficient and feasible.}

A different treatment is required when the assumptions of the preceding result are not satisfied.
Without loss of generality, assume that in state \(\theta_b\) the inequality \(l^b_b < l^b_g\)
holds.

Under such a condition and depending on the fee schedule designed (namely, on \(t_s\)), the action
chosen by the entrepreneur in certain states may not be the one that maximizes the value of the
project. Therefore, there is room for renegotiation of the contract, in order to guarantee that
there are no ex-post inefficiencies. It is also clear that all the gains from renegotiation revert
to the entrepreneur.

Since the issue of renegotiating or not depends on the design of the fee schedule, one first looks
for sufficient conditions that guarantee the absence of renegotiation and feasibility of the
contract. In this case one would like the contract to induce the entrepreneur to choose action
\(a_b\) in state \(\theta_b\). In other words, the fee schedule and the private benefits to the
entrepreneur should relate as

\[t_s y^b_b + l^b_b \geq t_s y^b_g + l^b_g\]

or \(t_s \geq (l^b_g - l^b_b)/(y^b_b - y^b_g)\). Let this lower bound be denoted by

\[\hat t \equiv \frac{l^b_g - l^b_b}{y^b_b - y^b_g}.\]

For the contract with \(t_s = \hat t\) to be feasible, then it is required that the expected payoff
to the investor be greater than \(K\) or, denoting by \(\pi_1\) the expected payoff,

\[\pi_1 \equiv q(1 - \hat t) y^g_g + (1 - q)(1 - \hat t) y^b_b \geq K.\]

This case is feasible only if \(K\) is low enough as compared to the expected payoff (that is what
the preceding condition means) and/or the fee schedule is high enough (that is the meaning of
\(t_s \geq \hat t\)), transferring to the entrepreneur a relatively high part of the monetary
proceedings of the project.

Consider now the case where \(t_s < \hat t\) for \(s = 0, 1\). In order to maximize the ex-ante
expected payoff, the investor will only accept a contract with \(t_s = 0\) whatever signal is
observed. The schedule \(t(r, s)\) in this contract becomes a simple fixed fee. This fact, together
with the assumption that \(l^b_b < l^b_g\), implies that the entrepreneur will always choose
\(a_g\), independently of the state of nature. Then, the ex-ante expected return to the investor is
given by

\[\pi_2 = q y^g_g + (1 - q) y^b_g.\]

For the contract to be feasible, one requires that \(\pi_2 \geq K\). In this situation, and given
that \(y^b_b > y^b_g\), there is an incentive to renegotiate. The investor is willing to pay up to
\(q y^g_g + (1 - q) y^b_b - \pi_2\) for the entrepreneur to choose action \(a_b\) in state
\(\theta_b\). It turns out that this is an interesting offer to the entrepreneur who in this way
increases its expected payoff. All the bargaining power being on the side of the entrepreneur, the
ex-ante expected payoff to the investor remains \(\pi_2\) after the renegotiation.

\emph{From Amaro de Matos, Notes to Chapter Five, note 1, p.~281 (PDF 299):}

\begin{enumerate}
\def\labelenumi{\arabic{enumi}.}
\tightlist
\item
  Take \(\gamma\) and \(\delta\) as the expected values conditioned on \(\theta\) and \(a\) of two
  random variables. If \(\gamma^g_g > \gamma^g_b\) and \(\gamma^b_b > \gamma^b_g\), then \(\delta\)
  is said to be comonotonic with \(\gamma\) if \(\delta^g_g > \delta^g_b\) and
  \(\delta^b_b > \delta^b_g\).
\end{enumerate}
\stmtend
%% END-STATEMENT

\asked Under entrepreneur control and a contract with $t_s<\hat t$ (so the investor sets $t_s=0$),
show that renegotiating the contract after the state is revealed makes the entrepreneur better off in
expectation than sticking to the original contract.

\approach This is the Aghion--Bolton incomplete-contract model: the contract cannot depend on $\theta$,
and with $t_s=0$ the entrepreneur picks the action with the larger private benefit, which is the wrong
action in $\theta_b$. Ex post both parties observe $\theta$, so there is a surplus from switching to
$a_b$ that can be shared by a side payment (Coasean renegotiation, as in the lecture's discussion of
renegotiation after default). Comparing her payoff with and without the side payment, given that she
has all the bargaining power, gives the result.

\notation
\begin{itemize}
  \item $q\in(0,1)$: probability of the good state $\theta_g$; $a_g$, $a_b$: the actions optimal in $\theta_g$, $\theta_b$.
  \item $y^i_j\in[0,1]$: expected return in state $\theta_i$ under action $a_j$; $l^i_j$: entrepreneur's private benefit.
  \item $t_s$, $t'_s\geq0$: variable and fixed parts of the transfer $t(s,r)=t_sr+t'_s$; $\hat t=(l^b_g-l^b_b)/(y^b_b-y^b_g)$.
  \item $\pi_2=qy^g_g+(1-q)y^b_g$: investor's expected return without renegotiation.
  \item $T$: transfer from the investor to the entrepreneur agreed in renegotiation; $\pi^E_{wt}$, $\pi^E_w$: entrepreneur's expected payoff without and with renegotiation.
\end{itemize}

\assume As in the text, $l^b_b<l^b_g$, $y^b_b>y^b_g$, and $a_b$ is optimal in $\theta_b$:
$y^b_b+l^b_b>y^b_g+l^b_g$; in $\theta_g$, $l^g_g>l^g_b$, so the entrepreneur picks $a_g$ there
(``the entrepreneur will always choose $a_g$''). The fixed fee is zero, as $\pi_2$ implies (a fixed fee
paid in both scenarios would cancel in the comparison). Renegotiation happens after $\theta$ is revealed
and before the action; the entrepreneur has all the bargaining power and offers to take $a_b$ against
a payment $T$ (the text's ``willing to pay up to'').

\step{1}{Payoff without renegotiation} This is the benchmark. With $t_s=0$ (and $t'_s=0$) her reward
$t_sy^i_j+t'_s+l^i_j$ reduces to $l^i_j$, so she takes $a_g$ in both states ($l^g_g>l^g_b$ and
$l^b_g>l^b_b$) and gets
\begin{align}
  \pi^E_{wt} = q\,l^g_g+(1-q)\,l^b_g. \label{t5e4-wt}
\end{align}

\step{2}{Room for renegotiation} We find the largest side payment the investor accepts. In $\theta_g$
the action is already optimal; nothing to renegotiate. In $\theta_b$ switching from $a_g$ to $a_b$
raises the investor's return from $y^b_g$ to $y^b_b$; paying $T$ leaves her $y^b_b-T$, which beats $y^b_g$
iff $T\leq y^b_b-y^b_g$. With all the bargaining power, the entrepreneur asks for the maximum:
\begin{align}
  T = y^b_b-y^b_g. \label{t5e4-T}
\end{align}

\step{3}{The entrepreneur accepts} She switches to $a_b$ if $T+l^b_b\geq l^b_g$. Substituting
\eqref{t5e4-T}:
\begin{align}
  (y^b_b-y^b_g)+l^b_b &\geq l^b_g \nonumber\\
  \Longleftrightarrow\quad y^b_b+l^b_b &\geq y^b_g+l^b_g, \label{t5e4-accept}
\end{align}
which holds strictly by optimality of $a_b$ in $\theta_b$.

\step{4}{Payoff with renegotiation} In $\theta_b$ she now gets $l^b_b+T$. Using \eqref{t5e4-T},
\begin{align}
  \pi^E_w &= q\,l^g_g+(1-q)\bigl[l^b_b+(y^b_b-y^b_g)\bigr]. \label{t5e4-w}
\end{align}
Subtracting \eqref{t5e4-wt} from \eqref{t5e4-w} (the $q\,l^g_g$ terms cancel):
\begin{align}
  \pi^E_w-\pi^E_{wt} &= (1-q)\bigl[l^b_b+(y^b_b-y^b_g)\bigr]-(1-q)\,l^b_g \nonumber\\
  &= (1-q)\bigl[l^b_b+y^b_b-y^b_g-l^b_g\bigr] \nonumber\\
  &= (1-q)\bigl[(y^b_b+l^b_b)-(y^b_g+l^b_g)\bigr]>0, \label{t5e4-gain}
\end{align}
by \eqref{t5e4-accept} and $q<1$. The investor's payoff is unchanged:
\begin{align}
  q\,y^g_g+(1-q)\bigl[y^b_b-(y^b_b-y^b_g)\bigr] = q\,y^g_g+(1-q)\,y^b_g = \pi_2. \label{t5e4-inv}
\end{align}
So renegotiation raises the entrepreneur's expected payoff by the full efficiency gain in the bad
state and leaves the investor at $\pi_2$. $\blacksquare$

\answer{(gain)} $\pi^E_w-\pi^E_{wt}=(1-q)\bigl[(y^b_b+l^b_b)-(y^b_g+l^b_g)\bigr]>0$, with side payment
$T=y^b_b-y^b_g$ in $\theta_b$; the investor's expected payoff stays at $\pi_2=qy^g_g+(1-q)y^b_g$.

\intuition With $t_s<\hat t$ the contract gives the entrepreneur no stake in the return, so in
the bad state she picks the action with the larger private benefit, which destroys value. Ex post
both parties see the loss and can split the gain by a side payment; with all bargaining power on
her side, she captures the whole surplus. Feasibility is unaffected because the investor's payoff
stays at $\pi_2$.
\pitfalls
\begin{itemize}
  \item Giving part of the renegotiation surplus to the investor: the text puts all bargaining power on
        the entrepreneur, so $T$ is the investor's full willingness to pay and her payoff stays $\pi_2$.
  \item Renegotiating in $\theta_g$ too: there $a_g$ is already optimal and there is no surplus.
  \item Skipping the entrepreneur's acceptance check $T+l^b_b\ge l^b_g$: it holds only because $a_b$ is
        optimal in $\theta_b$ ($y^b_b+l^b_b>y^b_g+l^b_g$).
\end{itemize}
\bookans Matches the book's answer (p.~266) in the result. The book's two middle lines carry sign
typos: ``$=ql^g_g+(1-q)(y^b_b-l^b_b)+(1-q)y^b_g$'' should read $ql^g_g+(1-q)(y^b_b+l^b_b)-(1-q)y^b_g$,
and ``$\geq ql^g_g+(1-q)(y^b_g-l^b_g)+(1-q)y^b_g$'' should read $ql^g_g+(1-q)(y^b_g+l^b_g)-(1-q)y^b_g$;
the book also writes $\geq$, while the inequality is strict under the text's strict optimality
assumption and $q<1$.
\refbib{Amaro de Matos Ch.~5, \S 5.1.1 ``The Model'' and \S 5.1.2 ``Entrepreneur Control'', pp.~127--131; Appendix C.5.1, p.~266.}

% =============================================================================
\exsec{5. Diamond (1991): short- vs.\ long-term debt, no liquidation --- Amaro de Matos 5.2.5 \#5, p.~148}
%% BEGIN-STATEMENT 5
\stmtbegin
\textbf{Source:} Amaro de Matos, Ch. 5, Exercise 5.2.5 \#5, p.~148 (PDF 166) · Answer in book: no
(p.~268 only cites example 1 of Diamond 1991, pp.~716--717)

\begin{enumerate}
\def\labelenumi{\arabic{enumi}.}
\setcounter{enumi}{4}
\tightlist
\item
  This exercise illustrates the case of Lemma 5.2.1 in the model of Diamond (1991). Take
  \(L = R = 1\), \(\pi = \frac{1}{2}\), \(X = 1\frac{3}{4}\), \(C = 0\), \(f = \frac{1}{2}\), and
  \(f^d = \frac{1}{3}\).
\end{enumerate}

\begin{enumerate}
\def\labelenumi{(\alph{enumi})}
\item
  What are the face values of a short-term and a long-term debt contract?
\item
  Show that in this example there will be no liquidation.
\item
  What is the payoff of a good project financed with long-term debt?
\item
  What is the probability that a good project receives a downgrade at \(t = 1\)?
\item
  What is the payoff of a good project financed with short-term debt?
\end{enumerate}

\paragraph{Context needed to solve it}

\emph{From Amaro de Matos §5.2.3 ``The Choice of the Duration'', pp.~139--141 (PDF 157--159):}

A model developed by Diamond (1991, 1993) illustrates in a very simple way the differences between
short-term and long-term debt financing. The model is based on information asymmetry and credit
rating.

Regaining the notation of Aghion and Bolton, consider an entrepreneur with a two-period project who
needs funds for an initial investment \(K = 1\). The project may be of one of two types: good or
bad. At time \(t = 2\) a good project has a payoff \(X > 0\), whereas a bad project pays that amount
with a probability \(\pi\), and pays zero with probability \(1 - \pi\). Let \(\theta\) denote the
type of project, \(\theta_g\) referring to a good project and \(\theta_b\) referring to a bad
project. The quality of the project is known to the entrepreneur but not to the investor. Given a
project, as of date \(t = 0\) the investor assigns a probability \(f\) that \(\theta = \theta_g\).
This probability is referred to as a \emph{credit rating} and reflects the lender's information
about the project. It follows that, as of time 0, the probability \(q\) that a borrower pays the
loan at date \(t = 2\) is given by

\[q \equiv f + \pi(1 - f) = \pi + f(1 - \pi). \qquad (5.9)\]

At time \(t = 1\) information about the project becomes available to the investor. This new
information is not verifiable, so it makes no sense to write contracts contingent upon that
information, since they cannot be enforced. However, the investors update their beliefs about the
quality of the project and the credit rating will be either upgraded or downgraded.

Under this revision of the credit rating, let \(f^u\) denote the probability that
\(\theta = \theta_g\) given an upgrade and let \(f^d\) denote the probability that
\(\theta = \theta_g\) given a downgrade. In that case, the probability that an entrepreneur who
receives a downgrade repays at time \(t = 2\) is

\[q^d = \pi + f^d(1 - \pi) > f^d. \qquad (5.10)\]

Moreover, it is here assumed that

\[f^d < f < f^u,\]

meaning that the probability that the project is good is greater given an upgrade than given a
downgrade. Furthermore, the simplifying assumption \(f^u = 1\) is also made. Thus, given an upgrade
at date \(t = 1\), it is certain that the project is good. However, it is not certain that all
projects of good quality are upgraded. There is room for one such project to be downgraded. Let
\(e\) denote the probability of such an event.

All projects are also assumed to produce a nonassignable control rent \(C\) if the management has
control at time \(t = 2\). Also, if at time \(t = 1\) the manager defaults on short-term debt,
lenders then have the right to remove the management from control, selling the assets for a
liquidation value \(L < K = 1\). All projects may be liquidated at that point in time, under those
conditions. If the project is liquidated, the entrepreneur is prevented from consumption of his/her
present or future cash flows and control rents. In this case, an optimal contract is a debt contract
enforced by the right to liquidate the project if the debt is not repaid in full.

Just as in the case of the borrowers, lenders are assumed be risk neutral. Moreover, lenders have
alternative investments at \(t = 0\) that provide a return \(R\) at \(t = 2\). The condition
\(L < X/R\) ensures that a good project yields a higher return when not liquidated.

A project may be financed with either long-term or short-term debt. In the context of this model,
long-term debt refers to a contract issued at \(t = 0\) with no refinancing at date \(t = 1\) and
maturing at \(t = 2\). If one such instrument provides \(K = 1\) at time \(t = 0\) with face value
\(\rho \le X\), the lender expects to receive at maturity \(\rho[\pi + f(1 - \pi)] = R^2\). Notice
that if \(\rho > X\), a borrower with credit rating \(f\) will not be able to provide the lender
with the required rate of return \(R^2\). Now, if the project is upgraded at time \(t = 1\), the
market value of the debt contract at that point in time becomes \(V_l^u = \rho/R\). Similarly, if
the project is downgraded, the market value of debt at time \(t = 1\) becomes \(V_l^d = q^d\rho/R\).

Short-term debt, on the other hand, is debt financed at \(t = 0\) with face value \(r_1\) and
maturity at time \(t = 1\). If short-term borrowers are upgraded at \(t = 1\), their debt contract
at that point in time is worth \(V_s^u = r_1\) and they may refinance their project by issuing a new
short-term debt contract with maturity at \(t = 2\) and face value \(r_1 R\) that will be paid with
certainty and guarantees a return of \(R\) per unit invested by the lender. If short-term borrowers
are downgraded, however, in principle they may refinance their debt with a new short-term contract
only if \(q^d X/R > r_1\), since \(X\) is the maximum that a borrower can pay at \(t = 2\). If that
is the case, at \(t = 1\) the market value of the short-term debt issued at \(t = 0\) would be
\(V_s^d = r_1\) if the borrower pays, possibly through the refinancement mechanism. But what happens
if \(q^d X/R < r_1\)? In that case there is still room for a refinancement provided that the
liquidation value \(L\) is less than the expected value to the lender from a refinancement,
\(q^d X/R\). Under a downgrading, the time \(t = 1\) value of the short-term debt is thus

\[V_s^d = \min\left[r_1, \max(L, q^d X/R)\right].\]

The face value of short-term debt is established such as to provide the lenders with an expected
return of \(R\) per period. It then follows that

\[r_1 = \frac{R - [1 - (1 - e)f]V_s^d}{(1 - e)f}. \qquad (5.11)\]

\emph{From Amaro de Matos §5.2.3, pp.~141--142 (PDF 159--160):}

Given this model, the point is to characterize under what conditions short-term or long-term
maturities will be offered. The maturity preferred by the entrepreneurs of good projects will be the
only one offered to all borrowers. This happens because if any other maturity is offered, it reveals
a project of bad quality and lenders will not be willing to finance it.

The following result implies that, in the absence of liquidation given a downgrade, the
entrepreneurs of good projects prefer short-term debt:

\textbf{Lemma 5.2.1 (Diamond, 1991).} \emph{Comparing two debt contracts issued at \(t = 0\),
neither of which involves liquidation on a downgrade at \(t = 1\), the one with the higher value at
that point in time given a downgrade, \(V^d\), is preferred by the entrepreneurs of good projects.}

\emph{Proof.} The expected payoff to the entrepreneurs of good projects from debt not yielding
liquidation is

\[e\left(X + C - \frac{RV^d}{q^d}\right) + (1 - e)(X + C - r_1 R).\]

Using expression (5.11) this payoff is rewritten as

\[X + C - \frac{R^2}{f} + RV^d e\left(\frac{1}{f^d} - \frac{1}{q^d}\right).\]

Given (5.10) this expression for the payoff is clearly an increasing function of \(V^d\). ■
\stmtend
%% END-STATEMENT

\asked (a) Compute the face value $\rho$ of long-term debt (due at $t=2$) and $r_1$ of short-term debt
(due at $t=1$) for the given numbers. (b) Show that no borrower is ever liquidated at $t=1$.
(c) Compute what a good project's entrepreneur gets with long-term debt. (d) Compute the probability $e$
that a good project is downgraded. (e) Compute what a good project's entrepreneur gets with short-term
debt (and compare with (c)).

\approach Diamond's (1991) maturity model: lenders price debt competitively from the public rating, and
at $t=1$ nonverifiable news (an upgrade or a downgrade) reprices short-term debt. Long-term debt is
priced once, $\rho q=R^2$; short-term debt by (5.11). Liquidation can only happen if a downgraded
borrower cannot roll over, i.e.\ if $q^dX/R<r_1$ and $L$ beats refinancing, so (b) is a check of that
condition. Lemma~5.2.1 then predicts the answer to (c) vs.\ (e): without liquidation, good types prefer the
debt with the higher downgrade value $V^d$, i.e.\ short-term debt, because it shifts repayment onto the
(mostly bad) downgraded pool.

\notation
\begin{itemize}
  \item $L=R=1$: liquidation value and gross per-period safe return; $\pi=\tfrac12$: success probability of a bad project; $X=\tfrac74$: payoff on success; $C=0$: control rent.
  \item $f=\tfrac12$, $f^u=1$, $f^d=\tfrac13$: prior and posterior probabilities of a good project; $e$: probability that a good project is downgraded.
  \item $q$, $q^d$: repayment probabilities at $t=0$ and after a downgrade, (5.9)--(5.10).
  \item $\rho$: face value of long-term debt due at $t=2$; $r_1$: face value of short-term debt due at $t=1$; $V^d$: $t=1$ value of the debt after a downgrade ($V_l^d$ long, $V_s^d$ short).
  \item $U_l$, $U_s$: expected payoff of a good project under long- and short-term debt.
\end{itemize}

\step{1}{Repayment probabilities and $e$} Every later step needs $q$, $q^d$ and $e$, so we compute them
first. From (5.9) and (5.10),
\begin{align}
  q &= \pi+f(1-\pi) = \tfrac12+\tfrac12\left(\tfrac12\right) = \tfrac34, \label{t5e5-q}\\
  q^d &= \pi+f^d(1-\pi) = \tfrac12+\tfrac13\left(\tfrac12\right) = \tfrac23. \label{t5e5-qd}
\end{align}
Bad projects are always downgraded ($f^u=1$) and good ones with probability $e$, so a downgrade has
probability $fe+(1-f)$ and Bayes' rule gives the downgrade posterior:
\begin{align}
  f^d &= \frac{fe}{fe+(1-f)}. \label{t5e5-bayes}
\end{align}
Solving \eqref{t5e5-bayes} for $e$ (multiply out, collect the $e$ terms, divide):
\begin{align}
  f^d\bigl[fe+(1-f)\bigr] &= fe, \nonumber\\
  f^d(1-f) &= fe-f^dfe = fe(1-f^d), \nonumber\\
  e &= \frac{f^d(1-f)}{f(1-f^d)} \label{t5e5-e}\\
    &= \frac{\left(\tfrac13\right)\left(\tfrac12\right)}{\left(\tfrac12\right)\left(\tfrac23\right)}
     = \frac{1/6}{1/3} = \tfrac12. \nonumber
\end{align}
Hence the probability of an upgrade is $(1-e)f=\left(\tfrac12\right)\left(\tfrac12\right)=\tfrac14$ and of a
downgrade $1-\tfrac14=\tfrac34$.

\step{2}{(a) Face values} Each face value comes from the lenders' break-even condition at its own
horizon. \emph{Long-term debt}: the lender lends $K=1$ for two periods and requires $\rho q=R^2$;
substituting \eqref{t5e5-q},
\begin{align}
  \rho\left(\tfrac34\right) &= (1)^2, \nonumber\\
  \rho &= \frac{1}{3/4}=\tfrac43\approx1.3333\leq X=1.75. \label{t5e5-rho}
\end{align}
\emph{Short-term debt}: guess $r_1\leq\max(L,q^dX/R)$, so that $V_s^d=\min[r_1,\max(L,q^dX/R)]=r_1$ (the
debt is repaid in full after a downgrade). Substituting $R=1$, $(1-e)f=\tfrac14$ and $V_s^d=r_1$ into (5.11):
\begin{align}
  r_1 &= \frac{1-\left[1-\tfrac14\right](r_1)}{\tfrac14}, \nonumber\\
  \tfrac14r_1 &= 1-\tfrac34r_1, \nonumber\\
  \tfrac14r_1+\tfrac34r_1 &= 1, \nonumber\\
  r_1 &= 1. \label{t5e5-r1}
\end{align}
Check of the guess: $q^dX/R=\left(\tfrac23\right)\left(\tfrac74\right)/1=\tfrac76\approx1.1667$, so
$\max(L,q^dX/R)=\tfrac76\geq r_1=1$. At $t=1$ the upgraded borrower rolls $r_1$ over into debt of face
value $r_1R=1$, and the downgraded borrower into debt of face value $r_1R/q^d=1/\tfrac23=\tfrac32$.

\answer{(a)} Short-term debt: $r_1=1$ (riskless, $r_1=R$). Long-term debt: $\rho=\tfrac43\approx1.3333$.

\step{3}{(b) No liquidation} Long-term debt has no payment due at $t=1$, so there is no default and
no liquidation. Short-term debt: after an upgrade the borrower refinances at face value $1$; after a
downgrade new lenders lend $r_1=1$ against face value $\tfrac32\leq X=\tfrac74$, and they break even
because $q^d\cdot\tfrac32=\left(\tfrac23\right)\left(\tfrac32\right)=1=Rr_1$. Equivalently the
refinancing condition $q^dX/R=\tfrac76>r_1=1$ holds. Every borrower repays $r_1$ at $t=1$, so the
liquidation right is never triggered. (Even after a default, lenders would refinance rather than
liquidate, since $q^dX/R=\tfrac76>L=1$.) Both contracts are of the kind covered by Lemma~5.2.1.

\answer{(b)} No liquidation: long-term debt has no $t=1$ payment, and with short-term debt every
downgraded borrower rolls over $r_1=1$ into face value $\tfrac32\le X=\tfrac74$ (equivalently
$q^dX/R=\tfrac76>r_1=1$).

\step{4}{(c) Good project, long-term debt} A good project pays $X$ for sure, keeps its control rent $C$
and repays $\rho$:
\begin{align}
  U_l &= X+C-\rho \label{t5e5-Ul}\\
      &= \left(\tfrac74\right)+(0)-\left(\tfrac43\right) = \tfrac{21-16}{12} = \tfrac{5}{12}\approx0.4167. \nonumber
\end{align}
Check with the Lemma's formula, using $V_l^d=q^d\rho/R=\left(\tfrac23\right)\left(\tfrac43\right)=\tfrac89$,
$R^2/f=1/\tfrac12=2$ and $1/f^d-1/q^d=3-\tfrac32=\tfrac32$:
\begin{align}
  X+C-\frac{R^2}{f}+RV^de\left(\frac1{f^d}-\frac1{q^d}\right)
  &= \tfrac74-2+(1)\left(\tfrac89\right)\left(\tfrac12\right)\left(\tfrac32\right) \nonumber\\
  &= -\tfrac14+\tfrac23 = \tfrac{-3+8}{12} = \tfrac{5}{12}. \nonumber
\end{align}

\answer{(c)} $U_l=X+C-\rho=\tfrac5{12}\approx0.4167$.

\step{5}{(d) Probability that a good project is downgraded} From \eqref{t5e5-e}, $e=\tfrac12$: half of
the good projects are downgraded at $t=1$ (and all bad ones are), which is what makes $f^d=\tfrac13$.
Check with \eqref{t5e5-bayes}: $fe/[fe+(1-f)]=\tfrac14/\bigl(\tfrac14+\tfrac12\bigr)=\tfrac13$.

\answer{(d)} $e=\tfrac12$ (the overall downgrade probability is $1-(1-e)f=\tfrac34$).

\step{6}{(e) Good project, short-term debt} With probability $1-e$ it is upgraded and repays $r_1R=1$
at $t=2$; with probability $e$ it is downgraded and repays $RV_s^d/q^d=\tfrac32$ (with $V_s^d=r_1=1$).
From the Lemma's payoff expression,
\begin{align}
  U_s &= e\left(X+C-\frac{RV_s^d}{q^d}\right)+(1-e)(X+C-r_1R) \label{t5e5-Us}\\
      &= \tfrac12\left(\tfrac74+0-\tfrac32\right)+\tfrac12\left(\tfrac74+0-1\right) \nonumber\\
      &= \tfrac12\left(\tfrac14\right)+\tfrac12\left(\tfrac34\right) = \tfrac12. \nonumber
\end{align}
Check: $X+C-R^2/f+RV_s^de(1/f^d-1/q^d)=-\tfrac14+(1)(1)\left(\tfrac12\right)\left(\tfrac32\right)=-\tfrac14+\tfrac34=\tfrac12$.
So $U_s=\tfrac12>U_l=\tfrac{5}{12}$ (gain $\tfrac12-\tfrac5{12}=\tfrac1{12}\approx0.0833$): good projects prefer
short-term debt, as Lemma~5.2.1 predicts since $V_s^d=1>V_l^d=\tfrac89$.

\answer{(e)} $U_s=\tfrac12$, above $U_l=\tfrac5{12}$ by $\tfrac1{12}\approx0.0833$: the good entrepreneur
prefers short-term debt.

\intuition Long-term debt locks in the date-0 pooled price $\rho=\tfrac43$, which subsidises the
bad projects. Short-term debt lets the price be reset at $t=1$: the upgraded good projects (half
of them) refinance at the riskless rate and pay only $1$; the downgraded ones pay more ($\tfrac32$),
but no project is liquidated, so the reset is pure gain for the good type on average. Short maturity
is attractive when refinancing risk is not a liquidation risk.
\pitfalls
\begin{itemize}
  \item Pricing long-term debt with $R$ instead of $R^2$, or with $q^d$ instead of $q$: at $t=0$ the
        relevant repayment probability is the prior one, $q=\tfrac34$, over two periods.
  \item Confusing $e$ (probability that a \emph{good} project is downgraded, $\tfrac12$) with the
        probability of a downgrade for an unknown type ($\tfrac34$).
  \item Charging the downgraded good borrower $r_1R=1$ instead of $RV_s^d/q^d=\tfrac32$: after a downgrade
        new lenders price the rollover at the pooled repayment probability $q^d$.
\end{itemize}
\bookans The book gives no worked answer: Appendix C (p.~268) only cites example~1 of Diamond (1991),
pp.~716--717, which is not in \texttt{Livros/}. The solution above is derived here from \S 5.2.3 and
Lemma~5.2.1. Note that the parameters set $L=1=K$, while the text assumes $L<K=1$; this does not
matter here because liquidation never occurs.
\refbib{Amaro de Matos Ch.~5, \S 5.2.3 ``The Choice of the Duration'', pp.~139--142 (Lemma~5.2.1, p.~142); Appendix C.5.2, p.~268.}

% =============================================================================
\exsec{6. Optimal contracts in the Bolton--Scharfstein model --- Tirole 7.7, p.~325}
%% BEGIN-STATEMENT 6
\stmtbegin
\textbf{Source:} Tirole, Ch. 7, Exercise 7.7, p.~325 (PDF 338) · Answer in book: yes, pp.~603--604

\textbf{Exercise 7.7 (optimal contracts in the Bolton--Scharfstein model).} Redo the Bolton and
Scharfstein analysis of Section 7.1.2, allowing for fully general contracts: the entrepreneur
receives \(r_b^S\) in the case of date-0 success but no refinancing, \(R_b^{SS}\) in the cases of
date-0 and date-1 success, and \(R_b^{FS}\) in the cases of date-1 success and date-0 failure (with
\(R_b^{SS}, R_b^{FS} \geqslant B/\Delta p\)). (Under risk neutrality, there is no point rewarding
failures unless it serves to deter predation. Hence, the exception \(R_b^{FS}\).) Generalize the
conditions (PD) and (IC) and show that \(r_b^S = 0\) and that \(R_b^{SS} \geqslant R_b^{FS}\)
(\(\geqslant B/\Delta p\)).

\paragraph{Context needed to solve it}

\emph{From Tirole §7.1.1.1 ``Basics: Profit Destruction and Benchmarking Effects'', p.~284 (PDF
297):}

The model is the basic, fixed-investment model of Section 3.2 except for the twist that a firm's
profit depends on how successful the other firm is. Namely, while a firm makes no profit if it fails
to develop the new technology, its profit when its succeeds in developing it depends on whether it
faces a competitor, that is, whether the other firm also succeeds. Thus the firm's profit is as
follows:

\[\text{profit} = \begin{cases} M & \text{if it is the only firm to succeed,} \\ D & \text{if both firms succeed,} \\ 0 & \text{if it fails,} \end{cases}\]

where

\[M \geqslant D \geqslant 0 \quad \text{and} \quad p_H M > I.\]

Here, \(M\) stands for ``monopoly profit'' and \(D\) for ``duopoly profit.'' {[}\ldots{]}

The familiar agency cost affects the development process. Each entrepreneur succeeds with
probability \(p_H\) if she behaves (in which case she receives no private benefit) and with
probability \(p_L = p_H - \Delta p\) if she misbehaves (and thereby receives private benefit \(B\)),
where \(\Delta p > 0\).

Each entrepreneur needs to raise \(I - A\) in order to finance her project, where \(I\) is the
investment cost and \(A\) her initial net worth.

\emph{From Tirole §7.1.2.1 ``Financial Muscle and Predation'', p.~290 (PDF 303):}

\emph{Simple-minded long-purse story: short-term financing arrangements.} There are two dates,
\(t = 0, 1\). There is no discounting between the two periods.

Consider a duopoly. Firms \(i = 1, 2\) are identical in all respects except the amount of wealth
they have access to in order to finance investments. Firm 1 (the predator, the financially strong
firm) has a large amount of wealth and never needs to go to the capital market to finance
investment. Firm 2 (the prey, the financially weak firm) has just enough wealth to finance the
date-0 investment.\textsuperscript{16}

{[}\ldots{]} The second period, date 1, is described exactly as in Section 7.1.1: for each firm, the
investment cost is \(I\). Entrepreneurs then engage in moral hazard. The probability of success of
the date-1 project is \(p_H\) if the entrepreneur behaves and \(p_L = p_H - \Delta p\) if she
misbehaves (in which case she obtains private benefit \(B\)). A firm's date-1 profit is \(M\) if it
alone succeeded, \(D\) if both firms succeeded, and 0 otherwise. {[}\ldots{]}

\textsuperscript{16} Alternatively, one could assume that it is able to secure financing for the
date-0 investment, where the loan is to be repaid from date-0 profits: that is, there is no
long-term financing arrangement.

\emph{From Tirole §7.1.2.1, ``Strategic security design (Bolton and Scharfstein)'', pp.~292--293
(PDF 305--306):}

Let us now investigate these points in more detail. To do so, we need to build upon the
simple-minded model by explicitly describing the date-0 actions and by allowing for long-term
financing. To simplify the analysis, let us assume that

\begin{itemize}
\tightlist
\item
  \(p_H = 1\);
\item
  furthermore, date 0 is identical to date 1 except for the private benefit from misbehaving---the
  latter is equal to \(B_0\) at date 0 and \(B\) at date 1;\textsuperscript{19}
\item
  the financially strong firm's act of predation in a given period\textsuperscript{20} results in 0
  profits (failure) for both firms in that period;
\item
  investors in the financially weak firm observe only that firm's profit at date
  0.\textsuperscript{21}
\end{itemize}

As earlier, the financially strong firm is self-financed. The financially weak firm signs a
long-term agreement with its investors that specifies:\textsuperscript{22}

\begin{itemize}
\tightlist
\item
  the probabilities \(z^S\) and \(z^F\) of refinancing (date-1 investment) in the case of date-0
  success (i.e., profit \(D\)) or failure (i.e., profit 0), respectively;
\item
  a reward \(R_b^i\) in the case of date-1 success (which assumes a date-1 reinvestment in the first
  place), where \(i \in \{S, F\}\) indexes the date-0 outcome---for all \(i\), \(R_b^i\) must
  satisfy the date-1 incentive constraint \((\Delta p)R_b^i \geqslant B\).\textsuperscript{23}
\end{itemize}

The timing is summarized in Figure 7.5.

Figure 7.5 (timeline, left to right; p.~293, PDF 306)

{\def\LTcaptype{none} % do not increment counter
\begin{longtable}[]{@{}
  >{\raggedright\arraybackslash}p{(\linewidth - 10\tabcolsep) * \real{0.1667}}
  >{\raggedright\arraybackslash}p{(\linewidth - 10\tabcolsep) * \real{0.1667}}
  >{\raggedright\arraybackslash}p{(\linewidth - 10\tabcolsep) * \real{0.1667}}
  >{\raggedright\arraybackslash}p{(\linewidth - 10\tabcolsep) * \real{0.1667}}
  >{\raggedright\arraybackslash}p{(\linewidth - 10\tabcolsep) * \real{0.1667}}
  >{\raggedright\arraybackslash}p{(\linewidth - 10\tabcolsep) * \real{0.1667}}@{}}
\toprule\noalign{}
\begin{minipage}[b]{\linewidth}\raggedright
Date 0: step 1
\end{minipage} & \begin{minipage}[b]{\linewidth}\raggedright
Date 0: step 2
\end{minipage} & \begin{minipage}[b]{\linewidth}\raggedright
Date 0: step 3
\end{minipage} & \begin{minipage}[b]{\linewidth}\raggedright
Date 1: step 4
\end{minipage} & \begin{minipage}[b]{\linewidth}\raggedright
Date 1: step 5
\end{minipage} & \begin{minipage}[b]{\linewidth}\raggedright
Date 1: step 6
\end{minipage} \\
\midrule\noalign{}
\endhead
\bottomrule\noalign{}
\endlastfoot
Financially weak firm signs initial contract. Financing of investment \(I\). Choice of \(z^S\),
\(z^F\), and \(R_b^i\). & Date-0 efforts. Probability of success: \(p_H = 1\) for high effort or
\(p_L\) for low effort (private benefit \(B_0\)) if no predation; 0 if predation. Strong firm preys?
& Profits: \(D\) if successful and no predation, 0 otherwise. & Reinvestment \(I\) by weak firm? &
Date-1 effort(s). Probability of success: \(p_H = 1\) for high effort, or \(p_L\) for low effort
(private benefit \(B\)). & Date-1 profit(s). \\
\end{longtable}
}

Let us make the following assumptions.

\textbf{Assumption 7.1.} \emph{Positive-NPV investment even under duopoly in the absence of
predation:}

\[D - I > 0.\]

(Recall: the probability of success in the case of good behavior is here equal to 1.)

\textbf{Assumption 7.2.} \emph{Dearth of pledgeable income in the case of
continuation:}\textsuperscript{24}

\[\left(D - \frac{B}{\Delta p}\right) - I < 0.\]

Preventing predation by firm 1 at date 0 benefits firm 2 in two ways. First, its date-0 income is
increased by \(D\), the duopoly profit. Second, it may avoid credit rationing, which from Assumption
7.2 is a threat, and from Assumption 7.1 undesirable.

\emph{Predation deterrence constraint.} To avoid predation firm 2 must choose its financial contract
so that firm 1's date-0 cost of predation, \(D\), exceed its date-1 gain from monopolization. To
compute the latter, note that it is not in firm 1's interest to prey in the last period. Hence,
preventing refinancing by firm 2 raises firm 1's profit from \(D\) to \(M\). And the probability
that firm 2 receives refinancing falls from \(z^S\) in the absence of predation to \(z^F\) under
predation (recall that the probability of refinancing is \(z^S\) if firm 2's date-0 profit is \(D\)
and \(z^F\) if it is equal to 0). The predation deterrence constraint is therefore

\[D \geqslant (z^S - z^F)(M - D). \qquad (\text{PD})\]

To deter predation, the weak firm's contract must make the continuation decision relatively
insensitive to that firm's date-0 profit performance. Note that (PD) can be rewritten as

\[\frac{D}{M - D} \geqslant z^S - z^F. \qquad (\text{PD}')\]

Suppose that competition between the two firms reduces industry profit:

\[M \geqslant 2D.\]

The left-hand side of (PD′) can then take any value between 0 (extreme, Bertrand competition) and 1
(perfect tacit collusion or noncompeting goods). In the latter case, (PD′) really does not constrain
the financial contract and there is little incentive to prey. By contrast, with Bertrand
competition, predation can only be deterred by a performance-insensitive continuation rule (on the
other hand, remaining in the market is also unattractive for firm 2).

\emph{Weak firm's date-0 incentive constraint.} The weak firm's contract must also induce its
entrepreneur to behave. Here, the entrepreneur's compensation is delayed. She receives \(R_b^i\) if
there is reinvestment and firm 2 has profit \(D\) in the last period. Let

\[\mathcal{R}_b^S \equiv z^S R_b^S \quad \text{and} \quad \mathcal{R}_b^F \equiv z^F R_b^F\]

denote the expected continuation payoffs for the entrepreneur in the cases of date-0 success and
failure, respectively. By misbehaving at date 0, the entrepreneur receives private benefit \(B_0\),
but reduces the probability of date-0 success by \(\Delta p\) (provided that the rival does not
prey, i.e., if constraint (PD) is satisfied).

Hence, the incentive constraint is

\[\mathcal{R}_b^S - \mathcal{R}_b^F \geqslant \frac{B_0}{\Delta p}. \qquad (\text{IC})\]

\begin{enumerate}
\def\labelenumi{(\alph{enumi})}
\tightlist
\item
  \emph{The no-predation benchmark.} Suppose, first, that the predator is unable to prey and so
  constraint (PD) is irrelevant. Let \(U_b(z^S)\) denote the NPV:
\end{enumerate}

\[U_b(z^S) \equiv D - I + z^S(D - I).\]

From Assumption 7.1, it increases in \(z^S\).

The investors' breakeven constraint can be written as

\[U_b(z^S) - \mathcal{R}_b^S + A \geqslant 0. \qquad (\text{IR}_l)\]

Finally, the entrepreneur's incentive compatibility constraint is

\[\mathcal{R}_b^S - \mathcal{R}_b^F \geqslant \frac{B_0}{\Delta p}. \qquad (\text{IC})\]

{[}\ldots{]}

\emph{From Tirole §7.1.2.1, p.~294 (PDF 307):}

\begin{enumerate}
\def\labelenumi{(\alph{enumi})}
\setcounter{enumi}{1}
\tightlist
\item
  \emph{Reintroducing the predation-deterrence constraint.} The best, predation-deterring financial
  contract is now obtained by maximizing firm 2's NPV subject to the predation-deterrence constraint
  (PD), the breakeven constraint (IR\(_l\)), and the incentive-compatibility (IC) constraint.
\end{enumerate}

If the solution in the no-predation benchmark case satisfied (PD), then it is also the solution when
predation if feasible. So, we will assume that (PD) is not satisfied by the benchmark solution.
{[}\ldots{]}

\emph{Footnotes 19--24, p.~292 (PDF 305):}

\textsuperscript{19} We will later assume that \(B_0 < B\). This distinction between \(B_0\) and
\(B\) allows us to shorten the analysis by ensuring that date-0 incentives can be provided through
the continuation policy. See below.

\textsuperscript{20} Predation can, of course, occur at date 0 only. The predator can only lose by
preying at date 1 because there is no future at that date.

\textsuperscript{21} The important assumption is that a court cannot ascertain whether the absence
of date-0 profit for firm 2 is due to a predatory act or to moral hazard (lack of luck is not a
possible third explanation because we assumed that \(p_H = 1\)). {[}\ldots{]}

\textsuperscript{22} Exercise 7.7 shows that there is only a slight loss of generality in
considering this class of contracts, and that the results are not qualitatively affected by our
focus on this class.

\textsuperscript{23} Note that if \((\Delta p)R_b^i < B\), then there is always renegotiation to an
\(R_b^i \geqslant B/\Delta p\) as long as \(p_H(R - B/\Delta p) > p_L R\). But even in the absence
of renegotiation, the analysis generalizes.

\textsuperscript{24} Given that rent \(p_H(B/\Delta p) = B/\Delta p\) must be given to the
entrepreneurs, investors would not finance firm 2 at date 1 in the absence of retained earnings and
long-term contracting obligations.
\stmtend
%% END-STATEMENT

\asked Redo the Bolton--Scharfstein contract with three reward levels instead of one per date-0
outcome. \emph{(PD)}: write the predation-deterrence constraint for this contract. \emph{(IC)}: write the
date-0 incentive constraint. \emph{($r_b^S$)}: show that the reward after date-0 success without
refinancing can be set to zero. \emph{(ordering)}: show that the reward after success--success is at least
the reward after failure--success, which is at least $B/\Delta p$.

\approach Bolton--Scharfstein is the lecture's threat-of-termination logic applied to a firm with a
predatory rival: the refinancing probabilities $z^S$, $z^F$ are the termination threat that disciplines
date-0 effort, while predation deterrence requires that threat to be weak. With risk neutrality only
\emph{expected} rewards matter, so the proof is a reshuffling argument: move reward to where it relaxes
the binding constraints and leaves the others unchanged.

\notation
\begin{itemize}
  \item $p_H=1$, $p_L<1$, $\Delta p=1-p_L$: success probabilities; $B_0$, $B$: private benefits from misbehaving at dates 0 and 1 ($B_0<B$).
  \item $D$, $M$: duopoly and monopoly profits ($M\geq D>0$); $I$: investment per date; $A$: firm 2's net worth.
  \item $z^S$, $z^F\in[0,1]$: probabilities of date-1 refinancing after date-0 success (profit $D$) / failure (profit 0).
  \item $r_b^S\geq0$: reward after date-0 success without refinancing; $R_b^{SS}$, $R_b^{FS}\geq B/\Delta p$: rewards after date-1 success following date-0 success / failure.
  \item $W^S\equiv z^SR_b^{SS}+(1-z^S)r_b^S$: entrepreneur's expected reward after date-0 success; $U_b(z^S)=D-I+z^S(D-I)$: NPV.
\end{itemize}

\step{1}{Date-1 behaviour} Whenever firm 2 is refinanced, $R_b^{SS},R_b^{FS}\geq B/\Delta p$ make good
behaviour optimal at date 1, so firm 2 succeeds with probability $p_H=1$ and the entrepreneur
collects $R_b^{SS}$ or $R_b^{FS}$ for sure.

\step{2}{Generalised (PD)} Firm 1's payoffs do not depend on how firm 2 pays its entrepreneur;
they depend only on whether firm 2 is active at date 1, which by Step~1 means profit $D$ for firm 1
instead of $M$. Preying costs firm 1 its date-0 profit $D$ and lowers firm 2's refinancing
probability from $z^S$ to $z^F$, so firm~1's expected date-1 profit rises by
$z^FD+(1-z^F)M-[z^SD+(1-z^S)M]=(z^S-z^F)(M-D)$. Preying is unprofitable iff the cost $D$ is at least this
gain. Hence (PD) is unchanged:
\begin{align}
  D \geq (z^S-z^F)(M-D). \label{t5e6-PD}
\end{align}

\answer{(PD)} $D\ge(z^S-z^F)(M-D)$, the same as with simple contracts: the rewards
$r_b^S,R_b^{SS},R_b^{FS}$ do not enter it.

\step{3}{Generalised (IC)} If the entrepreneur behaves at date 0 (and (PD) holds), date-0 success is
certain and she gets $W^S$. If she misbehaves she gets $B_0$, succeeds with probability $p_L$ (then
$W^S$), and fails with probability $1-p_L$ (then $z^FR_b^{FS}$). Behaving is optimal if
\begin{align}
  W^S &\geq B_0+p_LW^S+(1-p_L)z^FR_b^{FS}, \nonumber\\
  W^S-p_LW^S-(1-p_L)z^FR_b^{FS} &\geq B_0, \nonumber\\
  (1-p_L)\bigl(W^S-z^FR_b^{FS}\bigr) &\geq B_0, \nonumber\\
  (\Delta p)\bigl[z^SR_b^{SS}+(1-z^S)r_b^S-z^FR_b^{FS}\bigr] &\geq B_0, \label{t5e6-IC}
\end{align}
using $\Delta p=p_H-p_L=1-p_L$ and the definition of $W^S$ in the last line.
With $r_b^S=0$, $\mathcal R_b^S=z^SR_b^{SS}$ and $\mathcal R_b^F=z^FR_b^{FS}$ this is the text's (IC).

\answer{(IC)} $(\Delta p)\bigl[z^SR_b^{SS}+(1-z^S)r_b^S-z^FR_b^{FS}\bigr]\ge B_0$.

\step{4}{Generalised (IR$_l$) and the objective} On the equilibrium path date 0 succeeds; investors
receive $D$, pay $r_b^S$ if there is no refinancing, and if there is refinancing invest $I$, receive $D$
and pay $R_b^{SS}$:
\begin{align}
  D-(1-z^S)r_b^S+z^S(D-I-R_b^{SS}) &\geq I-A \nonumber\\
  \Longleftrightarrow\quad (D-I)+z^S(D-I)-\bigl[z^SR_b^{SS}+(1-z^S)r_b^S\bigr]+A &\geq 0 \nonumber\\
  \Longleftrightarrow\quad U_b(z^S)-W^S+A &\geq 0, \label{t5e6-IR}
\end{align}
where the second line moves $I-A$ to the left and groups terms, and the third uses
$U_b(z^S)=D-I+z^S(D-I)$ and the definition of $W^S$.
The entrepreneur's utility is $W^S-A$; with \eqref{t5e6-IR} binding it equals $U_b(z^S)$, which is
increasing in $z^S$ (Assumption~7.1). The contract maximises $U_b(z^S)$ subject to \eqref{t5e6-PD},
\eqref{t5e6-IC}, \eqref{t5e6-IR} and the floors $R_b^{SS},R_b^{FS}\geq B/\Delta p$.

\step{5}{$r_b^S=0$} The pair $(r_b^S,R_b^{SS})$ enters \eqref{t5e6-IC}, \eqref{t5e6-IR} and the
objective only through $W^S$, and does not enter \eqref{t5e6-PD}. Take a contract with
$z^S\in(0,1)$ and $r_b^S>0$, and replace it by
\begin{align}
  \tilde r_b^S = 0,\qquad \tilde R_b^{SS} = R_b^{SS}+\frac{(1-z^S)r_b^S}{z^S}. \label{t5e6-shift}
\end{align}
Then $z^S\tilde R_b^{SS}+(1-z^S)(0)=z^SR_b^{SS}+(1-z^S)r_b^S=W^S$, and $\tilde R_b^{SS}>R_b^{SS}\geq B/\Delta p$,
so every constraint and the objective are unchanged. If $z^S=1$, $r_b^S$ is never paid and can be set to
0. Finally $z^S=0$ is never optimal: from a feasible contract with $z^S=0$, moving to a small
$z^S=\varepsilon>0$ with $r_b^S=0$ and $R_b^{SS}=W^S/\varepsilon$ keeps $W^S$ (hence \eqref{t5e6-IC}), keeps
\eqref{t5e6-PD} for $\varepsilon\leq D/(M-D)$, meets the floor $W^S/\varepsilon\geq B/\Delta p$ for
small $\varepsilon$ (as $W^S\geq B_0/\Delta p>0$), relaxes \eqref{t5e6-IR} because $U_b(\varepsilon)>U_b(0)$,
and raises the objective. Hence, without loss of generality, $r_b^S=0$.

\answer{($r_b^S$)} $r_b^S=0$ without loss of generality: any $r_b^S>0$ can be moved into $R_b^{SS}$,
raising it by $(1-z^S)r_b^S/z^S$, with all constraints and the NPV unchanged.

\step{6}{$R_b^{SS}\geq R_b^{FS}\geq B/\Delta p$} $R_b^{FS}$ is paid only off the equilibrium path:
it appears in \eqref{t5e6-IC} with a negative sign and nowhere else. Lowering it to its floor
$B/\Delta p$ relaxes \eqref{t5e6-IC} and leaves everything else unchanged, so the optimal contract
can set
\begin{align}
  R_b^{FS} = \frac{B}{\Delta p} \leq R_b^{SS}. \label{t5e6-RFS}
\end{align}
The book adds that \eqref{t5e6-IC} binds when \eqref{t5e6-PD} binds (otherwise $z^F$ could be
raised to relax (PD)); then $R_b^{FS}>B/\Delta p$ would leave \eqref{t5e6-IC} slack, so \eqref{t5e6-RFS}
holds at every optimum. With (PD) binding, $z^F=z^S-D/(M-D)$ and $r_b^S=0$, substituting into
\eqref{t5e6-IC}:
\begin{align}
  z^SR_b^{SS}-\left(z^S-\frac{D}{M-D}\right)R_b^{FS} &= \frac{B_0}{\Delta p}, \nonumber\\
  z^S(R_b^{SS}-R_b^{FS})+\frac{D}{M-D}\,R_b^{FS} &= \frac{B_0}{\Delta p}, \label{t5e6-ICbind}
\end{align}
where the first line divides \eqref{t5e6-IC} (with equality) by $\Delta p$ and the second expands
$-(z^S-\frac{D}{M-D})R_b^{FS}$ and groups the $z^S$ terms.

\answer{(ordering)} $R_b^{FS}=B/\Delta p$ (its floor) and $R_b^{SS}\ge B/\Delta p=R_b^{FS}$, so
$R_b^{SS}\ge R_b^{FS}\ge B/\Delta p$.

\intuition Under risk neutrality only expected rewards matter, so the entrepreneur should be paid
where it helps incentives most --- in the continuation branch after success, which also ties her pay
to the date-1 output. Paying her after date-0 failure only weakens date-0 incentives; the contract
does it only because predation deterrence forces refinancing after failure ($z^F>0$), and then the
date-1 incentive floor $B/\Delta p$ is the least it can pay.
\pitfalls
\begin{itemize}
  \item Setting $R_b^{FS}=0$: once the firm is refinanced after failure, the date-1 incentive constraint
        still requires $R_b^{FS}\ge B/\Delta p$; the floor, not zero, is the minimum.
  \item Putting the rewards into (PD): firm~1 cares only about whether firm~2 is active at date~1, so
        (PD) depends on $z^S-z^F$ alone.
  \item Writing the IC with a generic $p_H$: here $p_H=1$, so the date-0 success probability gap is
        $\Delta p=1-p_L$ and success is certain on the equilibrium path.
\end{itemize}
\bookans Matches the book's answer (pp.~603--604): same (PD), (IC$'$) and (IR$'$), $r_b^S=0$ by
shifting reward into $R_b^{SS}$, and $R_b^{FS}=B/\Delta p$. The book argues $R_b^{SS}\geq R_b^{FS}$ via
the binding (IC$'$) and leaves aside $z^S=0$; Step~5 covers that case and Step~6 gives the direct argument.
\refbib{Tirole Ch.~7, \S 7.1.2.1 ``Financial Muscle and Predation'' (Bolton--Scharfstein), pp.~290--294; Answers, pp.~603--604.}

% =============================================================================
\exsec{7. Liquidation incentives with a private signal --- Tirole 3.14, pp.~149--150}
%% BEGIN-STATEMENT 7
\stmtbegin
\textbf{Source:} Tirole, Ch. 3, Exercise 3.14, pp.~149--150 (PDF 162) · Answer in book: yes,
pp.~577--578

\textbf{Exercise 3.14 (liquidation incentives).} This exercise extends the fixed-investment model of
Section 3.2 by adding a signal on the profitability of the project that (a) accrues after effort has
been chosen, and (b) is privately observed. (The following model is used as a building block in a
broader context by Dessi (2005).)

An entrepreneur has cash \(A\) and wants to invest \(I > A\) into a project. The project yields
\(R\) (success) or 0 (failure) at the end. An intermediate signal reveals the probability \(\gamma\)
that the project will succeed, with \(\gamma = \bar\gamma\) or \(\underline{\gamma}\)
(\(\bar\gamma = \underline{\gamma} + \Delta\gamma\) and \(\Delta\gamma > 0\)). The probability,
\(p\), that \(\gamma = \bar\gamma\) depends on the entrepreneur's effort. If the entrepreneur
behaves, then \(p = p_H\) and the entrepreneur receives no private benefit. If the entrepreneur
misbehaves, then \(p = p_L\) and the entrepreneur receives private benefit \(B\). Investors and
entrepreneur are risk neutral and the latter is protected by limited liability. The competitive rate
of return is equal to 0.

Introduce further an option to liquidate after the signal is realized but before the final profit
accrues. Liquidation yields \(L\), and \(L\) is entirely pledgeable to investors.

One will assume that

\[\bar\gamma R > L > \underline{\gamma} R,\]

so that it is efficient to liquidate if and only if the signal is bad; and that

\[p_H \bar\gamma R + (1 - p_H)L > I\]

(which will imply that the NPV is positive).

Figure 3.11 summarizes the timing.

Figure 3.11 (timeline, left to right; p.~150, PDF 163)

{\def\LTcaptype{none} % do not increment counter
\begin{longtable}[]{@{}
  >{\raggedright\arraybackslash}p{(\linewidth - 6\tabcolsep) * \real{0.2500}}
  >{\raggedright\arraybackslash}p{(\linewidth - 6\tabcolsep) * \real{0.2500}}
  >{\raggedright\arraybackslash}p{(\linewidth - 6\tabcolsep) * \real{0.2500}}
  >{\raggedright\arraybackslash}p{(\linewidth - 6\tabcolsep) * \real{0.2500}}@{}}
\toprule\noalign{}
\begin{minipage}[b]{\linewidth}\raggedright
Step 1
\end{minipage} & \begin{minipage}[b]{\linewidth}\raggedright
Step 2
\end{minipage} & \begin{minipage}[b]{\linewidth}\raggedright
Step 3
\end{minipage} & \begin{minipage}[b]{\linewidth}\raggedright
Step 4
\end{minipage} \\
\midrule\noalign{}
\endhead
\bottomrule\noalign{}
\endlastfoot
The entrepreneur borrows \(I - A\) and invests in a project with cost \(I\). & The entrepreneur
behaves (\(\Pr(\gamma = \bar\gamma) = p_H\); no private benefit), or misbehaves
(\(\Pr(\gamma = \bar\gamma) = p_L\); private benefit \(B\)). & State of the world
\(\gamma = \bar\gamma\) or \(\underline{\gamma}\) realized. & Continuation: success (\(R\)) with
probability \(\gamma\), failure (0) with probability \(1 - \gamma\). Or liquidation (yields
\(L\)). \\
\end{longtable}
}

\begin{enumerate}
\def\labelenumi{(\roman{enumi})}
\tightlist
\item
  Suppose first that \(\gamma\) is verifiable. Argue that the entrepreneur should be rewarded solely
  as a function of the realization of \(\gamma\). What is the pledgeable income? Show that the
  project is financed if and only if \(A \geqslant \bar A\), where
\end{enumerate}

\[p_H\left(\bar\gamma R - \frac{B}{\Delta p}\right) + (1 - p_H)L = I - \bar A.\]

\begin{enumerate}
\def\labelenumi{(\roman{enumi})}
\setcounter{enumi}{1}
\tightlist
\item
  Suppose now that \(\gamma\) is observed only by the entrepreneur. This implies that the
  entrepreneur must be induced to tell the truth about \(\gamma\). Without loss of generality,
  consider an incentive scheme in which the entrepreneur receives \(R_b\) in the case she announces
  \(\gamma = \bar\gamma\) (and therefore the project continues) \emph{and} the final profit is
  \(R\), \(L_b\) if she announces \(\gamma = \underline{\gamma}\) (and therefore the project is
  liquidated), and 0 otherwise.
\end{enumerate}

Show that the project is funded if and only if

\[A \geqslant \bar A + \underline{\gamma}\,\frac{B}{(\Delta p)(\Delta\gamma)}.\]

\paragraph{Context needed to solve it}

\emph{From Tirole §3.2.1 ``The Fixed-Investment Model'', pp.~115--116 (PDF 128--129):}

Variants of the following entrepreneurial model will be used in the following: an entrepreneur (also
called the ``insider'' or the ``borrower,'' ``she'') has a project. This project requires a fixed
investment \(I\). The entrepreneur initially has ``assets'' (``cash on hand'' or ``net worth'')
\(A < I\). For the moment we interpret these assets as being cash or liquid securities that can be
used toward covering the cost of investment. {[}\ldots{]} The entrepreneur's cash can either be
invested in the project or used for consumption. To implement the project the entrepreneur must
borrow \(I - A\) from lenders. {[}\ldots{]}

\emph{Project.} If undertaken, the project either succeeds, that is, yields verifiable income
\(R > 0\), or fails, that is, yields no income. The probability of success is denoted by \(p\). The
project is subject to moral hazard. The entrepreneur can ``behave'' (``work,'' ``exert effort,''
``take no private benefit'') or ``misbehave'' (``shirk,'' ``take a private benefit'') {[}\ldots{]}.

Behaving yields probability \(p = p_H\) of success and no private benefit to the entrepreneur, and
misbehaving results in probability \(p = p_L < p_H\) of success and private benefit \(B > 0\)
(measured in units of account) to the entrepreneur. In the ``effort interpretation,'' \(B\) can also
be interpreted as a disutility of effort saved by the entrepreneur when shirking. Let
\(\Delta p \equiv p_H - p_L\).

\emph{Preferences and the loan agreement.} Both the borrower and the potential lenders (or
``investors'') are risk neutral. For notational simplicity, there is no time preference; the rate of
return expected by investors (which is also the riskless rate, due to risk neutrality) is taken to
be 0. The borrower is protected by limited liability, and so her income cannot take negative values.

Lenders behave competitively in the sense that the loan, if any, makes zero profit. {[}\ldots{]}

Let us turn to the loan contract. A contract first stipulates whether the project is financed. If
so, it further specifies how the profit is shared between the lenders and the borrower. The
borrower's limited liability will imply that both sides receive 0 in the case of failure
{[}\ldots{]}. In the case of success, the two parties share the profit \(R\); \(R_b\) goes to the
borrower and \(R_l\) to the lenders. To sum up, we posit an incentive scheme for the entrepreneur of
the following form: \(R_b\) in the case of success, 0 in the case of failure.

{[}\ldots{]} We summarize the timing in Figure 3.1.

Figure 3.1 (timeline, left to right): Loan agreement (sharing rule) → Investment → Moral hazard →
Outcome.

We assume that the project is viable only in the absence of moral hazard. That is, the project has
positive NPV if the entrepreneur behaves,

\[p_H R - I > 0, \qquad (3.1)\]

but negative NPV, even if one includes the borrower's private benefit, if she does not,

\[p_L R - I + B < 0. \qquad (3.2)\]

\emph{From Tirole §3.2.2 ``The Lenders' Credit Analysis'', pp.~116--117 (PDF 129--130):}

{[}\ldots{]} the borrower will therefore behave if the following ``incentive compatibility
constraint'' is satisfied:

\[p_H R_b \geqslant p_L R_b + B \quad \text{or} \quad (\Delta p) R_b \geqslant B. \qquad (\text{IC}_b)\]

From this incentive compatibility constraint we infer that the highest income in the case of success
that can be pledged to the lenders without jeopardizing the borrower's incentives is

\[R - \frac{B}{\Delta p}.\]

The \emph{(expected) pledgeable income} is then

\[\mathcal{P} = p_H\left(R - \frac{B}{\Delta p}\right).\]

Because the lenders must break even in order to be willing to finance the project, a necessary
condition for the borrower to receive a loan is that the expected pledgeable income exceed the
lenders' initial outlay:

\[\mathcal{P} \equiv p_H\left(R - \frac{B}{\Delta p}\right) \geqslant I - A, \qquad (\text{IR}_l)\]

{[}\ldots{]} Thus a necessary condition for financing to be arranged is

\[A \geqslant \bar A = p_H \frac{B}{\Delta p} - (p_H R - I). \qquad (3.3)\]

To make things interesting, we will assume that

\[\bar A > 0 \iff p_H R - I < p_H \frac{B}{\Delta p}, \qquad (3.4)\]

otherwise even a borrower with no wealth of her own would be able to obtain credit. {[}\ldots{]}

Thus, \emph{the borrower must have enough assets in order to be granted a loan.} Note that, if
\(A < \bar A\), the project has positive NPV and yet is not funded. {[}\ldots{]} There is
\emph{credit rationing}. {[}\ldots{]}

Conversely, if \(A \geqslant \bar A\), the entrepreneur is able to secure financing, and so
condition (3.3) is both a necessary and a sufficient condition for financing. The entrepreneur
offers claim \(R_l\) to competitive investors so as not to leave them with a surplus:

\[p_H R_l = I - A.\]

Her stake,

\[R_b = R - R_l = R - \frac{I - A}{p_H} \geqslant R - \frac{I - \bar A}{p_H} = \frac{B}{\Delta p}, \qquad (3.5)\]

then induces her to behave.

{[}\ldots{]} Using the breakeven condition for the lenders (\(p_H R_l = I - A\)), the borrower
obtains net utility or payoff (where ``net'' means that we subtract the consumption utility, \(A\),
that the entrepreneur would get by not undertaking the project):

\[U_b = \begin{cases} 0 & \text{if } A < \bar A, \\ p_H R_b - A = p_H(R - R_l) - A = p_H R - I & \text{if } A \geqslant \bar A. \end{cases} \qquad (3.6)\]

As could have been expected from the zero-profit condition for the lenders, \emph{the borrower
receives the entire social surplus or net present value if the project is funded.}

So, the borrower's utility jumps up at \(A = \bar A\). {[}\ldots{]}
\stmtend
%% END-STATEMENT

\asked (i) With a verifiable interim signal $\gamma$: explain why the entrepreneur's pay should depend
only on $\gamma$, compute the pledgeable income, and show that the project is financed iff $A\ge\bar A$.
(ii) With $\gamma$ seen only by the entrepreneur: design the cheapest scheme $(R_b,L_b)$ that makes her
both behave and report $\gamma$ truthfully, and show that financing now requires
$A\ge\bar A+\underline\gamma B/[(\Delta p)(\Delta\gamma)]$.

\approach Fixed-investment model of Tirole \S3.2 (pledgeable income, $\mathcal P\ge I-A$) with an
interim liquidation option. In (i) the signal is a sufficient statistic for effort, so pay on $\gamma$
and liquidate after bad news. In (ii) the liquidation decision rests on information only the insider
has --- the Topic~5 problem of nonverifiable information --- so the contract must add truth-telling
constraints: the entrepreneur must not gain by hiding bad news to avoid termination. The extra rent
this requires is the increase in the net-worth threshold.

\notation
\begin{itemize}
  \item $A<I$: entrepreneur's cash and the investment; $R$: final profit on success; $L$: liquidation value, fully pledgeable.
  \item $\gamma\in\{\underline\gamma,\bar\gamma\}$: success probability revealed by the signal; $\Delta\gamma=\bar\gamma-\underline\gamma>0$; $\bar\gamma R>L>\underline\gamma R$.
  \item $p_H$, $p_L$, $\Delta p=p_H-p_L>0$: probability of $\gamma=\bar\gamma$ if she behaves / misbehaves; $B$: private benefit.
  \item (i) $w(\gamma)\geq0$: reward after signal $\gamma$; $\mathcal P$: pledgeable income; $\bar A$: financing threshold.
  \item (ii) $R_b$: reward if she announces $\bar\gamma$ and the profit is $R$; $L_b$: reward if she announces $\underline\gamma$ (liquidation); $\mathcal P'$: pledgeable income.
\end{itemize}

\step{1}{(i) Rewards depend on $\gamma$ only} Effort affects only the distribution of $\gamma$; once
$\gamma$ is known, the final outcome is pure noise for the entrepreneur ($\gamma$ is a sufficient
statistic for effort). With risk neutrality, any scheme that pays according to $\gamma$ and the final
profit acts on effort and on the investors' cost only through the expected reward given $\gamma$,
$w(\gamma)$. Paying $w(\gamma)$ directly leaves both unchanged, so rewards can be made contingent on
$\gamma$ alone.

\step{2}{(i) Liquidation policy} With $\gamma$ verifiable, the contract liquidates if and only if
$\gamma=\underline\gamma$: this raises investors' income by $L-\underline\gamma R>0$ after a bad signal,
and continuing after a good one gives $\bar\gamma R>L$. The rewards are transfers that do not depend
on this choice, so it maximises both NPV and pledgeable income.

\step{3}{(i) Incentive constraint and pledgeable income} Pledgeable income is total expected income
minus the smallest expected rent compatible with effort, so we first find that rent. Behaving is optimal if
\begin{align}
  p_Hw(\bar\gamma)+(1-p_H)w(\underline\gamma) &\geq p_Lw(\bar\gamma)+(1-p_L)w(\underline\gamma)+B, \nonumber\\
  (p_H-p_L)w(\bar\gamma)-\bigl[(1-p_L)-(1-p_H)\bigr]w(\underline\gamma) &\geq B, \nonumber\\
  (\Delta p)\bigl[w(\bar\gamma)-w(\underline\gamma)\bigr] &\geq B. \label{t5e7-ICi}
\end{align}
The expected rent $p_Hw(\bar\gamma)+(1-p_H)w(\underline\gamma)$ rises in both rewards, so the cheapest
scheme sets $w(\underline\gamma)=0$ (limited liability) and $w(\bar\gamma)=B/\Delta p$, an expected rent
$p_HB/\Delta p$. With the liquidation policy of Step~2, total expected income is
$p_H\bar\gamma R+(1-p_H)L$. Hence
\begin{align}
  \mathcal P = p_H\bar\gamma R+(1-p_H)L-p_H\frac{B}{\Delta p}. \label{t5e7-Pi}
\end{align}

\step{4}{(i) Financing condition} Investors lend if $\mathcal P\geq I-A$. Using \eqref{t5e7-Pi} and the
definition of $\bar A$, $p_H(\bar\gamma R-B/\Delta p)+(1-p_H)L=I-\bar A$:
\begin{align}
  p_H\left(\bar\gamma R-\frac{B}{\Delta p}\right)+(1-p_H)L &\geq I-A \nonumber\\
  \Longleftrightarrow\quad I-\bar A &\geq I-A \nonumber\\
  \Longleftrightarrow\quad A&\geq\bar A. \label{t5e7-Abar}
\end{align}

\answer{(i)} Pay on $\gamma$ only, $w(\bar\gamma)=B/\Delta p$, $w(\underline\gamma)=0$, liquidate iff
$\gamma=\underline\gamma$; $\mathcal P=p_H(\bar\gamma R-B/\Delta p)+(1-p_H)L$; financed iff $A\ge\bar A$.

\step{5}{(ii) Truth-telling constraints} Now only the entrepreneur sees $\gamma$. Announcing
$\bar\gamma$ gives continuation and an expected reward $\gamma R_b$; announcing $\underline\gamma$ gives
liquidation and $L_b$. Truth-telling in each state requires
\begin{align}
  \bar\gamma R_b &\geq L_b, \label{t5e7-TTg}\\
  L_b &\geq \underline\gamma R_b. \label{t5e7-TTb}
\end{align}

\step{6}{(ii) Effort constraint} Given truthful reports, she gets $\bar\gamma R_b$ after a good signal
and $L_b$ after a bad one, so, as in \eqref{t5e7-ICi},
\begin{align}
  (\Delta p)\bigl(\bar\gamma R_b-L_b\bigr) \geq B. \label{t5e7-ICii}
\end{align}

\step{7}{(ii) Cheapest scheme} Investors receive $p_H\bar\gamma(R-R_b)+(1-p_H)(L-L_b)$, so they
minimise the expected rent $p_H\bar\gamma R_b+(1-p_H)L_b$. Both $R_b$ and $L_b$ should be as low as
\eqref{t5e7-TTb} and \eqref{t5e7-ICii} allow; making both bind,
\begin{align}
  L_b &= \underline\gamma R_b, \label{t5e7-Lb}\\
  (\Delta p)\bigl(\bar\gamma R_b-(\underline\gamma R_b)\bigr) &= B, \nonumber\\
  R_b &= \frac{B}{(\Delta p)(\Delta\gamma)}, \label{t5e7-Rb}\\
  L_b &= \underline\gamma\,\frac{B}{(\Delta p)(\Delta\gamma)}. \nonumber
\end{align}
(Formally: with $\delta\equiv\bar\gamma R_b-L_b\geq B/\Delta p$, \eqref{t5e7-TTb} is
$L_b\geq\underline\gamma\delta/\Delta\gamma$, so the rent $L_b+p_H\delta\geq(\underline\gamma/\Delta\gamma+p_H)\delta$
is minimised at $\delta=B/\Delta p$.) \eqref{t5e7-TTg} holds since $\bar\gamma R_b-L_b=B/\Delta p>0$.

\step{8}{(ii) Pledgeable income and financing condition} Substituting \eqref{t5e7-Lb}--\eqref{t5e7-Rb}
into the rent, and writing $\bar\gamma=\underline\gamma+\Delta\gamma$ in the second line,
\begin{align}
  p_H\bar\gamma R_b+(1-p_H)L_b &= \bigl[p_H\bar\gamma+(1-p_H)\underline\gamma\bigr]\frac{B}{(\Delta p)(\Delta\gamma)} \nonumber\\
  &= \bigl[p_H\underline\gamma+p_H\Delta\gamma+\underline\gamma-p_H\underline\gamma\bigr]\frac{B}{(\Delta p)(\Delta\gamma)} \nonumber\\
  &= \bigl[\underline\gamma+p_H\Delta\gamma\bigr]\frac{B}{(\Delta p)(\Delta\gamma)} \nonumber\\
  &= p_H\frac{B}{\Delta p}+\underline\gamma\,\frac{B}{(\Delta p)(\Delta\gamma)}. \label{t5e7-rent}
\end{align}
So, subtracting this rent from total expected income and using \eqref{t5e7-Pi},
\begin{align}
  \mathcal P' &= p_H\bar\gamma R+(1-p_H)L-p_H\frac{B}{\Delta p}-\underline\gamma\,\frac{B}{(\Delta p)(\Delta\gamma)} \nonumber\\
  &= \mathcal P-\underline\gamma\,\frac{B}{(\Delta p)(\Delta\gamma)}, \label{t5e7-Pii}
\end{align}
and $\mathcal P'\geq I-A$ is, using $\mathcal P=I-\bar A$ from \eqref{t5e7-Abar},
\begin{align}
  I-\bar A-\underline\gamma\,\frac{B}{(\Delta p)(\Delta\gamma)} &\geq I-A \nonumber\\
  \Longleftrightarrow\quad A &\geq \bar A+\underline\gamma\,\frac{B}{(\Delta p)(\Delta\gamma)}. \qquad\blacksquare \label{t5e7-Aii}
\end{align}
(Continuing after both signals would avoid the truth-telling problem but costs the same rent,
since then $(\Delta p)(\Delta\gamma)R_b\geq B$, and loses $(1-p_H)(L-\underline\gamma R)$ of pledgeable income;
so truthful liquidation is better. Implicitly $R_b\leq R$ and $L_b\leq L$.)

\answer{(ii)} $R_b=\dfrac{B}{(\Delta p)(\Delta\gamma)}$, $L_b=\underline\gamma R_b$; pledgeable income falls
by $\underline\gamma B/[(\Delta p)(\Delta\gamma)]$, so the project is funded iff
$A\ge\bar A+\underline\gamma\,\dfrac{B}{(\Delta p)(\Delta\gamma)}$.

\intuition When the signal is private, the entrepreneur must be paid to admit bad news, because
admitting it kills the project in which she has a stake: she gets $L_b=\underline\gamma R_b$, a
``golden parachute'' worth what she would expect by hiding the news. That raises the stake needed
for effort from $B/\Delta p$ to $B/[(\Delta p)(\Delta\gamma)]$ and the expected rent by
$\underline\gamma B/[(\Delta p)(\Delta\gamma)]$; the net-worth threshold rises one for one. It vanishes
when $\underline\gamma=0$ (nothing to hide) and grows as the signal becomes less informative ($\Delta\gamma\to0$).
\pitfalls
\begin{itemize}
  \item Setting $L_b=0$ in (ii): then an entrepreneur with bad news reports $\bar\gamma$ (she gets
        $\underline\gamma R_b>0$ by continuing), the project is never liquidated and the $L$ income is lost.
  \item Writing the effort constraint as $(\Delta p)\bar\gamma R_b\ge B$: after a bad signal she still gets
        $L_b$, so what effort buys is the \emph{difference} $\bar\gamma R_b-L_b$.
  \item Counting $L$ as pledgeable after a good signal or $\bar\gamma R$ after a bad one: the efficient
        policy continues iff $\gamma=\bar\gamma$ because $\bar\gamma R>L>\underline\gamma R$.
\end{itemize}
\bookans Matches the book's answer (pp.~577--578): $L_b=\underline\gamma R_b$, $(\Delta p)(\bar\gamma R_b-L_b)=B$,
and the threshold $\bar A+\underline\gamma B/[(\Delta p)(\Delta\gamma)]$.
\refbib{Tirole Ch.~3, \S 3.2 ``The Fixed-Investment Model'' and \S 3.2.2 ``The Lenders' Credit Analysis'', pp.~115--117; Exercise 3.14, pp.~149--150; Answers, pp.~577--578.}

% =============================================================================
\exsec{8. Decreasing returns to scale and debt maturity --- Tirole Exercise 5.7, p.~232}
%% BEGIN-STATEMENT 8
\stmtbegin
\textbf{Source:} Tirole, Ch. 5, Exercise 5.7, p.~232 (PDF 245) · Answer in book: yes, p.~593

Exercise 5.7 (decreasing returns to scale). Extend the treatment of Section 5.6.1 to the case of
decreasing returns to scale: the payoff in the case of continuation and success is \(R(I)\), with
\(R(0) = 0\), \(R' > 0\), \(R'' < 0\), \(R'(0) = \infty\), and \(R'(\infty) = 0\). The rest is
unchanged (the short-term income is \(rI\), the reinvestment need is \(\rho I\), and the private
benefit is \(BI\)).

\begin{enumerate}
\def\labelenumi{(\roman{enumi})}
\item
  What are the first-order conditions yielding the optimal investment level \(I\) and cutoff
  \(\rho^*\)?
\item
  Assuming that \(r > \rho^*\), and that \((R(I)/I - R'(I))\) is increasing in \(I\) (a condition
  satisfied, for example, by \(R(I)\) quadratic), derive the impact of the strength of the balance
  sheet as measured, say, by \(A\) on debt maturity.
\end{enumerate}

\paragraph{Context needed to solve it}

\emph{From Tirole §3.4.1 ``The Continuous-Investment Model'', p.~127 (PDF 140):}

The continuous-investment model of this section is the polar opposite of the fixed-investment one.
The fixed-investment model depicts a situation in which returns are sharply decreasing beyond some
investment level. In contrast, we now assume that there are constant returns to scale in the
investment technology. An investment \(I \in [0, \infty)\) yields income \(RI\), proportional to
\(I\), in the case of success, and 0 in the case of failure. The borrower's private benefit from
misbehaving is also taken to be proportional to investment. As before, the borrower has a choice
between behaving, in which case she derives no private benefit and the probability of success is
\(p_H\), and misbehaving, that is, enjoying private benefit \(BI\), and reducing the probability of
success to \(p_L = p_H - \Delta p < p_H\). {[}\ldots{]}

The borrower initially has cash \(A\), and must therefore borrow \(I - A\) to finance a project of
size \(I\). A loan agreement specifies that the lenders (who as before are assumed to make no
profit) and the borrower receive 0 each in the case of failure, and \(R_l\) and \(R_b\),
respectively, in the case of success, where \(R_l + R_b = RI\).

\emph{From Tirole §5.3.2 ``Continuum of Liquidity Shocks'', pp.~209--210 (PDF 222--223):}

After the (endogenous size) investment \(I\) is sunk at date 0 and before the borrower works on the
project, some exogenous shock occurs at date 1 that determines a per-unit-of-investment level
\(\rho \in [0, \infty)\) of ``cost overruns.'' That is, a cash infusion equal to \(\rho I\) is
needed in order for the project to continue. If \(\rho I\) is not invested, the project is abandoned
altogether and thus yields no income. As in Section 5.2, the fraction \(\rho\) is \emph{a priori}
distributed according to the continuous distribution \(F(\rho)\) on \([0, \infty)\), with density
\(f(\rho)\). (As we already observed, the model of Section 3.4 is therefore a special case, with
\(F\) being a spike at \(\rho = 0\).)

Regardless of the required amount of the cash infusion, the project, if pursued, is still a project
of size \(I\), in that the income in the case of success is \(RI\) and the borrower's private
benefit from misbehaving is \(BI\). One cannot increase the size of the project after the initial
stage.

The timing is summarized in Figure 5.4.

Figure 5.4 (timing, transcribed): \textbf{Date 0} --- Loan agreement; Investment \(I\). \textbf{Date
1} --- Need for cash infusion \(\rho I\) realized; then either Disbursement → Moral hazard, or No
disbursement → Project is abandoned. \textbf{Date 2} --- Outcome.

{[}\ldots{]} It is easy to show that it is optimal to have a ``cutoff rule'' for infusing cash.
There exists an optimal threshold \(\rho^*\) such that one should continue if and only if

\[\rho \leqslant \rho^*. \qquad (5.2)\]

The incentive constraint in the case of continuation is the same as in the absence of a liquidity
shock (see Section 3.4):

\[(\Delta p)R_b \geqslant BI. \qquad (\text{IC}_b)\]

The breakeven condition is slightly altered by the presence of liquidity shocks:

\[F(\rho^*)[p_H(RI - R_b)] \geqslant I - A + \int_0^{\rho^*}\rho If(\rho)\,\mathrm{d}\rho. \qquad (\text{IR}_l)\]

That is, the lenders receive a return only if the project is pursued, which has probability
\(F(\rho^*)\). The left-hand side of (IR\(_l\)) is the expected pledgeable income. Furthermore,
there is a new term, representing the expected outlay on overruns, on the right-hand side. From
these two constraints, we deduce the borrowing capacity (or, more precisely, the maximum investment
that allows the lenders to break even):

\[I = k(\rho^*)A,\]

where

\[k(\rho^*) = \frac{1}{1 + \int_0^{\rho^*}\rho f(\rho)\,\mathrm{d}\rho - F(\rho^*)[p_HR - p_HB/\Delta p]} = \frac{1}{1 + \int_0^{\rho^*}\rho f(\rho)\,\mathrm{d}\rho - F(\rho^*)\rho_0} \qquad (5.3)\]

involves a straightforward modification relative to the no-liquidity-shock multiplier \(k\).
{[}\ldots{]} Note that the borrower's borrowing capacity is maximal when the threshold \(\rho^*\) is
equal to the expected per-unit pledgeable income \(\rho_0 \equiv p_H(R - B/\Delta p)\).

Given that the competitive lenders make no profits, the borrower's net utility is as usual the
social surplus brought about by the project, namely,

\[U_b = m(\rho^*)I = m(\rho^*)k(\rho^*)A, \qquad (5.4)\]

where

\[m(\rho^*) \equiv F(\rho^*)p_HR - 1 - \int_0^{\rho^*}\rho f(\rho)\,\mathrm{d}\rho\]

is the margin per unit of investment.

{[}\ldots{]} Next, we observe that \emph{this optimal threshold lies between the expected
per-unit-of-investment pledgeable income and income}:

\[\rho_0 = p_H\left(R - \frac{B}{\Delta p}\right) < \rho^* < \rho_1 = p_HR. \qquad (5.8)\]

\emph{From Tirole §5.6.1 ``Optimal Claim Maturity'', p.~226 (PDF 239):}

Let us return to the continuous-investment, continuous-shock version of Section 5.3.2, but with a
short-term income (the analysis is not really new and is therefore only sketched): see Figure 5.10.
Because the short-term income \(rI\) is fully pledgeable to the lenders, everything is \emph{as if}
the unit investment cost were equal to \(1 - r\) instead of 1. The lenders' breakeven condition,
that is, the equality between expected revenue and expected investment cost, becomes

\[rI + F(\rho^*)\rho_0I = I - A + \left[\int_0^{\rho^*}\rho f(\rho)\,\mathrm{d}\rho\right]I\]

and so

\[k(\rho^*) = \frac{1}{1 + \int_0^{\rho^*}\rho f(\rho)\,\mathrm{d}\rho - [r + F(\rho^*)\rho_0]}. \qquad (5.3'')\]

The margin (expected profit of the firm per unit of investment) becomes

\[m(\rho^*) = [r + F(\rho^*)\rho_1] - \left[1 + \int_0^{\rho^*}\rho f(\rho)\,\mathrm{d}\rho\right]. \qquad (5.4'')\]

And thus the borrower's (gross) utility becomes

\[U_b = m(\rho^*)k(\rho^*)A = \frac{\rho_1 - c(\rho^*)}{c(\rho^*) - \rho_0}A,\]

where the expected unit cost of effective investment, \(c(\rho^*)\), is given by

\[c(\rho^*) = \frac{1 - r + \int_0^{\rho^*}\rho f(\rho)\,\mathrm{d}\rho}{F(\rho^*)}. \qquad (5.5'')\]

So the optimal threshold is given by

\[\int_0^{\rho^*}F(\rho)\,\mathrm{d}\rho = 1 - r \qquad (5.6'')\]

and the borrower's utility by

\[U_b = \frac{\rho_1 - \rho^*}{\rho^* - \rho_0}A. \qquad (5.7'')\]

Figure 5.10 (timing, transcribed): \textbf{Date 0} --- Loan agreement. Investment \(I\).
\textbf{Date 1} --- Accrual of short-term income \(rI\). Realization \(\rho I\) of investment need.
Then either Disbursement → Moral hazard, or No disbursement → Project is abandoned. \textbf{Date 2}
--- Date-2 income (\(RI\) or 0).

It is important to note that \emph{the short-term income, even though it is deterministic and fully
pledgeable, is not equivalent to an increase in the borrower's cash on hand \(A\).} Such an increase
in equity would result in a larger investment (as is the case here), but not in a modification of
the continuation rule. By contrast, condition (5.6\(''\)) shows that the larger the short-term
profit, the lower the optimal threshold \(\rho^*\). {[}\ldots{]}

\begin{enumerate}
\def\labelenumi{(\alph{enumi})}
\tightlist
\item
  \emph{Liquidity management.} Let us now turn to the implementation of the optimum and thus to the
  claim maturity. To this purpose we make the following assumption. (Footnote 57: Of course, it must
  also be the case that \(\rho^* > \rho_0\) (otherwise, the borrower's borrowing capacity and
  utility would be infinite in this constant-returns-to-scale model). Because
  \(\mathrm{d}\rho^*/\mathrm{d}r < -1\), we must thus also assume that \(r\) is not ``too large.'')
\end{enumerate}

\emph{Free-cash-flow assumption:} \(r > \rho^*\). Under the free-cash-flow assumption, and given
that the entrepreneur cannot steal the intermediate income, the entrepreneur would reinvest
excessively if she were not asked to pay out money to investors at date 1. Namely, she would
reinvest as long as \(\rho \leqslant r\).

To obtain the optimal amount of reinvestment, an amount \(P_1 \equiv (r - \rho^*)I\) must be pumped
out of the firm, and the entrepreneur must be denied the right to dilute initial investors.
\stmtend
%% END-STATEMENT

\asked (i) In the liquidity model with short-term income $rI$ and a continuum of reinvestment shocks
$\rho I$, but with a concave success payoff $R(I)$, derive the first-order conditions that pin down the
scale $I$ and the continuation cutoff $\rho^*$. (ii) Under free cash flow ($r>\rho^*$), find how a
stronger balance sheet (higher $A$) changes debt maturity, i.e.\ the short-term debt the firm must pay at
date~1 per unit of investment.

\approach Holmstr\"om--Tirole liquidity model (Tirole \S5.3.2 and \S5.6.1): the borrower maximises the
NPV subject to the lenders' break-even, a Lagrangian problem in $(I,\rho^*)$. Continuing after a shock
$\rho$ costs $\rho I$ and is worth $\rho_1$ in value but only $\rho_0$ in pledgeable income, so the
optimal cutoff lies between them. Under free cash flow the cutoff is implemented by short-term debt that
pumps $(r-\rho^*)I$ out of the firm --- the termination threat of the lecture, here a liquidity sweep ---
so maturity moves with $\rho^*$, and comparative statics in $A$ answer (ii).

\notation
\begin{itemize}
  \item $I>0$: investment (endogenous); $A\ge 0$: the entrepreneur's cash; lenders are competitive, rate of return 0.
  \item $R(I)$: date-2 income in case of continuation and success, $R(0)=0$, $R'>0$, $R''<0$.
  \item $rI$: short-term (date-1) income, verifiable and fully pledgeable, $r>0$.
  \item $\rho\ge 0$: per-unit liquidity shock (reinvestment need $\rho I$), c.d.f.\ $F$, density $f$ on $[0,\infty)$; $\rho^*$: continuation cutoff (continue iff $\rho\le\rho^*$).
  \item $p_H>p_L$: success probabilities when behaving / misbehaving, $\Delta p=p_H-p_L$; $BI$: private benefit from misbehaving; $R_b$: entrepreneur's stake in success.
  \item $\rho_1(I)\equiv p_H R(I)/I$ and $\rho_0(I)\equiv p_H[R(I)/I-B/\Delta p]$: per-unit expected income and per-unit pledgeable income at scale $I$; $\rho_1-\rho_0=p_HB/\Delta p$.
  \item $\mu\ge 0$: Lagrange multiplier of the lenders' breakeven constraint.
\end{itemize}

\step{1}{The program} We set up the borrower's problem. As in Section 5.6.1, the incentive constraint
$(\Delta p)R_b\ge BI$ binds (any larger stake only lowers pledgeable income), so the income pledgeable in
case of success is $R(I)-BI/\Delta p$. With competitive lenders the entrepreneur gets the NPV: short-term
income, plus expected date-2 income when continuing (probability $F(\rho^*)$), minus the investment and the
expected reinvestment,
\begin{align}
  U_b &= rI + F(\rho^*)\,p_H R(I) - \Big[1+\int_0^{\rho^*}\rho f(\rho)\,d\rho\Big] I, \label{t5e8-Ub}
\end{align}
subject to the lenders' breakeven constraint
\begin{align}
  rI + F(\rho^*)\,p_H\Big[R(I)-\frac{BI}{\Delta p}\Big] &\ge I - A + \Big[\int_0^{\rho^*}\rho f(\rho)\,d\rho\Big] I. \label{t5e8-IR}
\end{align}
The Lagrangian is $\mathcal L = U_b + \mu\,[\text{LHS}-\text{RHS of }\eqref{t5e8-IR}]$.

\step{2}{(i) First-order conditions} We differentiate $\mathcal L$ with respect to each choice variable.
By the fundamental theorem of calculus, $\frac{d}{d\rho^*}\int_0^{\rho^*}\rho f\,d\rho=\rho^* f(\rho^*)$, and
$\frac{d}{d\rho^*}F(\rho^*)=f(\rho^*)$. Term by term,
\begin{align}
  \frac{\partial U_b}{\partial\rho^*} &= f(\rho^*)\,p_HR(I)-\rho^*f(\rho^*)\,I, \nonumber\\
  \frac{\partial[\text{LHS}-\text{RHS of }\eqref{t5e8-IR}]}{\partial\rho^*} &= f(\rho^*)\,p_H\Big[R(I)-\frac{BI}{\Delta p}\Big]-\rho^*f(\rho^*)\,I, \nonumber
\end{align}
so $\partial\mathcal L/\partial\rho^*=0$ reads
\begin{align}
  0 &= f(\rho^*)\big[p_HR(I)-\rho^*I\big] + \mu f(\rho^*)\Big[p_HR(I)-\frac{p_HBI}{\Delta p}-\rho^*I\Big]. \nonumber
\end{align}
Dividing by $f(\rho^*)I>0$ and using $\rho_1(I)=p_HR(I)/I$ and $\rho_0(I)=p_H[R(I)/I-B/\Delta p]$,
\begin{align}
  0 &= \big[\rho_1(I)-\rho^*\big] + \mu\big[\rho_0(I)-\rho^*\big]. \label{t5e8-focrho}
\end{align}
Collecting the $\rho^*$ terms, $(1+\mu)\rho^*=\rho_1(I)+\mu\rho_0(I)$, so
$\rho^*=\frac{\rho_1(I)+\mu\rho_0(I)}{1+\mu}$, a weighted average that lies between $\rho_0(I)$ and $\rho_1(I)$,
the analogue of (5.8). With respect to $I$, the derivatives are
\begin{align}
  \frac{\partial U_b}{\partial I} &= r+F(\rho^*)p_HR'(I)-1-\int_0^{\rho^*}\!\rho f\,d\rho, \nonumber\\
  \frac{\partial[\text{LHS}-\text{RHS of }\eqref{t5e8-IR}]}{\partial I} &= r+F(\rho^*)p_H\Big(R'(I)-\frac{B}{\Delta p}\Big)-1-\int_0^{\rho^*}\!\rho f\,d\rho, \nonumber
\end{align}
so $\partial\mathcal L/\partial I=0$ is
\begin{align}
  0 &= \Big[r+F(\rho^*)p_HR'(I)-1-\int_0^{\rho^*}\!\rho f\,d\rho\Big]
       + \mu\Big[r+F(\rho^*)p_H\Big(R'(I)-\frac{B}{\Delta p}\Big)-1-\int_0^{\rho^*}\!\rho f\,d\rho\Big]. \label{t5e8-focI}
\end{align}
These two conditions, together with \eqref{t5e8-IR} holding with equality, determine $(I,\rho^*,\mu)$.

\step{3}{(i) Eliminating $\mu$} The multiplier is not observable, so we combine the two conditions into
one that involves only $(I,\rho^*)$. Write the second bracket of \eqref{t5e8-focI} as the first minus
$F(\rho^*)p_HB/\Delta p$ and collect:
\begin{align}
  (1+\mu)\Big[r+F(\rho^*)p_HR'(I)-1-\int_0^{\rho^*}\!\rho f\,d\rho\Big] &= \mu F(\rho^*)\frac{p_HB}{\Delta p}. \label{t5e8-focI2}
\end{align}
From \eqref{t5e8-focrho}, $\mu(\rho^*-\rho_0)=\rho_1-\rho^*$, so
$\mu=\frac{\rho_1-\rho^*}{\rho^*-\rho_0}$ and $1+\mu=\frac{\rho^*-\rho_0+\rho_1-\rho^*}{\rho^*-\rho_0}=\frac{\rho_1-\rho_0}{\rho^*-\rho_0}$,
hence $\frac{\mu}{1+\mu}=\frac{\rho_1-\rho^*}{\rho_1-\rho_0}$. Since
$\rho_1-\rho_0=p_H\frac{R(I)}{I}-p_H\Big(\frac{R(I)}{I}-\frac{B}{\Delta p}\Big)=\frac{p_HB}{\Delta p}$,
\begin{align}
  \frac{\mu}{1+\mu}\,\frac{p_HB}{\Delta p} &= \frac{\rho_1(I)-\rho^*}{\rho_1(I)-\rho_0(I)}\,\bigl(\rho_1(I)-\rho_0(I)\bigr) = \rho_1(I)-\rho^*. \label{t5e8-mu}
\end{align}
Dividing \eqref{t5e8-focI2} by $1+\mu$ and substituting \eqref{t5e8-mu}:
\begin{align}
  r+F(\rho^*)p_HR'(I)-1-\int_0^{\rho^*}\!\rho f\,d\rho &= F(\rho^*)\Big(\frac{p_HR(I)}{I}-\rho^*\Big). \nonumber
\end{align}
Move the $R$ terms to the left and the rest to the right:
\begin{align}
  F(\rho^*)\,p_H\Big[R'(I)-\frac{R(I)}{I}\Big] &= 1-r+\int_0^{\rho^*}\!\rho f\,d\rho-\rho^*F(\rho^*) \nonumber\\
  &= 1-r-\Big[\rho^*F(\rho^*)-\int_0^{\rho^*}\!\rho f\,d\rho\Big]. \nonumber
\end{align}
The bracket is $\int_0^{\rho^*}(\rho^*-\rho)f(\rho)\,d\rho$; integrating by parts with $F(0)=0$,
$\int_0^{\rho^*}(\rho^*-\rho)f\,d\rho=\big[(\rho^*-\rho)F(\rho)\big]_0^{\rho^*}+\int_0^{\rho^*}F(\rho)\,d\rho=\int_0^{\rho^*}F(\rho)\,d\rho$.
Dividing by $F(\rho^*)>0$,
\begin{align}
  p_H\Big[R'(I)-\frac{R(I)}{I}\Big] &= \frac{1-r-\int_0^{\rho^*}F(\rho)\,d\rho}{F(\rho^*)}. \label{t5e8-key}
\end{align}
The binding \eqref{t5e8-IR}, solved for $A$ and with $p_H[R(I)-BI/\Delta p]=I\rho_0(I)$, reads
\begin{align}
  A &= I-rI+\Big[\int_0^{\rho^*}\rho f(\rho)\,d\rho\Big]I-F(\rho^*)\,I\rho_0(I) \nonumber\\
    &= I\Big[1-r+\int_0^{\rho^*}\rho f(\rho)\,d\rho - F(\rho^*)\rho_0(I)\Big]. \label{t5e8-IRb}
\end{align}
Check: with constant returns, $R(I)=RI$, the left side of \eqref{t5e8-key} is 0 and \eqref{t5e8-key} becomes
$\int_0^{\rho^*}F=1-r$, which is (5.6$''$). With strict concavity and $R(0)=0$, $R(I)/I>R'(I)$, so the left
side is negative and $\int_0^{\rho^*}F>1-r$: the marginal unit of scale is worth less than the average one,
so at the margin the firm spends more on continuation and less on scale than (5.6$''$) would say.

\answer{(i)} The FOCs are \eqref{t5e8-focrho}, $\rho_1(I)-\rho^*+\mu[\rho_0(I)-\rho^*]=0$, and
\eqref{t5e8-focI}; eliminating $\mu$, the optimum $(I,\rho^*)$ solves
$p_H[R'(I)-R(I)/I]=\bigl[1-r-\int_0^{\rho^*}F(\rho)\,d\rho\bigr]/F(\rho^*)$ together with the binding
break-even \eqref{t5e8-IRb}, with $\rho_0(I)<\rho^*<\rho_1(I)$.

\step{4}{(ii) $I$ and $\rho^*$ move together} Condition \eqref{t5e8-key} links the two choices; we find
the sign of that link. Write $\varphi(I)\equiv p_H[R(I)/I-R'(I)]>0$, increasing by assumption, and
$h(\rho^*)$ for the right side of \eqref{t5e8-key}, so \eqref{t5e8-key} is $-\varphi(I)=h(\rho^*)$. The
numerator of $h$ is $N(\rho^*)=1-r-\int_0^{\rho^*}F$, with $N'(\rho^*)=-F(\rho^*)$; by the quotient rule,
\begin{align}
  h'(\rho^*) &= \frac{N'(\rho^*)F(\rho^*)-N(\rho^*)f(\rho^*)}{F(\rho^*)^2} \nonumber\\
             &= \frac{-F(\rho^*)^2 - f(\rho^*)\big[1-r-\int_0^{\rho^*}F\big]}{F(\rho^*)^2} \nonumber\\
             &= -1 - \frac{f(\rho^*)}{F(\rho^*)}\,h(\rho^*). \label{t5e8-hprime}
\end{align}
\assume Tirole's answer takes $h$ decreasing. Since $h<0$ at the optimum, \eqref{t5e8-hprime} is negative iff
$f(\rho^*)\big[\int_0^{\rho^*}F-(1-r)\big]<F(\rho^*)^2$. This holds, e.g., for $F$ uniform on $[0,\bar\rho]$
($f=1/\bar\rho$, $F(\rho^*)=\rho^*/\bar\rho$, $\int_0^{\rho^*}F=\rho^{*2}/(2\bar\rho)$): multiplying by $\bar\rho^2$, the
condition reads $\rho^{*2}/2-\bar\rho(1-r)<\rho^{*2}$, always true for $r\le1$. We assume it.

Then along \eqref{t5e8-key} an increase in $I$ lowers $-\varphi(I)$, which requires a lower $h$, i.e.\ a higher
$\rho^*$. Totally differentiating $-\varphi(I)=h(\rho^*)$ gives $-\varphi'(I)\,dI=h'(\rho^*)\,d\rho^*$, so
\eqref{t5e8-key} defines $\rho^*=g(I)$ with
\begin{align}
  g'(I) &= \frac{-\varphi'(I)}{h'(\rho^*)} > 0. \label{t5e8-g}
\end{align}

\step{5}{(ii) Effect of $A$} The binding break-even \eqref{t5e8-IRb} tells how much net worth a given
$(I,\rho^*)$ requires; we differentiate its right side. With respect to $\rho^*$,
\begin{align}
  \frac{\partial A}{\partial\rho^*} &= I\big[\rho^*f(\rho^*)-f(\rho^*)\rho_0(I)\big] \nonumber\\
  &= I f(\rho^*)\big[\rho^*-\rho_0(I)\big] > 0, \label{t5e8-dArho}
\end{align}
positive because $\rho^*>\rho_0(I)$ by \eqref{t5e8-focrho}. With respect to $I$ (using $\frac{d}{dI}[I\rho_0(I)]=p_H[R'(I)-B/\Delta p]$),
\begin{align}
  \frac{\partial A}{\partial I} &= 1-r+\int_0^{\rho^*}\!\rho f\,d\rho - F(\rho^*)p_H\Big[R'(I)-\frac{B}{\Delta p}\Big]. \nonumber
\end{align}
The second bracket of \eqref{t5e8-focI} is $-\partial A/\partial I$ and the first is $-\partial A/\partial I+F(\rho^*)p_HB/\Delta p$,
so \eqref{t5e8-focI} reads $-(1+\mu)\,\partial A/\partial I + F(\rho^*)p_HB/\Delta p=0$, i.e.
\begin{align}
  \frac{\partial A}{\partial I} &= \frac{F(\rho^*)\,p_HB}{(1+\mu)\Delta p} > 0. \label{t5e8-dAI}
\end{align}
Along the optimum, $\rho^*=g(I)$, so by \eqref{t5e8-g}, \eqref{t5e8-dArho} and \eqref{t5e8-dAI}
\begin{align}
  \frac{dA}{dI} &= \frac{\partial A}{\partial I}+\frac{\partial A}{\partial\rho^*}\,g'(I) > 0, \label{t5e8-dAdI}
\end{align}
hence $dI/dA=1/(dA/dI)>0$ and $d\rho^*/dA=g'(I)\,dI/dA>0$: a weaker balance sheet lowers both the scale
and the cutoff.

\step{6}{(ii) Debt maturity} Under the free-cash-flow assumption $r>\rho^*$ the optimum is implemented as in
Section 5.6.1(a): the firm keeps $\rho^*I$ to meet shocks $\rho\le\rho^*$, must pay out
$P_1=(r-\rho^*)I$ at date 1 (short-term debt), is barred from diluting the initial investors, and the investors' long-term
claim is $R(I)-BI/\Delta p$ in case of success. Per unit of investment the short-term debt is $r-\rho^*$, and by Step 5
\begin{align}
  \frac{d(r-\rho^*)}{dA} &= -\frac{d\rho^*}{dA} < 0. \label{t5e8-mat}
\end{align}
A firm with a weaker balance sheet (lower $A$) issues more short-term debt per unit of investment: its debt
maturity is shorter. (The level $(r-\rho^*)I$ can move either way, because $I$ falls too.)

\answer{(ii)} $dI/dA>0$ and $d\rho^*/dA>0$, so the short-term debt per unit of investment, $r-\rho^*$,
falls with $A$: weaker balance sheets carry shorter maturity (more of the date-1 cash flow is paid out),
provided $h'(\rho^*)<0$ (true, e.g., for uniform shocks).

\intuition A weak balance sheet makes pledgeable income scarce. The firm raises it by continuing less often,
i.e.\ by keeping less liquidity in the firm, which means sweeping more of the date-1 cash flow out to creditors
through short-term debt. Decreasing returns add one more channel: a lower scale raises $R(I)/I$ and hence per-unit
pledgeable income, and the two margins (scale and continuation) are cut together.

\pitfalls
\begin{itemize}
  \item Carrying over the constant-returns result $\int_0^{\rho^*}F=1-r$, under which $\rho^*$ does not
        depend on $A$: with $R''<0$ the left side of \eqref{t5e8-key} is negative and moves with $I$, which
        is what makes $\rho^*$ depend on $A$.
  \item Stating the FOCs without the multiplier on the break-even: the cutoff is a weighted average of
        $\rho_1$ (value) and $\rho_0$ (pledgeable income), not $\rho_1$ alone.
  \item Reading ``shorter maturity'' off the level $(r-\rho^*)I$: since $I$ also falls with $A$, the
        unambiguous statement is per unit of investment.
\end{itemize}

\bookans Matches the book's answer (p.~593): same condition \eqref{t5e8-key} (the book writes its numerator as
$1-r-\int_0^{\rho^*}(\rho^*-\rho)f(\rho)d\rho$, equal to ours by integration by parts) and the same conclusion
(lower $A$ $\Rightarrow$ lower $I$ and $\rho^*$ $\Rightarrow$ more short-term debt). The book states without proof that the
right side of \eqref{t5e8-key} is decreasing in $\rho^*$; \eqref{t5e8-hprime} shows this needs a condition on $F$
(true for the uniform).

\refbib Tirole \S5.3.2 ``Continuum of Liquidity Shocks'', pp.~209--210; \S5.6.1 ``Optimal Claim Maturity'', p.~226;
Exercise 5.7, p.~232; answer p.~593.

% =============================================================================
\exsec{9. The priority game: uncoordinated lending and short-term bias --- Tirole Exercise 5.9, p.~233}
%% BEGIN-STATEMENT 9
\stmtbegin
\textbf{Source:} Tirole, Ch. 5, Exercise 5.9, p.~233 (PDF 246) · Answer in book: yes, p.~594

Exercise 5.9 (the priority game: uncoordinated lending leads to a short-term bias). This chapter,
like Chapters 3 and 4, has assumed that the firm's balance sheet is transparent. In particular, each
investor has perfect knowledge of loans made by other lenders and of the firm's obligations to them.

This exercise argues that uncoordinated lending leads to financing that is too oriented to the short
term. In a nutshell, lenders, by cashing out early, exert a negative externality on other lenders.
Because this externality is not internalized, the resulting financial structure contains too much
short-term debt.

We consider a three-period model: \(t = 0, 1, 2\). The entrepreneur has no cash (\(A = 0\)), is risk
neutral, and is protected by limited liability. At date 0, a fixed investment \(I\) is made. The
project yields a known return \(r > 0\) at date 1, and an uncertain return (\(R\) or 0) at date 2.
Because the point is quite general and does not require credit constraints, we assume away moral
hazard; or, equivalently, the private benefit from misbehaving is 0. The probability of a date 2
success is

\[p + \tau(I_1),\]

where \(I_1\) is the date-1 deepening investment, equal to \(r\) minus the level of short-term debt
repaid to lenders and the date-1 payment to the entrepreneur (the firm does not return to the
capital market at date 1), and \(\tau\) is an increasing and concave function (with
\(\tau'(0) = \infty\)). Assume that \(\tau'(r)R < 1\).

We assume that the entrepreneur cannot engage in ``fraud,'' that is, cannot fail to honor the
short-term debt and, if the project succeeds at date 2, the long-term debt. By contrast, obligations
to lenders, and in particular \(I_1\), cannot be verified as the firm's balance sheet is opaque.

\begin{enumerate}
\def\labelenumi{(\roman{enumi})}
\item
  Derive the first-best investment \(I_1^*\). Show how this allocation can be implemented by a
  mixture of short- and long-term debt (note that in this model without moral hazard the structure
  of compensation for the entrepreneur exhibits a degree of indeterminacy).
\item
  Assume that \(r - I_1^* < I\) (creditors must hold long-term debt). Suppose next that financing is
  not transparent. Start from the first-best solution, with a large number (a continuum of mass 1)
  of lenders, with the representative lender owning short-term claim \(r_l\) and contingent
  long-term claim \(R_l\) on the firm.
\end{enumerate}

Show that the entrepreneur has an incentive to secretly collude with any lender to increase the
latter's short-term claim in exchange for a smaller long-term claim.

Given the constraint that financing is provided by many lenders and that the latter do not observe
each other's contracts, is the indeterminacy mentioned in question (i) resolved?
\stmtend
%% END-STATEMENT

\asked (i) Find the first-best date-1 deepening investment $I_1^*$ and show that a mix of short- and
long-term debt implements it, with some indeterminacy in how the claims (and the entrepreneur's pay)
are split. (ii) With an opaque balance sheet and many lenders, show that the entrepreneur gains by
secretly swapping long-term claim for short-term claim with one lender, and say whether this pins down
the split left open in (i).

\approach Without moral hazard the borrower captures the NPV, so (i) is a plain optimisation plus
counting equations against unknowns. In (ii) the date-1 payout to one lender is paid out of the
reinvestment fund, lowering the success probability of \emph{every} long-term claim: an externality
that neither party to the side deal internalises (a common-agency, ``priority'' problem). Only
collusion-proof structures survive, which is the lecture's point that maturity is chosen to discipline
what the insider can do with unverifiable cash.

\notation
\begin{itemize}
  \item $I$: fixed date-0 investment; $A=0$; no moral hazard (private benefit 0); risk-free rate 0.
  \item $r>0$: known date-1 return; $R$ or 0: date-2 return, success probability $p+\tau(I_1)$.
  \item $I_1\ge 0$: date-1 deepening investment; $\tau$ increasing, concave, $\tau'(0)=\infty$, $\tau'(r)R<1$.
  \item $r_l,r_b\ge 0$: date-1 payments to lenders (short-term debt) and to the entrepreneur;
        $R_l,R_b\ge0$: success-contingent date-2 payments, $R_l+R_b=R$.
  \item $I_1^*$: first-best deepening investment; $\delta r_l,\delta R_l>0$: sizes of a secret deal.
\end{itemize}

\step{1}{(i) First best} The first best maximises total surplus over $I_1$. The NPV is date-1 cash not
reinvested plus expected date-2 income minus the initial investment:
\begin{align}
  \text{NPV}(I_1) &= r - I_1 + [p+\tau(I_1)]R - I. \label{t5e9-npv}
\end{align}
Its derivatives are $\text{NPV}'(I_1)=-1+\tau'(I_1)R$ and $\text{NPV}''(I_1)=\tau''(I_1)R\le0$, so it is
concave and the first-order condition $\text{NPV}'(I_1^*)=0$ is sufficient:
\begin{align}
  -1+\tau'(I_1^*)\,R &= 0 \quad\Longleftrightarrow\quad \tau'(I_1^*)\,R = 1. \label{t5e9-fb}
\end{align}
$\tau'(0)=\infty$ gives $I_1^*>0$; $\tau'(r)R<1=\tau'(I_1^*)R$ and $\tau'$ decreasing give $I_1^*<r$, so
$r-I_1^*>0$ can be paid out at date 1.

\step{2}{(i) Implementation and indeterminacy} With transparent books the contract fixes the date-1
payouts, and hence $I_1=r-r_l-r_b$. The first best requires
\begin{align}
  r_b+r_l+I_1^* &= r, \label{t5e9-c1}\\
  R_b+R_l &= R, \label{t5e9-c2}\\
  I &= r_l+[p+\tau(I_1^*)]R_l, \label{t5e9-c3}
\end{align}
the last being the lenders' breakeven condition. Four unknowns $(r_b,r_l,R_b,R_l)$ and three equations: one
degree of freedom. For instance (when $r-I_1^*<I$), maximal short-term debt $r_b=0$, $r_l=r-I_1^*$ with
$R_l=\frac{I-(r-I_1^*)}{p+\tau(I_1^*)}$; or pure long-term debt $r_l=0$, $R_l=\frac{I}{p+\tau(I_1^*)}$ (if this is $\le R$),
with the entrepreneur taking $r_b=r-I_1^*$. Without moral hazard the entrepreneur gets the first-best NPV in all
cases, so the split of her pay between $r_b$ and $R_b$, and of the lenders' claims between short and long
term, is undetermined.

\answer{(i)} $I_1^*$ solves $\tau'(I_1^*)R=1$, with $0<I_1^*<r$; any $(r_b,r_l,R_b,R_l)\ge0$ with
$r_b+r_l=r-I_1^*$, $R_b+R_l=R$ and $r_l+[p+\tau(I_1^*)]R_l=I$ implements it --- one degree of freedom
in the short-/long-term split.

\step{3}{(ii) The secret deal} With $r-I_1^*<I$, \eqref{t5e9-c3} and $r_l\le r-I_1^*$ force $R_l>0$, hence $R_b<R$
by \eqref{t5e9-c2}. Starting from any first-best structure, the entrepreneur offers one lender $+\delta r_l$ in short-term
claim against the transfer to herself of $\delta R_l$ of his long-term claim. The lender is atomistic, takes $I_1$ as given,
and accepts iff
\begin{align}
  \delta r_l &\ge [p+\tau(I_1)]\,\delta R_l. \label{t5e9-acc}
\end{align}
Date-1 cash is fixed and the other claims cannot be defrauded, so the extra payout comes out of deepening investment:
\begin{align}
  \delta I_1 &= -\delta r_l. \label{t5e9-dI1}
\end{align}
The entrepreneur's utility is $U_b=r_b+[p+\tau(I_1)]R_b$; in the deal $r_b$ is unchanged and she
receives the long-term claim the lender gives up, $\delta R_b=\delta R_l$. To first order (product rule),
\begin{align}
  \delta U_b &= \tau'(I_1)\,(\delta I_1)\,R_b + [p+\tau(I_1)]\,(\delta R_l). \nonumber
\end{align}
Substituting \eqref{t5e9-dI1} and \eqref{t5e9-acc} at equality ($\delta R_l=\delta r_l/[p+\tau(I_1)]$),
\begin{align}
  \delta U_b &= \tau'(I_1)\,(-\delta r_l)\,R_b + [p+\tau(I_1)]\Big(\frac{\delta r_l}{p+\tau(I_1)}\Big) \nonumber\\
             &= \big[1-\tau'(I_1)R_b\big]\,\delta r_l. \label{t5e9-dU}
\end{align}
At $I_1=I_1^*$, \eqref{t5e9-fb} gives $\tau'(I_1^*)=1/R$, so $\tau'(I_1^*)R_b=R_b/R<1$ and
$\delta U_b=(1-R_b/R)\,\delta r_l=\frac{R-R_b}{R}\,\delta r_l=(R_l/R)\,\delta r_l>0$: the deal pays. The gain is
exactly the loss $\tau'(I_1^*)R_l\,\delta r_l=(R_l/R)\delta r_l$ of the other long-term claimholders, whose success probability
falls; total value is unchanged to first order because $I_1^*$ maximises \eqref{t5e9-npv}.

\step{4}{(ii) The indeterminacy is resolved} Lenders anticipate such deals, so only collusion-proof structures are
equilibria. By \eqref{t5e9-dU} no deal pays iff
\begin{align}
  \tau'(I_1)\,R_b &\ge 1. \label{t5e9-cp}
\end{align}
The entrepreneur bears the NPV loss ex ante, and \eqref{t5e9-npv} rises with $I_1$ below $I_1^*$; since $\tau'$ falls with $I_1$,
she sets $I_1$ as high as \eqref{t5e9-cp} allows, so \eqref{t5e9-cp} binds: $\tau'(I_1)R_b=1$. Because $R_b<R$,
$\tau'(I_1)>\tau'(I_1^*)$ and $I_1<I_1^*$. A larger $R_b$ permits a larger $I_1$, so $R_b$ must be the largest compatible with
breakeven, which requires paying her nothing at date 1, $r_b=0$, and all date-1 free cash to lenders, $r_l=r-I_1$:
\begin{align}
  I &= r-I_1+[p+\tau(I_1)](R-R_b), \label{t5e9-be}\\
  \tau'(I_1)\,R_b &= 1. \label{t5e9-cpb}
\end{align}
These two equations pin down $(I_1,R_b)$, hence $r_l=r-I_1$ and $R_l=R-R_b$. So yes, opacity resolves the
indeterminacy: the entrepreneur is paid only long-term ($r_b=0$), lenders hold more short-term debt
($r-I_1>r-I_1^*$), and deepening investment is below the first best.

\answer{(ii)} At the first best a secret swap gives $\delta U_b=[1-\tau'(I_1^*)R_b]\delta r_l=(R_l/R)\delta r_l>0$.
Yes, the indeterminacy is resolved: $r_b=0$, $r_l=r-I_1$, and $(I_1,R_b)$ solve $\tau'(I_1)R_b=1$ and
$I=r-I_1+[p+\tau(I_1)](R-R_b)$, with $I_1<I_1^*$ (a short-term bias).

\intuition A lender who cashes out early is paid out of the firm's reinvestment fund, and the cost (a lower success
probability) falls on every long-term claimant. Neither the colluding lender nor the entrepreneur internalises the
loss of the other lenders, so with opaque books maturity is biased toward the short term. The entrepreneur can only
commit not to strip the firm by holding a large long-term stake herself, and $\tau'(I_1)R_b=1$ shows how large.

\pitfalls
\begin{itemize}
  \item Letting the colluding lender internalise the fall in $I_1$: he is atomistic, takes $I_1$ as given
        and accepts any deal with $\delta r_l\ge[p+\tau(I_1)]\delta R_l$.
  \item Counting the deal as value-creating: at $I_1^*$ total value is unchanged to first order; the
        entrepreneur's gain is the other long-term holders' loss.
  \item Concluding that ``more long-term debt'' fixes the bias: what deters stripping is the
        entrepreneur's own long-term stake $R_b$, which is why $r_b=0$.
\end{itemize}

\bookans Matches the book's answer (p.~594): $\tau'(I_1^*)R=1$; one degree of freedom; $\delta U_b=[1-\tau'(I_1)R_b]\delta r_l>0$;
resolution with $r_b=0$, $I=r-I_1+[p+\tau(I_1)](R-R_b)$ and $\tau'(I_1)R_b=1$, $I_1<I_1^*$.

\refbib Tirole \S5.6.1 ``Optimal Claim Maturity'', p.~226; Exercise 5.9, p.~233; answer p.~594.

% =============================================================================
\exsec{10. One vs.\ two creditors: strategic and liquidity default --- Amaro de Matos 5.2.5 \#2, p.~148}
%% BEGIN-STATEMENT 10
\stmtbegin
\textbf{Source:} Amaro de Matos, Ch. 5, Exercise 5.2.5 \#2, p.~148 (PDF 166) · Answer in book: no
(p.~268 only cites pages 10--13 of Bolton and Scharfstein 1996)

Show that two creditors are paid more in the case of strategic default than one creditor, but are
paid less in the case of liquidity default.

\paragraph{Context needed to solve it}

\emph{From Amaro de Matos §5.2.1 ``Extending the Model'', pp.~134--137 (PDF 152--155):}

As before, all the agents are risk neutral and the risk-free rate is assumed to be zero. The time
schedule refers to a two-period model and is characterized as follows: At time \(t = 0\) an amount
\(K\) is invested in the project. The entrepreneur, as before, is assumed to have no initial wealth.
In this specification of the model, the state of nature is characterized by the cash flow at
\(t = 1\). It may be \(\theta_g = x\) with probability \(q\) and \(\theta_b = 0\) with probability
\(1 - q\). The cash flow is observable for both the investor and the entrepreneur, but not for any
third party, and is thus not verifiable. Therefore, as before, a contract between the investor and
the entrepreneur is incomplete.

The manager is assumed to have the ability to divert cash flow without the control of the investor.
In that case a zero payoff would be declared at \(t = 1\) when the actual payoff was \(x\). If the
declared cash flow at \(t = 1\) is \(x\), the contract states a payment of \(R_x\) to the
creditholder and, with probability \(\beta_x\), the right of the creditholder to liquidate a part of
the assets that is worth \(L_x\). If the cash flow is null, the values are, respectively,
\(R_\circ\), \(\beta_\circ\), and \(L_\circ\).

Notice that taking \(\beta_x = 0\), \(\beta_\circ = 1\), and \(L_x = L_\circ\), being equal to the
value of the assets of the project at time \(t = 1\), this reduces to a usual debt contract with
face value \(R_x < x\) and one-period maturity. The issue in designing a contract is to find the
optimal values for the preceding parameters. Therefore, it is possible in this simple setting to
analyze the conditions under which a debt contract would result as optimal.

In the second period, two things may happen. Either the manager continues the project or the project
is liquidated, in the sense that it is separated from the manager. If the manager's activity
continues, the cash flow at time \(t = 2\) is \(y\). In the case of no liquidation, the project thus
has a net present value of \(qx + y - K\).

The expected return to the firm is then written as

\[\pi_f = q[x - R_x + (1 - \beta_x)y] + (1 - q)[-R_\circ + (1 - \beta_\circ)y]. \qquad (5.3)\]

On the other hand, the investor expects to get

\[\pi_i = q[R_x + \beta_x L_x] + (1 - q)[R_\circ + \beta_\circ L_\circ] - K. \qquad (5.4)\]

Suppose for a moment that at time \(t = 1\) the cash flow is \(x\). The expected return to the firm
would then be \(x - R_x + (1 - \beta_x)y\), decomposed as follows: The cash flow \(x\) goes to the
entrepreneur, who has to pay \(R_x\) to the investor. With probability \(\beta_x\) the project is
liquidated at date \(t = 1\), leading to a continuation value of zero, and with probability
\(1 - \beta_x\) the project continues until time \(t = 2\), generating a cash flow \(y\) at that
point. If, however, the manager hides from the outer investor the cash flow \(x\), then only
\(R_\circ\) will be paid at time \(t = 1\), and liquidation occurs with probability \(\beta_\circ\),
leaving to the firm a possible payoff \(S\), not necessarily zero, as before. This happens because
the liquidation value \(L_\circ\) does not correspond to the true value of the assets in that state
of nature. In that case the expected return to the firm is
\(x - R_\circ + \beta_\circ S + (1 - \beta_\circ)y\).

For the manager to be induced to pay the due \(R_x\) when the cash flow is \(x\), one must therefore
impose

\[x - R_x + (1 - \beta_x)y \geq x - R_\circ + \beta_\circ S + (1 - \beta_\circ)y. \qquad (5.5)\]

Noticing that \(R_\circ = \beta_x = 0\) is optimal in the sense that these values maximize the
profit to the firm and facilitates the verification of (5.5), the preceding inequality reads

\[R_x \leq \beta_\circ(y - S).\]

Among all these possible values for \(R_x\) that ensure a trustworthy behavior from the manager of
the project, any investor would maximize its expected return by choosing the greatest value. Thus,
at the optimum, this restriction is binding. Substituting the optimal values of the parameters in
expression (5.4) for the expected result of the investor, one gets

\[\pi_i = \beta_\circ[q(y - S) + (1 - q)L_\circ] - K. \qquad (5.6)\]

Using the optimal values of the parameters, namely those minimizing the profit to the investor, in
the profit of the firm (5.3), it follows that

\[\pi_f = (qx + y - K) - (1 - q)\beta_\circ(y - L_\circ). \qquad (5.7)\]

This result can be interpreted as the net present value of the project in the absence of liquidation
minus a cost due to contract incompleteness.

For the profit in (5.6) to be positive, ensuring the minimum condition of feasibility of the
contract, one must have

\[1 \geq \beta_\circ \geq \frac{K}{q(y - S) + (1 - q)L_\circ},\]

where the first inequality comes from the fact that \(\beta_\circ\) is a probability and the second
inequality comes from the condition \(\pi_i \geq 0\). Notice that \(\pi_f\) is decreasing in
\(\beta_\circ\) whereas \(\pi_i\) is monotonically increasing in that parameter. Maximizing the
expected return to the firm conditioned to a positive expected return to the investor defines the
optimal value of \(\beta_\circ\) as the lowest possible. Thus

\[\beta_\circ = \frac{K}{q(y - S) + (1 - q)L_\circ}.\]

Both \(S\) and \(L_\circ\) (and thus also \(\beta_\circ\)) are still to be specified. There is no
need to specify \(L_x\), since the main feature of the optimal contract just described is that there
is no liquidation in the good state. Liquidation occurs only, if at all, in the case of a poor
performance. In order to understand what determines \(S\) and \(L_\circ\) some additional
assumptions are required.

It is assumed that the entrepreneur is not able to continue managing the project in the case of
bankruptcy after a zero cash flow. This is the case of a \emph{liquidity default}. In that case, the
creditor may sell the position to a third party who obtains at \(t = 2\) a cash flow \(\alpha y\),
with \(\alpha \in [0, 1]\), meaning that the incoming manager cannot get better results out of the
assets than the manager who left the firm.

In the case of bankruptcy after a nonzero cash flow, the original manager will negotiate the buyback
of the firm with the creditholder, sharing the second-period return \(y\). In that situation,
\(S = y/2\). This is the case of a \emph{strategic default}.

It is also assumed that, in the case of liquidity default, the buyer of the assets incurs an
informational cost \(c \in [0, \bar c]\). This cost is modeled as a uniformly distributed random
variable. Therefore, by a negotiating process similar to the preceding, the remaining value of the
assets, \(\alpha y\), is equally divided between the creditor and the buyer. The buyer then
negotiates only if \(c \leq \alpha y/2\). Thus, assuming that \(\alpha y/2 \leq \bar c\), the
probability that this happens is

\[\Pr\left(c \leq \frac{\alpha y}{2}\right) = \frac{\alpha y}{2\bar c}.\]

Therefore, the expected value of the liquidation payoff in the case of a zero cash flow at \(t = 1\)
is

\[L_\circ = \frac{\alpha^2 y^2}{4\bar c},\]

and the optimal \(\beta_\circ\) reads

\[\beta_\circ = \frac{K}{q\frac{y}{2} + (1 - q)\frac{\alpha^2 y^2}{4\bar c}}. \qquad (5.8)\]

{[}\ldots{]}

\emph{From Amaro de Matos §5.2.2 ``The Case of Many Creditholders'', pp.~137--138 (PDF 155--156):}

Bolton and Scharfstein (1996) bring to this discussion a related issue: The crucial element is the
number, nature, and type of creditors---in other words, whether credit is private or public, whether
it is granted by banks or obtained by issuing bonds in the marketplace, and also the nature of the
governance structure among the creditors.

In this setting assume that at time \(t = 0\) the amount \(K\) is invested to acquire the assets
\(A\) and \(B\). These are the assets that constitute the investment project. If there are two
creditors, assume that each of the assets \(A\) and \(B\) provides collateral for each of them.
Denoting by \(y^A\) and \(y^B\) the cash flows for assets \(A\) and \(B\), as managed separately,
one defines \(\Delta = y - y^A - y^B\), which is assumed to be positive. This means that the cash
flow generated by both assets is assumed to be greater than the sum of the cash flow generated by
each of the assets individually. It is then possible to show that the expected value of liquidation
at time \(t = 1\) is given by

\[\overline{L}_\circ = \frac{\alpha^2 y^2}{4\bar c} - \frac{\alpha^2 \Delta^2}{36\bar c} < \frac{\alpha^2 y^2}{4\bar c} = L_\circ.\]

In other words, in the case of two creditors, the expected value of liquidation when there is a zero
cash flow at \(t = 1\) is lower than the value expected with only one creditor. Moreover, in the
state with probability \(q\),

\[S = \frac{y}{2} - \frac{\Delta}{6},\]

leading to an optimal value of \(\beta_\circ\) with a denominator given by

\[q\left(\frac{y}{2} - \frac{\Delta}{6}\right) + (1 - q)\left(\frac{\alpha^2 y^2}{4\bar c} - \frac{\alpha^2 \Delta^2}{36\bar c}\right),\]

which can be greater or less than the similar value found before with only one creditor.
{[}\ldots{]}
\stmtend
%% END-STATEMENT

\asked \emph{(strategic)}: show that when the manager defaults strategically and buys the firm back,
the creditors collect more if the assets are split between two creditors than if one creditor holds
them all. \emph{(liquidity)}: show that when the default is a genuine liquidity default and the assets are
sold to an outside buyer, two creditors collect less in expectation than one.

\approach This is Bolton and Scharfstein (1996) inside the strategic-default model of \S5.2.1: the
investor's recovery is $y-S$ after a strategic default and $L_\circ$ after a liquidity default, and both
come from bargaining. With two creditors each holds an asset the other party needs to realise the
synergy $\Delta$, so each has hold-up power. Splitting the surplus by the Shapley value (the three-player
version of the equal split $S=y/2$) gives both formulas on p.~138; comparing them with $y/2$ and
$\alpha^2y^2/(4\bar c)$ answers the two parts.

\notation
\begin{itemize}
  \item $K$: investment; $q$: probability of cash flow $x$ at $t=1$ (else 0); $y$: $t=2$ cash flow if the manager continues.
  \item $S$: the manager's payoff when she strategically defaults and buys the firm back; $y-S$: what the creditor(s) get.
  \item $\alpha\in[0,1]$: efficiency of an outside buyer ($\alpha y$); $c\sim U[0,\bar c]$: buyer's information cost.
  \item $L_\circ$ / $\overline L_\circ$: expected liquidation payoff to creditors after a zero cash flow, with one / two creditors.
  \item $y^A,y^B\ge0$: cash flows of assets $A$ and $B$ run separately; $\Delta=y-y^A-y^B\in(0,y]$: synergy.
  \item $\beta_\circ$: liquidation probability after a zero report, $\beta_\circ=K/[q(y-S)+(1-q)L_\circ]$ from (5.6).
\end{itemize}
\assume With two creditors, creditor 1 holds asset $A$ and creditor 2 asset $B$. Multilateral bargaining follows the
Shapley value, the three-player version of the equal split used for one creditor; disagreement payoffs are 0
(as behind $S=y/2$ in the one-creditor case). The book only cites Bolton and Scharfstein (1996), pp.~10--13;
the derivation below is ours, and it reproduces both formulas printed on p.~138.

\step{1}{One creditor (benchmark)} We first record the one-creditor recoveries the comparison is made
against. Manager and creditor split $y$ equally: $S=y/2$, so in a strategic default the creditor gets
\begin{align}
  y-S &= \frac{y}{2}. \label{t5e10-s1}
\end{align}
After a zero cash flow the buyer gets $\alpha y/2$, buys iff $c\le\alpha y/2$ (probability $\alpha y/(2\bar c)$), and the creditor gets $\alpha y/2$:
\begin{align}
  L_\circ &= \Big(\frac{\alpha y}{2\bar c}\Big)\Big(\frac{\alpha y}{2}\Big) = \frac{\alpha^2y^2}{4\bar c}. \label{t5e10-L1}
\end{align}

\step{2}{Two creditors, strategic default} Players: manager $M$, creditors 1 and 2. Worth of a coalition: $y$ if it contains $M$ and both assets,
$y^A$ ($y^B$) if it contains $M$ and only asset $A$ ($B$), and 0 without $M$. Over the $3!=6$ orderings, $M$'s marginal contribution is 0
when she comes first (orderings $M12$, $M21$), $y^A$ in ordering $1M2$ and $y^B$ in $2M1$ (she joins a
coalition holding one asset, worth 0 without her), and $y$ when last ($12M$, $21M$). Averaging over the six
orderings and using $y^A+y^B=y-\Delta$:
\begin{align}
  S &= \frac{0+0+y^A+y^B+y+y}{6} = \frac{2y+y^A+y^B}{6} \nonumber\\
    &= \frac{2y+(y-\Delta)}{6} = \frac{3y-\Delta}{6} \nonumber\\
    &= \frac{y}{2}-\frac{\Delta}{6}, \label{t5e10-S2}
\end{align}
which is the book's $S$. The two creditors together get
\begin{align}
  y-S &= y-\Big(\frac{y}{2}-\frac{\Delta}{6}\Big) \nonumber\\
      &= \frac{y}{2}+\frac{\Delta}{6} \;>\; \frac{y}{2}. \label{t5e10-s2}
\end{align}
Comparing with \eqref{t5e10-s1}: two creditors are paid more in a strategic default, by $\Delta/6$. Each creditor can hold up the synergy $\Delta$,
so the manager keeps less when she buys the firm back.

\answer{(strategic)} Two creditors get $y-S=\dfrac y2+\dfrac\Delta6>\dfrac y2$ (one creditor): more, by
$\Delta/6$.

\step{3}{Two creditors, liquidity default} Replace $M$ by the outside buyer, whose coalition worths are $\alpha y$, $\alpha y^A$, $\alpha y^B$
(the cost $c$ is sunk and not bargained over, as in Step 1). By \eqref{t5e10-S2} the buyer gets $\alpha(y/2-\Delta/6)$, buys iff
$c\le\alpha(y/2-\Delta/6)$, which for $c\sim U[0,\bar c]$ has probability $\alpha(y/2-\Delta/6)/\bar c$
(assume $\alpha(y/2-\Delta/6)\le\bar c$), and the creditors share $\alpha y-\alpha(y/2-\Delta/6)=\alpha(y/2+\Delta/6)$.
Expected liquidation payoff = probability of a sale times the creditors' share:
\begin{align}
  \overline L_\circ &= \Big(\frac{\alpha(3y-\Delta)}{6\bar c}\Big)\Big(\frac{\alpha(3y+\Delta)}{6}\Big) \nonumber\\
                    &= \frac{\alpha^2(3y-\Delta)(3y+\Delta)}{36\bar c} = \frac{\alpha^2(9y^2-\Delta^2)}{36\bar c} \nonumber\\
                    &= \frac{9\alpha^2y^2}{36\bar c}-\frac{\alpha^2\Delta^2}{36\bar c}
                     = \frac{\alpha^2y^2}{4\bar c}-\frac{\alpha^2\Delta^2}{36\bar c}, \label{t5e10-L2}
\end{align}
the book's formula. By \eqref{t5e10-L1}, $\overline L_\circ=L_\circ-\alpha^2\Delta^2/(36\bar c)<L_\circ$: two creditors are paid less in a liquidity
default. The creditors' extra share does not compensate the lower probability that a buyer shows up, which falls
from $\alpha y/(2\bar c)$ to $\alpha(y/2-\Delta/6)/\bar c$.

\answer{(liquidity)} Two creditors expect $\overline L_\circ=\dfrac{\alpha^2y^2}{4\bar c}-\dfrac{\alpha^2\Delta^2}{36\bar c}
<L_\circ=\dfrac{\alpha^2y^2}{4\bar c}$: less, by $\alpha^2\Delta^2/(36\bar c)$.

\step{4}{The $\beta_\circ$ denominator and the printed sign} By (5.6) the denominator of $\beta_\circ$ is $q(y-S)+(1-q)L$. With two creditors,
substituting \eqref{t5e10-s2} and \eqref{t5e10-L2},
\begin{align}
  D_2 &= q\Big(\frac{y}{2}+\frac{\Delta}{6}\Big) + (1-q)\Big(\frac{\alpha^2y^2}{4\bar c}-\frac{\alpha^2\Delta^2}{36\bar c}\Big). \label{t5e10-D2}
\end{align}
The book prints $q(\frac y2-\frac\Delta6)$ in the first term (p.~138), i.e.\ $qS$ instead of $q(y-S)$ (the two coincide only when
$S=y/2$, the one-creditor case). That is a sign error: with the printed term the two-creditor denominator would be below the one-creditor
one for every $q$, contradicting the same page's remarks that it ``can be greater or less'' and that two creditors are paid more in strategic default.
With \eqref{t5e10-D2} and $D_1=q\frac y2+(1-q)L_\circ$,
\begin{align}
  D_2-D_1 &= q\Big(\frac y2+\frac\Delta6-\frac y2\Big)+(1-q)\big(\overline L_\circ-L_\circ\big) \nonumber\\
          &= q\,\frac{\Delta}{6}-(1-q)\,\frac{\alpha^2\Delta^2}{36\bar c}. \label{t5e10-dD}
\end{align}
It is positive (two creditors liquidate less often) iff, multiplying by $36\bar c/\Delta>0$,
$6\bar c\,q>(1-q)\alpha^2\Delta$, i.e.\ $q(6\bar c+\alpha^2\Delta)>\alpha^2\Delta$, i.e.\
$q>\bar q\equiv\frac{\alpha^2\Delta}{6\bar c+\alpha^2\Delta}$: when bankruptcy risk is low, as the book concludes.

\step{5}{Numerical illustration (ours)} Take $y=10$, $y^A=4$, $y^B=3$ ($\Delta=10-4-3=3$), $\alpha=0.8$, $\bar c=6$, $q=0.6$, $K=2$. Then
$S=10/2-3/6=5-0.5=4.5$ and $y-S=5.5$ versus $5$ with one creditor. In liquidity default the buyer gets $0.8(4.5)=3.6$, buys with probability
$3.6/6=0.6$ (versus $4/6=0.666667$), and the creditors get $0.8(5.5)=4.4$:
\begin{align}
  \overline L_\circ &= (0.6)(4.4) = 2.64 \;<\; L_\circ=\frac{(0.64)(100)}{24}=2.666667. \nonumber
\end{align}
Denominators: $D_1=(0.6)(5)+(0.4)(2.666667)=4.066667$ and $D_2=(0.6)(5.5)+(0.4)(2.64)=4.356$, so $\beta_\circ$ falls from
$2/4.066667=0.491803$ to $2/4.356=0.459137$. The misprinted term would give $(0.6)(4.5)+(0.4)(2.64)=3.756$ and $\beta_\circ=0.532481$.
Here $\bar q=1.92/37.92=0.050633$; at $q=0.02$, \eqref{t5e10-dD} gives $-0.016133<0$ and one creditor is better.

\intuition Splitting the assets among creditors gives each a veto over the synergy $\Delta$. This hold-up power makes strategic default
costly for the manager, so more can be credibly repaid (good when default is mostly strategic, high $q$). The same hold-up power scares
off outside buyers when the firm genuinely fails, lowering the liquidation value (bad when default is mostly for liquidity, low $q$).

\pitfalls
\begin{itemize}
  \item Comparing $S$ instead of $y-S$: $S$ is what the \emph{manager} keeps; the creditors' recovery is
        $y-S$, which rises with two creditors.
  \item Keeping the buyer's probability of purchase fixed at $\alpha y/(2\bar c)$: the buyer's share
        falls to $\alpha(y/2-\Delta/6)$, so fewer buyers show up, and that loss dominates.
  \item Copying the printed two-creditor denominator $q(y/2-\Delta/6)+\dots$: its first term must be
        $q(y-S)=q(y/2+\Delta/6)$.
\end{itemize}

\bookans The book (p.~268) only cites Bolton and Scharfstein (1996), pp.~10--13; derived here. Our $S$ and $\overline L_\circ$ match p.~138.
The printed two-creditor denominator on p.~138 has a sign error: its first term should be $q(y/2+\Delta/6)=q(y-S)$, not $q(y/2-\Delta/6)$.

\refbib Amaro de Matos \S5.2.1 ``Extending the Model'', pp.~134--137; \S5.2.2 ``The Case of Many Creditholders'', pp.~137--138; exercise p.~148.

% =============================================================================
\exsec{11. When the good entrepreneur prefers liquidating debt --- Amaro de Matos 5.2.5 \#4, p.~148}
%% BEGIN-STATEMENT 11
\stmtbegin
\textbf{Source:} Amaro de Matos, Ch. 5, Exercise 5.2.5 \#4, p.~148 (PDF 166) · Answer in book: no
(p.~268 only cites the appendix of Diamond 1991, pp.~732--733)

\begin{enumerate}
\def\labelenumi{\arabic{enumi}.}
\setcounter{enumi}{3}
\tightlist
\item
  Prove that inequality (5.12) provides a sufficient condition for the good entrepreneur to prefer
  the debt structure that leads to liquidation under the conditions of the lemma to which it refers.
\end{enumerate}

\paragraph{Context needed to solve it}

\emph{From Amaro de Matos §5.2.3 ``The Choice of the Duration'', pp.~139--141 (PDF 157--159):}

A model developed by Diamond (1991, 1993) illustrates in a very simple way the differences between
short-term and long-term debt financing. The model is based on information asymmetry and credit
rating.

Regaining the notation of Aghion and Bolton, consider an entrepreneur with a two-period project who
needs funds for an initial investment \(K = 1\). The project may be of one of two types: good or
bad. At time \(t = 2\) a good project has a payoff \(X > 0\), whereas a bad project pays that amount
with a probability \(\pi\), and pays zero with probability \(1 - \pi\). Let \(\theta\) denote the
type of project, \(\theta_g\) referring to a good project and \(\theta_b\) referring to a bad
project. The quality of the project is known to the entrepreneur but not to the investor. Given a
project, as of date \(t = 0\) the investor assigns a probability \(f\) that \(\theta = \theta_g\).
This probability is referred to as a \emph{credit rating} and reflects the lender's information
about the project. It follows that, as of time 0, the probability \(q\) that a borrower pays the
loan at date \(t = 2\) is given by

\[q \equiv f + \pi(1 - f) = \pi + f(1 - \pi). \qquad (5.9)\]

At time \(t = 1\) information about the project becomes available to the investor. This new
information is not verifiable, so it makes no sense to write contracts contingent upon that
information, since they cannot be enforced. However, the investors update their beliefs about the
quality of the project and the credit rating will be either upgraded or downgraded.

Under this revision of the credit rating, let \(f^u\) denote the probability that
\(\theta = \theta_g\) given an upgrade and let \(f^d\) denote the probability that
\(\theta = \theta_g\) given a downgrade. In that case, the probability that an entrepreneur who
receives a downgrade repays at time \(t = 2\) is

\[q^d = \pi + f^d(1 - \pi) > f^d. \qquad (5.10)\]

Moreover, it is here assumed that

\[f^d < f < f^u,\]

meaning that the probability that the project is good is greater given an upgrade than given a
downgrade. Furthermore, the simplifying assumption \(f^u = 1\) is also made. Thus, given an upgrade
at date \(t = 1\), it is certain that the project is good. However, it is not certain that all
projects of good quality are upgraded. There is room for one such project to be downgraded. Let
\(e\) denote the probability of such an event.

All projects are also assumed to produce a nonassignable control rent \(C\) if the management has
control at time \(t = 2\). Also, if at time \(t = 1\) the manager defaults on short-term debt,
lenders then have the right to remove the management from control, selling the assets for a
liquidation value \(L < K = 1\). All projects may be liquidated at that point in time, under those
conditions. If the project is liquidated, the entrepreneur is prevented from consumption of his/her
present or future cash flows and control rents. In this case, an optimal contract is a debt contract
enforced by the right to liquidate the project if the debt is not repaid in full.

Just as in the case of the borrowers, lenders are assumed be risk neutral. Moreover, lenders have
alternative investments at \(t = 0\) that provide a return \(R\) at \(t = 2\). The condition
\(L < X/R\) ensures that a good project yields a higher return when not liquidated.

A project may be financed with either long-term or short-term debt. In the context of this model,
long-term debt refers to a contract issued at \(t = 0\) with no refinancing at date \(t = 1\) and
maturing at \(t = 2\). If one such instrument provides \(K = 1\) at time \(t = 0\) with face value
\(\rho \le X\), the lender expects to receive at maturity \(\rho[\pi + f(1 - \pi)] = R^2\). Notice
that if \(\rho > X\), a borrower with credit rating \(f\) will not be able to provide the lender
with the required rate of return \(R^2\). Now, if the project is upgraded at time \(t = 1\), the
market value of the debt contract at that point in time becomes \(V_l^u = \rho/R\). Similarly, if
the project is downgraded, the market value of debt at time \(t = 1\) becomes \(V_l^d = q^d\rho/R\).

Short-term debt, on the other hand, is debt financed at \(t = 0\) with face value \(r_1\) and
maturity at time \(t = 1\). If short-term borrowers are upgraded at \(t = 1\), their debt contract
at that point in time is worth \(V_s^u = r_1\) and they may refinance their project by issuing a new
short-term debt contract with maturity at \(t = 2\) and face value \(r_1 R\) that will be paid with
certainty and guarantees a return of \(R\) per unit invested by the lender. If short-term borrowers
are downgraded, however, in principle they may refinance their debt with a new short-term contract
only if \(q^d X/R > r_1\), since \(X\) is the maximum that a borrower can pay at \(t = 2\). If that
is the case, at \(t = 1\) the market value of the short-term debt issued at \(t = 0\) would be
\(V_s^d = r_1\) if the borrower pays, possibly through the refinancement mechanism. But what happens
if \(q^d X/R < r_1\)? In that case there is still room for a refinancement provided that the
liquidation value \(L\) is less than the expected value to the lender from a refinancement,
\(q^d X/R\). Under a downgrading, the time \(t = 1\) value of the short-term debt is thus

\[V_s^d = \min\left[r_1, \max(L, q^d X/R)\right].\]

The face value of short-term debt is established such as to provide the lenders with an expected
return of \(R\) per period. It then follows that

\[r_1 = \frac{R - [1 - (1 - e)f]V_s^d}{(1 - e)f}. \qquad (5.11)\]

\emph{From Amaro de Matos §5.2.3, pp.~141--142 (PDF 159--160):}

Given this model, the point is to characterize under what conditions short-term or long-term
maturities will be offered. The maturity preferred by the entrepreneurs of good projects will be the
only one offered to all borrowers. This happens because if any other maturity is offered, it reveals
a project of bad quality and lenders will not be willing to finance it.

The following result implies that, in the absence of liquidation given a downgrade, the
entrepreneurs of good projects prefer short-term debt:

\textbf{Lemma 5.2.1 (Diamond, 1991).} \emph{Comparing two debt contracts issued at \(t = 0\),
neither of which involves liquidation on a downgrade at \(t = 1\), the one with the higher value at
that point in time given a downgrade, \(V^d\), is preferred by the entrepreneurs of good projects.}

\emph{Proof.} The expected payoff to the entrepreneurs of good projects from debt not yielding
liquidation is

\[e\left(X + C - \frac{RV^d}{q^d}\right) + (1 - e)(X + C - r_1 R).\]

Using expression (5.11) this payoff is rewritten as

\[X + C - \frac{R^2}{f} + RV^d e\left(\frac{1}{f^d} - \frac{1}{q^d}\right).\]

Given (5.10) this expression for the payoff is clearly an increasing function of \(V^d\). ■

\emph{From Amaro de Matos §5.2.3, p.~142 (PDF 160):}

This implies that when there is no liquidation after a downgrade, short-term debt is preferred. The
idea behind this result is simple. If there is no risk of liquidation, in order to reduce their
expected financing costs, borrowers of the good type, who are less likely than average to be
downgraded, prefer to extract larger payments from those who are downgraded. However, if there is
the possibility of liquidation, the increase in the payment extracted from those who get a downgrade
incorporates, for the good entrepreneurs, the risk of losing control rents. In that case the
following result characterizes the trade-off incurred:

\textbf{Lemma 5.2.2 (Diamond, 1991).} \emph{Consider a debt structure that does not imply
liquidation and leads to a value of debt \(V^d\), given a downgrade. Compare it with an alternative
structure that implies liquidation. Entrepreneurs of good projects prefer the latter if and only if
the control rent \(C < \widehat C\), where the upper bound is given by}

\[\widehat C = \frac{LR}{f^d} - X - RV^d\left(\frac{1}{f^d} - \frac{1}{q^d}\right) > 0.\]

It follows that a sufficient condition for the good entrepreneur to prefer liquidation is that
\(C = 0\), representing the absence of control rents. It is also sufficient for the good
entrepreneur to prefer liquidation if that state does not provide sufficiently high control rents
\(C\) to the entrepreneur compared to the proceeds \(L\) to the investor or

\[\frac{q^d X + f^d C}{R} \le L. \qquad (5.12)\]
\stmtend
%% END-STATEMENT

\asked Prove that if $(q^dX+f^dC)/R\le L$, condition (5.12), then the entrepreneur of a good project
prefers a debt structure that liquidates the firm after a downgrade to any structure that does not
(the comparison of Lemma~5.2.2).

\approach Diamond (1991): liquidation after a downgrade costs the good entrepreneur $X+C$ in that state
but lets lenders recover $L$ from the whole downgraded pool, most of which is bad, so the upgraded (good)
firms pay less. Lemma~5.2.2 reduces the comparison to $C<\widehat C$. The proof therefore (1) derives
both payoffs from the break-even (5.11) and Bayes' rule, (2) obtains $U^L-U^N=e(\widehat C-C)$, and
(3) bounds $\widehat C$ from below using the fact that no non-liquidating debt can be worth more than
$q^dX/R$ after a downgrade.

\notation
\begin{itemize}
  \item $X$: $t=2$ payoff (good project surely, bad with probability $\pi$); $C\ge0$: nonassignable control rent; $L$: liquidation value; $R$: gross return per period.
  \item $f$: initial rating; $f^u=1$, $f^d$: posterior of ``good'' after an upgrade / downgrade; $e\in(0,1)$: probability a good project is downgraded.
  \item $P_d$: probability of a downgrade for an unknown type; $q^d=\pi+f^d(1-\pi)$: repayment probability after a downgrade.
  \item $V^d$: $t=1$ value of a non-liquidating debt after a downgrade; $r_1$: face value of short-term debt; $U^N$, $U^L$: good entrepreneur's payoffs under non-liquidating and liquidating debt.
\end{itemize}

\step{1}{Bayes} Both payoffs depend on the downgrade probability, which we link to $f^d$. Since $f^u=1$,
bad projects are always downgraded and good ones with probability $e$, so
\begin{align}
  P_d &= fe+(1-f) = 1-(1-e)f, \label{t5e11-Pd}\\
  f^d &= \frac{fe}{P_d}\quad\Longleftrightarrow\quad \frac{P_d}{f}=\frac{e}{f^d}. \label{t5e11-fd}
\end{align}

\step{2}{Non-liquidating debt} We compute the good entrepreneur's payoff when a downgrade does not
trigger liquidation. After a downgrade, the debt worth $V^d$ at $t=1$ requires a $t=2$ payment
$RV^d/q^d$ (new or existing lenders break even at the pooled repayment probability $q^d$), which the good
firm pays for sure; after an upgrade it pays $r_1R$. (For long-term debt read $r_1=\rho/R$: then the
break-even $P_u(\rho/R)+P_dV_l^d=R$ is (5.11), since $P_u+P_dq^d=q$ and $\rho q=R^2$.) Rearranging (5.11),
$(1-e)f\,r_1=R-P_dV^d$; multiplying by $R/f$ and using \eqref{t5e11-fd},
\begin{align}
  (1-e)\,r_1R &= \frac{R\,(R-P_dV^d)}{f} \nonumber\\
              &= \frac{R^2}{f}-RV^d\,\frac{P_d}{f} \nonumber\\
              &= \frac{R^2}{f}-\frac{eRV^d}{f^d}. \label{t5e11-r1}
\end{align}
So, expanding and substituting \eqref{t5e11-r1},
\begin{align}
  U^N &= e\Big(X+C-\frac{RV^d}{q^d}\Big)+(1-e)(X+C-r_1R) \nonumber\\
      &= (X+C)-\frac{eRV^d}{q^d}-(1-e)\,r_1R \nonumber\\
      &= (X+C)-\frac{eRV^d}{q^d}-\frac{R^2}{f}+\frac{eRV^d}{f^d} \nonumber\\
      &= X+C-\frac{R^2}{f}+eRV^d\Big(\frac{1}{f^d}-\frac{1}{q^d}\Big), \label{t5e11-UN}
\end{align}
the expression in the proof of Lemma 5.2.1.

\step{3}{Liquidating debt} Now a downgraded firm is liquidated and the lenders get $L$, so \eqref{t5e11-r1} holds with $V^d=L$, and the good
entrepreneur gets nothing after a downgrade:
\begin{align}
  U^L &= (1-e)(X+C) - (1-e)\,r_1R \nonumber\\
      &= (1-e)(X+C)-\Big(\frac{R^2}{f}-\frac{eRL}{f^d}\Big) \nonumber\\
      &= (1-e)(X+C)-\frac{R^2}{f}+\frac{eRL}{f^d}. \label{t5e11-UL}
\end{align}

\step{4}{Lemma 5.2.2} We now compare the two. Subtracting \eqref{t5e11-UN} from \eqref{t5e11-UL}, the
$R^2/f$ terms cancel and $(1-e)(X+C)-(X+C)=-e(X+C)$:
\begin{align}
  U^L-U^N &= -e(X+C)+\frac{eRL}{f^d}-eRV^d\Big(\frac{1}{f^d}-\frac{1}{q^d}\Big) \nonumber\\
          &= e\Big[\frac{LR}{f^d}-X-RV^d\Big(\frac{1}{f^d}-\frac{1}{q^d}\Big)-C\Big] \nonumber\\
          &= e\,(\widehat C-C), \label{t5e11-diff}
\end{align}
so (with $e>0$) the good entrepreneur prefers liquidation iff $C<\widehat C$.

\step{5}{A bound on $V^d$} A downgraded firm can pay at most $X$ at $t=2$, and repays with probability $q^d$; so any non-liquidating debt has
\begin{align}
  V^d &\le \frac{q^dX}{R}. \label{t5e11-Vb}
\end{align}
(Long-term: $V_l^d=q^d\rho/R$ with $\rho\le X$. Short-term without liquidation: $V_s^d=r_1\le q^dX/R$, the refinancing condition.)

\step{6}{(5.12) is sufficient} Since $\frac1{f^d}-\frac1{q^d}>0$ by (5.10), $\widehat C$ falls with $V^d$; substituting \eqref{t5e11-Vb} into the definition of $\widehat C$,
\begin{align}
  \widehat C &\ge \frac{LR}{f^d}-X-R\Big(\frac{q^dX}{R}\Big)\Big(\frac{1}{f^d}-\frac{1}{q^d}\Big) \nonumber\\
             &= \frac{LR}{f^d}-X-\frac{q^dX}{f^d}+\frac{q^dX}{q^d} \nonumber\\
             &= \frac{LR}{f^d}-X-\frac{q^dX}{f^d}+X \nonumber\\
             &= \frac{LR-q^dX}{f^d}. \label{t5e11-Cb}
\end{align}
Inequality (5.12), $(q^dX+f^dC)/R\le L$, i.e.\ $q^dX+f^dC\le LR$, is equivalent (subtract $q^dX$, divide by $f^d>0$) to
\begin{align}
  C &\le \frac{LR-q^dX}{f^d}. \label{t5e11-512}
\end{align}
Chaining \eqref{t5e11-512} and \eqref{t5e11-Cb}: $C\le\widehat C$, and by \eqref{t5e11-diff} $U^L\ge U^N$. The preference is strict unless both
\eqref{t5e11-512} and \eqref{t5e11-Vb} hold with equality (a knife-edge with $\rho=X$ or $r_1=q^dX/R$). This proves the claim.

\answer{(proof)} For any non-liquidating structure, $V^d\le q^dX/R$ gives $\widehat C\ge(LR-q^dX)/f^d$;
(5.12) is $C\le(LR-q^dX)/f^d$; hence $C\le\widehat C$ and $U^L-U^N=e(\widehat C-C)\ge0$: the good entrepreneur
prefers the liquidating structure.

\step{7}{Remark on $\widehat C>0$} \eqref{t5e11-Cb} also shows $\widehat C\ge0$ whenever $LR\ge q^dX$, i.e.\ whenever lenders really prefer liquidating
a downgraded firm to rolling it over --- the condition for a liquidating structure to exist. Without it, $\widehat C$ can be negative:
with the numbers of Amaro de Matos exercise 5.2.5 \#5 ($L=R=1$, $X=\tfrac74$, $f^d=\tfrac13$, $q^d=\tfrac23$) and short-term debt
($V^d=1$), $\widehat C=3-1.75-(1)(1)\bigl(3-\tfrac32\bigr)=-0.25$; so the lemma's ``$>0$'' presumes credible liquidation.

\intuition Liquidation costs the good entrepreneur her firm (and $C$) in the downgrade state, but it shifts financing costs onto the pool
of downgraded firms, most of whom are bad: lenders recover $L$ from every downgraded firm, so the upgraded, i.e.\ good, firms pay less.
(5.12) says the continuation value of a downgraded firm --- what lenders could ever collect, $q^dX$, plus the control rent weighted by the
share of good types, $f^dC$ --- is below $L$; then the cross-subsidy gain outweighs the lost rent.

\pitfalls
\begin{itemize}
  \item Forgetting that a liquidated good entrepreneur loses both $X$ and the control rent $C$ in the
        downgrade state: $U^L$ has $(1-e)(X+C)$, not $X+C$.
  \item Signing $\frac1{f^d}-\frac1{q^d}$ wrongly: $q^d>f^d$ by (5.10), so it is positive and $\widehat C$
        \emph{decreases} in $V^d$; the bound $V^d\le q^dX/R$ is therefore the worst case for liquidation.
  \item Dropping $R$ when moving between (5.12) and $C\le(LR-q^dX)/f^d$: (5.12) is stated per unit of $R$.
\end{itemize}

\bookans The book (p.~268) only cites the appendix of Diamond (1991), pp.~732--733; proved here.

\refbib Amaro de Matos \S5.2.3 ``The Choice of the Duration'', pp.~139--142 (Lemmas 5.2.1--5.2.2, (5.9)--(5.12)); exercise p.~148.

% =============================================================================
\exsec{12. Diamond (1991): liquidation after a downgrade --- Amaro de Matos 5.2.5 \#6 (with \#5), p.~148}
%% BEGIN-STATEMENT 12
\stmtbegin
\textbf{Source:} Amaro de Matos, Ch. 5, Exercise 5.2.5 \#6, p.~148 (PDF 166) · Answer in book: no
(p.~268 only cites examples 2 and 3 of Diamond 1991, pp.~717--719)

\begin{enumerate}
\def\labelenumi{\arabic{enumi}.}
\setcounter{enumi}{5}
\tightlist
\item
  This exercise illustrates the case of Lemma 5.2.2 in the model of Diamond (1991). Keep everything
  from the previous exercise except that here \(X = 1.4\) and \(C = 0.35\). Show that in this case,
  if the project is financed with short-term debt, it will be liquidated in the case of a downgrade
  at \(t = 1\). Show also that this leads the entrepreneur of a good project to prefer long-term to
  short-term debt. Finally, show that this conclusion is reversed if the initial credit rating is
  \(f = \frac{3}{4}\) instead of the initial \(\frac{1}{2}\).
\end{enumerate}

\paragraph{Context needed to solve it}

\textbf{Part 0.} \emph{From Amaro de Matos, Ch. 5, Exercise 5.2.5 \#5, p.~148 (PDF 166) --- ``the
previous exercise'':}

\begin{enumerate}
\def\labelenumi{\arabic{enumi}.}
\setcounter{enumi}{4}
\tightlist
\item
  This exercise illustrates the case of Lemma 5.2.1 in the model of Diamond (1991). Take
  \(L = R = 1\), \(\pi = \frac{1}{2}\), \(X = 1\frac{3}{4}\), \(C = 0\), \(f = \frac{1}{2}\), and
  \(f^d = \frac{1}{3}\).
\end{enumerate}

\begin{enumerate}
\def\labelenumi{(\alph{enumi})}
\item
  What are the face values of a short-term and a long-term debt contract?
\item
  Show that in this example there will be no liquidation.
\item
  What is the payoff of a good project financed with long-term debt?
\item
  What is the probability that a good project receives a downgrade at \(t = 1\)?
\item
  What is the payoff of a good project financed with short-term debt?
\end{enumerate}

\emph{From Amaro de Matos §5.2.3 ``The Choice of the Duration'', pp.~139--141 (PDF 157--159):}

A model developed by Diamond (1991, 1993) illustrates in a very simple way the differences between
short-term and long-term debt financing. The model is based on information asymmetry and credit
rating.

Regaining the notation of Aghion and Bolton, consider an entrepreneur with a two-period project who
needs funds for an initial investment \(K = 1\). The project may be of one of two types: good or
bad. At time \(t = 2\) a good project has a payoff \(X > 0\), whereas a bad project pays that amount
with a probability \(\pi\), and pays zero with probability \(1 - \pi\). Let \(\theta\) denote the
type of project, \(\theta_g\) referring to a good project and \(\theta_b\) referring to a bad
project. The quality of the project is known to the entrepreneur but not to the investor. Given a
project, as of date \(t = 0\) the investor assigns a probability \(f\) that \(\theta = \theta_g\).
This probability is referred to as a \emph{credit rating} and reflects the lender's information
about the project. It follows that, as of time 0, the probability \(q\) that a borrower pays the
loan at date \(t = 2\) is given by

\[q \equiv f + \pi(1 - f) = \pi + f(1 - \pi). \qquad (5.9)\]

At time \(t = 1\) information about the project becomes available to the investor. This new
information is not verifiable, so it makes no sense to write contracts contingent upon that
information, since they cannot be enforced. However, the investors update their beliefs about the
quality of the project and the credit rating will be either upgraded or downgraded.

Under this revision of the credit rating, let \(f^u\) denote the probability that
\(\theta = \theta_g\) given an upgrade and let \(f^d\) denote the probability that
\(\theta = \theta_g\) given a downgrade. In that case, the probability that an entrepreneur who
receives a downgrade repays at time \(t = 2\) is

\[q^d = \pi + f^d(1 - \pi) > f^d. \qquad (5.10)\]

Moreover, it is here assumed that

\[f^d < f < f^u,\]

meaning that the probability that the project is good is greater given an upgrade than given a
downgrade. Furthermore, the simplifying assumption \(f^u = 1\) is also made. Thus, given an upgrade
at date \(t = 1\), it is certain that the project is good. However, it is not certain that all
projects of good quality are upgraded. There is room for one such project to be downgraded. Let
\(e\) denote the probability of such an event.

All projects are also assumed to produce a nonassignable control rent \(C\) if the management has
control at time \(t = 2\). Also, if at time \(t = 1\) the manager defaults on short-term debt,
lenders then have the right to remove the management from control, selling the assets for a
liquidation value \(L < K = 1\). All projects may be liquidated at that point in time, under those
conditions. If the project is liquidated, the entrepreneur is prevented from consumption of his/her
present or future cash flows and control rents. In this case, an optimal contract is a debt contract
enforced by the right to liquidate the project if the debt is not repaid in full.

Just as in the case of the borrowers, lenders are assumed be risk neutral. Moreover, lenders have
alternative investments at \(t = 0\) that provide a return \(R\) at \(t = 2\). The condition
\(L < X/R\) ensures that a good project yields a higher return when not liquidated.

A project may be financed with either long-term or short-term debt. In the context of this model,
long-term debt refers to a contract issued at \(t = 0\) with no refinancing at date \(t = 1\) and
maturing at \(t = 2\). If one such instrument provides \(K = 1\) at time \(t = 0\) with face value
\(\rho \le X\), the lender expects to receive at maturity \(\rho[\pi + f(1 - \pi)] = R^2\). Notice
that if \(\rho > X\), a borrower with credit rating \(f\) will not be able to provide the lender
with the required rate of return \(R^2\). Now, if the project is upgraded at time \(t = 1\), the
market value of the debt contract at that point in time becomes \(V_l^u = \rho/R\). Similarly, if
the project is downgraded, the market value of debt at time \(t = 1\) becomes \(V_l^d = q^d\rho/R\).

Short-term debt, on the other hand, is debt financed at \(t = 0\) with face value \(r_1\) and
maturity at time \(t = 1\). If short-term borrowers are upgraded at \(t = 1\), their debt contract
at that point in time is worth \(V_s^u = r_1\) and they may refinance their project by issuing a new
short-term debt contract with maturity at \(t = 2\) and face value \(r_1 R\) that will be paid with
certainty and guarantees a return of \(R\) per unit invested by the lender. If short-term borrowers
are downgraded, however, in principle they may refinance their debt with a new short-term contract
only if \(q^d X/R > r_1\), since \(X\) is the maximum that a borrower can pay at \(t = 2\). If that
is the case, at \(t = 1\) the market value of the short-term debt issued at \(t = 0\) would be
\(V_s^d = r_1\) if the borrower pays, possibly through the refinancement mechanism. But what happens
if \(q^d X/R < r_1\)? In that case there is still room for a refinancement provided that the
liquidation value \(L\) is less than the expected value to the lender from a refinancement,
\(q^d X/R\). Under a downgrading, the time \(t = 1\) value of the short-term debt is thus

\[V_s^d = \min\left[r_1, \max(L, q^d X/R)\right].\]

The face value of short-term debt is established such as to provide the lenders with an expected
return of \(R\) per period. It then follows that

\[r_1 = \frac{R - [1 - (1 - e)f]V_s^d}{(1 - e)f}. \qquad (5.11)\]

\emph{From Amaro de Matos §5.2.3, pp.~141--142 (PDF 159--160):}

Given this model, the point is to characterize under what conditions short-term or long-term
maturities will be offered. The maturity preferred by the entrepreneurs of good projects will be the
only one offered to all borrowers. This happens because if any other maturity is offered, it reveals
a project of bad quality and lenders will not be willing to finance it.

The following result implies that, in the absence of liquidation given a downgrade, the
entrepreneurs of good projects prefer short-term debt:

\textbf{Lemma 5.2.1 (Diamond, 1991).} \emph{Comparing two debt contracts issued at \(t = 0\),
neither of which involves liquidation on a downgrade at \(t = 1\), the one with the higher value at
that point in time given a downgrade, \(V^d\), is preferred by the entrepreneurs of good projects.}

\emph{Proof.} The expected payoff to the entrepreneurs of good projects from debt not yielding
liquidation is

\[e\left(X + C - \frac{RV^d}{q^d}\right) + (1 - e)(X + C - r_1 R).\]

Using expression (5.11) this payoff is rewritten as

\[X + C - \frac{R^2}{f} + RV^d e\left(\frac{1}{f^d} - \frac{1}{q^d}\right).\]

Given (5.10) this expression for the payoff is clearly an increasing function of \(V^d\). ■

\emph{From Amaro de Matos §5.2.3, p.~142 (PDF 160):}

This implies that when there is no liquidation after a downgrade, short-term debt is preferred. The
idea behind this result is simple. If there is no risk of liquidation, in order to reduce their
expected financing costs, borrowers of the good type, who are less likely than average to be
downgraded, prefer to extract larger payments from those who are downgraded. However, if there is
the possibility of liquidation, the increase in the payment extracted from those who get a downgrade
incorporates, for the good entrepreneurs, the risk of losing control rents. In that case the
following result characterizes the trade-off incurred:

\textbf{Lemma 5.2.2 (Diamond, 1991).} \emph{Consider a debt structure that does not imply
liquidation and leads to a value of debt \(V^d\), given a downgrade. Compare it with an alternative
structure that implies liquidation. Entrepreneurs of good projects prefer the latter if and only if
the control rent \(C < \widehat C\), where the upper bound is given by}

\[\widehat C = \frac{LR}{f^d} - X - RV^d\left(\frac{1}{f^d} - \frac{1}{q^d}\right) > 0.\]

It follows that a sufficient condition for the good entrepreneur to prefer liquidation is that
\(C = 0\), representing the absence of control rents. It is also sufficient for the good
entrepreneur to prefer liquidation if that state does not provide sufficiently high control rents
\(C\) to the entrepreneur compared to the proceeds \(L\) to the investor or

\[\frac{q^d X + f^d C}{R} \le L. \qquad (5.12)\]

\emph{Omitted: Proposition 5.2.1 (p.~142, PDF 160), which gives the credit-rating bound above which
good entrepreneurs prefer short-term debt; applied to these numbers it states the exercise's
conclusion.}
\stmtend
%% END-STATEMENT

\asked \emph{Part 0} (exercise \#5, needed as the base case): (a) face values of short- and long-term
debt; (b) show there is no liquidation; (c) the good entrepreneur's payoff with long-term debt; (d) the
probability $e$ that a good project is downgraded; (e) her payoff with short-term debt.
\emph{(i)} With $X=1.4$, $C=0.35$, show that short-term debt now leads to liquidation after a downgrade.
\emph{(ii)} Show that the good entrepreneur then prefers long-term debt. \emph{(iii)} Show that this
preference reverses when the initial rating is $f=\tfrac34$.

\approach Diamond's (1991) maturity trade-off. Short-term debt is repriced at $t=1$ on nonverifiable
rating news: without liquidation that only shifts repayment onto the downgraded pool, which helps good
types (Lemma~5.2.1); with liquidation it also costs the good type $X+C$ whenever she is downgraded
(Lemma~5.2.2: liquidating debt wins iff $C<\widehat C$). A downgraded borrower is liquidated when the
most she can promise, $q^dX/R$, is below both $r_1$ and $L$. The rating $f$ matters through $e$: a higher
$f$ (with $f^d$ fixed) means good types are rarely downgraded, so the liquidation cost is rarely paid.

\notation
\begin{itemize}
  \item $X$: $t=2$ payoff (good project surely, bad with probability $\pi$); $C\ge0$: control rent; $L$: liquidation value; $R$: gross return per period.
  \item $f$: initial rating; $f^u=1$, $f^d$: posterior of ``good'' after an upgrade / downgrade; $e$: probability that a good project is downgraded;
        $P_d$: probability of a downgrade for an unknown type; $q^d=\pi+f^d(1-\pi)$; $V^d$: $t=1$ value of the debt after a downgrade ($V_s^d$ short, $V_l^d$ long).
  \item $\rho$: face value of long-term debt; $r_1$: face value of short-term debt; $q$: $t=0$ repayment probability (5.9);
        $F^d$: face value of the $t=1$ refinancing after a downgrade; $U_{LT}$, $U_{ST}$: good entrepreneur's payoff with long- / short-term debt.
  \item Part 0 (\#5): $L=R=1$, $\pi=\frac12$, $X=\frac74$, $C=0$, $f=\frac12$, $f^d=\frac13$. Exercise \#6: $X=1.4$, $C=0.35$, then $f=\frac34$.
\end{itemize}
\assume (a) $L=1=K$ contradicts the text's $L<K$; it is harmless here but makes short-term debt riskless ($r_1=R$).
(b) When $f$ becomes $\frac34$ we keep $f^d=\frac13$, as the statement says (``keep everything'') and as Proposition 5.2.1 does; then $e$ adjusts.
The book only cites Diamond (1991), examples 1--3; all numbers below are derived here.

\step{1}{Part 0 (a): face values} Each face value is set by the lenders' break-even at its horizon. By (5.9),
\begin{align}
  q &= \pi+f(1-\pi) = \tfrac12+\big(\tfrac12\big)\big(\tfrac12\big) = \tfrac34. \label{t5e12-q}
\end{align}
Long-term debt satisfies $\rho q=R^2$:
\begin{align}
  \rho &= \frac{R^2}{q} = \frac{(1)}{(3/4)} = \frac43 \;\le X. \label{t5e12-rho}
\end{align}
For short-term debt we need $e$, $P_d$, $q^d$. Bad projects are always downgraded ($f^u=1$) and good ones
with probability $e$, so $P_d=fe+(1-f)$ and Bayes' rule gives
\begin{align}
  f^d &= \frac{fe}{fe+(1-f)}. \label{t5e12-bayes}
\end{align}
Solving for $e$: $f^d[fe+(1-f)]=fe$, so $f^d(1-f)=fe(1-f^d)$ and
\begin{align}
  e &= \frac{f^d(1-f)}{f(1-f^d)} = \frac{(1/3)(1/2)}{(1/2)(2/3)} = \frac{1/6}{1/3} = \frac12, \label{t5e12-e}\\
  P_d &= 1-(1-e)f = 1-\big(\tfrac12\big)\big(\tfrac12\big) = \tfrac34, \label{t5e12-Pd}\\
  q^d &= \pi+f^d(1-\pi) = \tfrac12+\big(\tfrac13\big)\big(\tfrac12\big)=\tfrac23. \label{t5e12-qd}
\end{align}
A downgraded firm can roll over up to $q^dX/R=(2/3)(7/4)=7/6>L$. Conjecture no liquidation, $V_s^d=r_1$; then (5.11),
multiplied by $(1-e)f$, gives
\begin{align}
  (1-e)f\,r_1 &= R-P_d\,r_1, \nonumber\\
  r_1\big[(1-e)f+P_d\big] &= R, \nonumber\\
  r_1 &= R = 1, \label{t5e12-r1}
\end{align}
since $(1-e)f+P_d=1$. Indeed $r_1=1\le 7/6$, so $V_s^d=\min[1,\max(1,7/6)]=1=r_1$: consistent.

\answer{0(a)} $r_1=1$, $\rho=\tfrac43$.

\step{2}{Part 0 (b): no liquidation} Long-term debt has no $t=1$ payment. A downgraded short-term borrower refinances $r_1=1$ with face
\begin{align}
  F^d &= \frac{Rr_1}{q^d} = \frac{(1)(1)}{(2/3)} = \frac32 \;\le X=\frac74, \label{t5e12-Fd}
\end{align}
and the lenders would not liquidate anyway since $q^dX/R=7/6>L=1$.

\answer{0(b)} No liquidation: the downgraded borrower refinances at face $F^d=\tfrac32\le X$.

\step{3}{Part 0 (c)--(e)} (c) The good project always pays $\rho$:
$U_{LT}=X+C-\rho=\frac74-\frac43=\frac{21-16}{12}=\frac5{12}=0.416667$. (d) By \eqref{t5e12-e}, $e=\frac12$.
(e) Upgraded firms roll over at $r_1R=1$, downgraded ones pay $F^d$:
\begin{align}
  U_{ST} &= e(X+C-F^d)+(1-e)(X+C-r_1R) \nonumber\\
         &= \big(\tfrac12\big)\big(\tfrac74-\tfrac32\big)+\big(\tfrac12\big)\big(\tfrac74-1\big) \nonumber\\
         &= \big(\tfrac12\big)\big(\tfrac14\big)+\big(\tfrac12\big)\big(\tfrac34\big) = \tfrac18+\tfrac38 = \tfrac12. \label{t5e12-UST0}
\end{align}
Short-term debt is preferred ($\frac12>\frac5{12}$), as Lemma 5.2.1 predicts: by the Lemma's payoff expression,
with $R^2/f=2$ and $e(\frac1{f^d}-\frac1{q^d})=\frac12\bigl(3-\frac32\bigr)=\frac34$,
\begin{align}
  U &= X+C-\frac{R^2}{f}+eRV^d\Big(\frac1{f^d}-\frac1{q^d}\Big) = \tfrac74+0-2+\tfrac34V^d = -\tfrac14+\tfrac34V^d, \label{t5e12-lemma1}
\end{align}
with $V_s^d=1$ (giving $-\frac14+\frac34=\frac12$) $>V_l^d=q^d\rho/R=\frac89$ (giving $-\frac14+\frac23=\frac5{12}$).

\answer{0(c)--(e)} (c) $U_{LT}=\tfrac5{12}\approx0.416667$; (d) $e=\tfrac12$; (e) $U_{ST}=\tfrac12$.

\step{4}{\#6, $f=\frac12$: short-term debt is liquidated after a downgrade} Now $X=1.4$, $C=0.35$; $q,\rho,e,P_d,q^d$ are unchanged, and
$\rho=\frac43\le1.4$ keeps long-term debt feasible. A downgraded firm can now roll over at most
\begin{align}
  \frac{q^dX}{R} &= \big(\tfrac23\big)(1.4) = \tfrac{14}{15} = 0.933333 \;<\; L=1. \label{t5e12-roll}
\end{align}
If it were not liquidated, \eqref{t5e12-r1} would give $r_1=1>14/15$: it could not refinance. So it defaults, and the lenders, offered at
most $14/15$ by continuation against $L=1$ from liquidation, liquidate. Hence $V_s^d=\max(L,q^dX/R)=1$ and by (5.11)
\begin{align}
  r_1 &= \frac{R-P_dV_s^d}{(1-e)f} = \frac{1-(3/4)(1)}{(1/4)} = 1, \label{t5e12-r1b}
\end{align}
consistent with $V_s^d=\min[1,\max(1,14/15)]=1$.

\answer{(i)} $q^dX/R=\tfrac{14}{15}<r_1=1$ and $<L=1$: a downgraded short-term borrower cannot refinance
and is liquidated.

\step{5}{\#6, $f=\frac12$: long-term debt is preferred} The good firm keeps control only after an upgrade:
\begin{align}
  U_{ST} &= (1-e)(X+C-r_1R) = \big(\tfrac12\big)(1.4+0.35-1) = \big(\tfrac12\big)(0.75) = 0.375, \label{t5e12-UST1}\\
  U_{LT} &= X+C-\rho = 1.75-\tfrac43 = 0.416667. \label{t5e12-ULT1}
\end{align}
$U_{LT}>U_{ST}$. To check this against Lemma~5.2.2, write both payoffs with (5.11) and Bayes
($P_d/f=e/f^d$, from \eqref{t5e12-bayes}). The liquidating short-term debt has $V_s^d=L$, so
$(1-e)r_1R=R(R-P_dL)/f=R^2/f-eRL/f^d$ and
$U_{ST}=(1-e)(X+C)-R^2/f+eRL/f^d$; the long-term debt has the Lemma~5.2.1 payoff
$U_{LT}=X+C-R^2/f+eRV_l^d(1/f^d-1/q^d)$. Subtracting,
\begin{align}
  U_{ST}-U_{LT} &= -e(X+C)+\frac{eRL}{f^d}-eRV_l^d\Big(\frac1{f^d}-\frac1{q^d}\Big)
  = e\big(\widehat C-C\big), \label{t5e12-diff}
\end{align}
with $\widehat C$ as in Lemma~5.2.2. With $V^d=V_l^d=\frac89$:
$\widehat C=\frac{(1)(1)}{1/3}-1.4-(1)\big(\tfrac89\big)\big(3-\tfrac32\big)=1.6-\tfrac43=0.266667<C=0.35$, and \eqref{t5e12-diff} gives
$U_{ST}-U_{LT}=(0.5)(0.266667-0.35)=-0.041667$. (Condition (5.12) fails: $(q^dX+f^dC)/R=0.933333+0.116667=1.05>L$.)

\answer{(ii)} $U_{LT}=0.416667>U_{ST}=0.375$ (gap $-0.041667=e(\widehat C-C)$ with $\widehat C=0.266667<C$):
long-term debt is preferred.

\step{6}{\#6, $f=\frac34$: the conclusion reverses} With $f^d=\frac13$ kept, \eqref{t5e12-e}--\eqref{t5e12-Pd} and \eqref{t5e12-q}--\eqref{t5e12-rho} give
\begin{align}
  e &= \frac{(1/3)(1/4)}{(3/4)(2/3)} = \frac{1/12}{1/2} = \frac16, \quad P_d = 1-\big(\tfrac56\big)\big(\tfrac34\big)=1-\tfrac58=\tfrac38, \nonumber\\
  q &= \tfrac12+\big(\tfrac34\big)\big(\tfrac12\big)=\tfrac78, \quad \rho=\frac{(1)}{(7/8)}=\tfrac87=1.142857. \nonumber
\end{align}
$q^d=\frac23$ is unchanged, so \eqref{t5e12-roll} still holds and short-term debt is still liquidated after a downgrade;
with $(1-e)f=\big(\tfrac56\big)\big(\tfrac34\big)=\tfrac58$, (5.11) gives $r_1=\frac{1-(3/8)(1)}{(5/8)}=\frac{5/8}{5/8}=1$. Then
\begin{align}
  U_{ST} &= (1-e)(X+C-r_1R) = \big(\tfrac56\big)(0.75) = 0.625, \label{t5e12-UST2}\\
  U_{LT} &= X+C-\rho = 1.75-\tfrac87 = 0.607143, \label{t5e12-ULT2}
\end{align}
so now short-term (liquidating) debt is preferred. Check: $V_l^d=q^d\rho/R=\big(\tfrac23\big)\big(\tfrac87\big)=\frac{16}{21}$,
$\widehat C=1.6-\big(\tfrac{16}{21}\big)\big(\tfrac32\big)=1.6-\tfrac87=0.457143>C=0.35$, and by \eqref{t5e12-diff}
$e(\widehat C-C)=\frac16(0.107143)=0.017857=U_{ST}-U_{LT}$.
Proposition 5.2.1 (p.~142) agrees: its bound is
$\frac{\pi[R^2-LR+f^d(C+X-R^2)]}{(1-\pi)[LR-f^d(C+X)]}=\frac{(1/2)[1-1+(1/3)(0.75)]}{(1/2)[1-(1/3)(1.75)]}=\frac{(1/2)(0.25)}{(1/2)(0.416667)}=0.6$, and $f=\frac12<0.6<\frac34$.

\answer{(iii)} At $f=\tfrac34$ ($e=\tfrac16$): $U_{ST}=0.625>U_{LT}=0.607143$ (gap $0.017857$): short-term,
liquidating debt is now preferred.

\intuition With a low rating the good firm is downgraded often ($e=\frac12$); liquidation then costs it the control rent too often, and
long-term debt, which never triggers liquidation, wins. With a high rating, good firms are rarely downgraded ($e=\frac16$), while the
downgraded pool is mostly bad firms; liquidating that pool lets good firms borrow cheaply after the upgrade, which now outweighs the rare loss of control.
(If instead $e=\frac12$ were kept at $f=\frac34$, then $f^d=\frac35$, $q^dX/R=1.12>L$, there is no liquidation, and short-term debt still wins by Lemma 5.2.1, $0.625$ vs $0.607143$.)

\pitfalls
\begin{itemize}
  \item Keeping $e=\tfrac12$ when $f$ rises to $\tfrac34$: the statement keeps everything else, including
        $f^d=\tfrac13$, so $e$ must be recomputed by Bayes ($e=\tfrac16$).
  \item Giving the liquidated good entrepreneur something after a downgrade: liquidation removes $X$ and
        $C$, so $U_{ST}=(1-e)(X+C-r_1R)$ only.
  \item Checking only $q^dX/R<r_1$ for liquidation: lenders liquidate because, in addition, $L$ beats the
        best refinancing offer, $L=1>q^dX/R=\tfrac{14}{15}$.
\end{itemize}

\bookans The book (p.~268) only cites examples 1--3 of Diamond (1991), pp.~716--719; derived here. The three claims of \#6 are confirmed
(liquidation after downgrade; long-term preferred at $f=\frac12$; reversal at $f=\frac34$), consistent with Proposition 5.2.1.

\refbib Amaro de Matos \S5.2.3 ``The Choice of the Duration'', pp.~139--142 ((5.9)--(5.12), Lemmas 5.2.1--5.2.2, Proposition 5.2.1); exercises p.~148.

% =============================================================================
\exsec{13. A borrower-friendly bankruptcy court --- Tirole Exercise 4.9, pp.~190--191}
%% BEGIN-STATEMENT 13
\stmtbegin
\textbf{Source:} Tirole, Ch. 4, Exercise 4.9, pp.~190--191 (PDF 203) · Answer in book: yes, p.~586

\textbf{Exercise 4.9 (borrower-friendly bankruptcy court).} Consider the timing described in Figure
4.7.

Figure 4.7 (timeline, left to right; p.~191, PDF 204)

{\def\LTcaptype{none} % do not increment counter
\begin{longtable}[]{@{}
  >{\raggedright\arraybackslash}p{(\linewidth - 6\tabcolsep) * \real{0.2500}}
  >{\raggedright\arraybackslash}p{(\linewidth - 6\tabcolsep) * \real{0.2500}}
  >{\raggedright\arraybackslash}p{(\linewidth - 6\tabcolsep) * \real{0.2500}}
  >{\raggedright\arraybackslash}p{(\linewidth - 6\tabcolsep) * \real{0.2500}}@{}}
\toprule\noalign{}
\begin{minipage}[b]{\linewidth}\raggedright
Date 0
\end{minipage} & \begin{minipage}[b]{\linewidth}\raggedright
Date 1
\end{minipage} & \begin{minipage}[b]{\linewidth}\raggedright
\end{minipage} & \begin{minipage}[b]{\linewidth}\raggedright
Date 2
\end{minipage} \\
\midrule\noalign{}
\endhead
\bottomrule\noalign{}
\endlastfoot
Investment, contract. Entrepreneur has cash \(A\), must borrow \(I - A\). & Short-term profit
realized. Liquidate → Liquidation value \(L\) (known); or Continue → & Moral hazard: Effort decision
(high or low). & Final profit: Success (profit \(R\)) or failure (profit 0)? \\
\end{longtable}
}

The project, if financed, yields random and verifiable short-term profit \(r \in [0, \bar r]\) (with
a continuous density and \emph{ex ante} mean \(E[r]\)). After \(r\) is realized and cashed in, the
firm either liquidates (sells its assets), yielding some known liquidation value \(L > 0\), or
continues. Note that (the random) \(r\) and (the deterministic) \(L\) are not subject to moral
hazard. If the firm continues, its prospects improve with \(r\) (so \(r\) is ``good news'' about the
future). Namely, the probability of success is \(p_H(r)\) if the entrepreneur works between dates 1
and 2 and \(p_L(r)\) if the entrepreneur shirks. Assume that \(p_H' > 0\), \(p_L' > 0\), and

\[p_H(r) - p_L(r) \equiv \Delta p\]

is independent of \(r\) (so shirking reduces the probability of success by a fixed amount
independent of prospects). As usual, one will want to induce the entrepreneur to work if
continuation obtains. It is convenient to use the notation

\[\rho_1(r) \equiv p_H(r)R \quad \text{and} \quad \rho_0(r) \equiv p_H(r)\left[R - \frac{B}{\Delta p}\right].\]

Investors are competitive and demand an expected rate of return equal to 0. Assume

\[\rho_1(r) > L \quad \text{for all } r \qquad (1)\]

and

\[E[r] + L > I - A > E[r + \rho_0(r)]. \qquad (2)\]

\begin{enumerate}
\def\labelenumi{(\roman{enumi})}
\tightlist
\item
  Argue informally that, in the optimal contract for the borrower, the short-term profit and the
  liquidation value (if the firm is liquidated) ought to be given to investors.
\end{enumerate}

Argue that, in the case of continuation, \(R_b = B/\Delta p\). (If you are unable to show why, take
this fact for granted in the rest of the question.)

Interpret conditions (1) and (2).

\begin{enumerate}
\def\labelenumi{(\roman{enumi})}
\setcounter{enumi}{1}
\tightlist
\item
  Write the borrower's optimization program.
\end{enumerate}

Assume (without loss of generality) that the firm continues if and only if \(r \geqslant r^*\) for
some \(r^* \in (0, r)\). Exhibit the equation defining \(r^*\).

\begin{enumerate}
\def\labelenumi{(\roman{enumi})}
\setcounter{enumi}{2}
\tightlist
\item
  Argue that this optimal contract can be implemented using, inter alia, a short-term debt contract
  at level \(d = r^*\). Interpret ``liquidation'' as a ``bankruptcy.''
\end{enumerate}

How does short-term debt vary with the borrower's initial equity? Explain.

\begin{enumerate}
\def\labelenumi{(\roman{enumi})}
\setcounter{enumi}{3}
\tightlist
\item
  Suppose that, when the decision to liquidate is taken, the firm must go to a bankruptcy court. The
  judge mechanically splits the bankruptcy proceeds \(L\) equally between investors and the
  borrower.
\end{enumerate}

Define \(\hat r\) by

\[\rho_0(\hat r) \equiv \tfrac{1}{2}L.\]

Assume first that

\[r^* > \hat r\]

(where \(r^*\) is the value found in question (ii)).

Show that the borrower-friendly court actually prevents the borrower from having access to
financing. (Note: a diagram may help.)

\begin{enumerate}
\def\labelenumi{(\alph{enumi})}
\setcounter{enumi}{21}
\tightlist
\item
  Continuing on question (iv), show that when
\end{enumerate}

\[r^* < \hat r,\]

the borrower-friendly court either prevents financing or increases the probability of bankruptcy,
and in all cases hurts the borrower and not the lenders.

\paragraph{Context needed to solve it}

The statement defines its own model; ``as usual'' refers to the fixed-investment model's
moral-hazard setup (private benefit \(B\), limited liability, incentive constraint), copied here.

\emph{From Tirole §3.2.1 ``The Fixed-Investment Model'', pp.~115--116 (PDF 128--129):}

Variants of the following entrepreneurial model will be used in the following: an entrepreneur (also
called the ``insider'' or the ``borrower,'' ``she'') has a project. This project requires a fixed
investment \(I\). The entrepreneur initially has ``assets'' (``cash on hand'' or ``net worth'')
\(A < I\). For the moment we interpret these assets as being cash or liquid securities that can be
used toward covering the cost of investment. {[}\ldots{]} The entrepreneur's cash can either be
invested in the project or used for consumption. To implement the project the entrepreneur must
borrow \(I - A\) from lenders. {[}\ldots{]}

\emph{Project.} If undertaken, the project either succeeds, that is, yields verifiable income
\(R > 0\), or fails, that is, yields no income. The probability of success is denoted by \(p\). The
project is subject to moral hazard. The entrepreneur can ``behave'' (``work,'' ``exert effort,''
``take no private benefit'') or ``misbehave'' (``shirk,'' ``take a private benefit'') {[}\ldots{]}.

Behaving yields probability \(p = p_H\) of success and no private benefit to the entrepreneur, and
misbehaving results in probability \(p = p_L < p_H\) of success and private benefit \(B > 0\)
(measured in units of account) to the entrepreneur. In the ``effort interpretation,'' \(B\) can also
be interpreted as a disutility of effort saved by the entrepreneur when shirking. Let
\(\Delta p \equiv p_H - p_L\).

\emph{Preferences and the loan agreement.} Both the borrower and the potential lenders (or
``investors'') are risk neutral. For notational simplicity, there is no time preference; the rate of
return expected by investors (which is also the riskless rate, due to risk neutrality) is taken to
be 0. The borrower is protected by limited liability, and so her income cannot take negative values.

Lenders behave competitively in the sense that the loan, if any, makes zero profit. {[}\ldots{]}

Let us turn to the loan contract. A contract first stipulates whether the project is financed. If
so, it further specifies how the profit is shared between the lenders and the borrower. The
borrower's limited liability will imply that both sides receive 0 in the case of failure
{[}\ldots{]}. In the case of success, the two parties share the profit \(R\); \(R_b\) goes to the
borrower and \(R_l\) to the lenders. To sum up, we posit an incentive scheme for the entrepreneur of
the following form: \(R_b\) in the case of success, 0 in the case of failure.

\emph{From Tirole §3.2.2 ``The Lenders' Credit Analysis'', pp.~116--117 (PDF 129--130):}

{[}\ldots{]} the borrower will therefore behave if the following ``incentive compatibility
constraint'' is satisfied:

\[p_H R_b \geqslant p_L R_b + B \quad \text{or} \quad (\Delta p) R_b \geqslant B. \qquad (\text{IC}_b)\]

From this incentive compatibility constraint we infer that the highest income in the case of success
that can be pledged to the lenders without jeopardizing the borrower's incentives is

\[R - \frac{B}{\Delta p}.\]

The \emph{(expected) pledgeable income} is then

\[\mathcal{P} = p_H\left(R - \frac{B}{\Delta p}\right).\]
\stmtend
%% END-STATEMENT

\asked (i) Explain why the date-1 profit $r$ and the liquidation value $L$ should go to investors, why the
borrower's stake on continuation is exactly $B/\Delta p$, and what conditions (1) and (2) mean. (ii) Write
the borrower's optimisation over the continuation rule and give the equation for the cutoff $r^*$.
(iii) Show that short-term debt $d=r^*$ implements the optimum (liquidation = bankruptcy) and say how
$d$ moves with the borrower's equity $A$. (iv) If a court gives the borrower half of $L$ in bankruptcy and
$r^*>\hat r$, show that the borrower can no longer be financed. (v) If $r^*<\hat r$, show that the court
either blocks financing or makes bankruptcy more likely, hurting the borrower and not the lenders.

\approach Tirole's fixed-investment model with an interim, verifiable signal $r$ and a liquidation
option. Continuation is always efficient but pledges only $\rho_0(r)<\rho_1(r)$, while liquidation
pledges all of $L$; when pledgeable income is scarce, the optimal contract buys it with a \emph{threat of
termination} after bad news, which short-term debt implements (default $\Rightarrow$ creditors take control
and liquidate). Anything that lowers what investors get in bankruptcy lowers pledgeable income ex ante;
competitive lenders still break even, so the cost falls on the borrower.

\notation
\begin{itemize}
  \item $I$, $A<I$: investment and cash; $r\in[0,\bar r]$: verifiable date-1 profit, density $f$, c.d.f.\ $F$, mean $E[r]$.
  \item $L>0$: liquidation value; $R$: date-2 success profit; $p_H(r)$, $p_L(r)$: success probabilities (work / shirk), $p_H'>0$, $\Delta p$ constant; $B$: private benefit.
  \item $\rho_1(r)=p_H(r)R$, $\rho_0(r)=p_H(r)[R-B/\Delta p]$: expected and pledgeable continuation income, both increasing in $r$.
  \item $r^*$: continuation cutoff; $d$: short-term debt; $\hat r$: $\rho_0(\hat r)=L/2$; $r^{**}$: cutoff under the borrower-friendly court.
  \item $\mathcal P(x)$: investors' expected income with cutoff $x$; $\text{NPV}(x)$: total surplus with cutoff $x$.
\end{itemize}

\step{1}{(i) Who gets $r$ and $L$; $R_b=B/\Delta p$; conditions (1)--(2)} $r$ and $L$ are verifiable and not subject to moral hazard, so
pledging them to investors costs no incentives. Any part of them left to the borrower lowers pledgeable income, which (Step 2) must then be
recovered by liquidating more often, an efficiency loss by (1). For the same reason, in continuation the incentive constraint
$[p_H(r)-p_L(r)]R_b\ge B$, i.e.\ $(\Delta p)R_b\ge B$, should bind: $R_b=B/\Delta p$. Allowing $R_b(r)\ge B/\Delta p$ in the program of Step 2 leaves the NPV
unchanged (a transfer) but lowers the left side of the breakeven constraint by $\int_{r^*}^{\bar r}p_H(r)[R_b(r)-B/\Delta p]f\,dr$, forcing a higher $r^*$.

Interpretation. (1) $\rho_1(r)>L$: continuation is always efficient. (2), right inequality $I-A>E[r+\rho_0(r)]$: always continuing
leaves too little pledgeable income to finance the project. (2), left inequality $E[r]+L>I-A$: always liquidating would raise enough.
So financing requires some inefficient liquidation.

\answer{(i)} Give $r$ and $L$ to investors (they cost no incentives, and anything left to the borrower
must be recovered by liquidating more often); $R_b=B/\Delta p$ on continuation. (1): continuation is always
efficient; (2): the project cannot be financed without some liquidation, but can be with enough of it.

\step{2}{(ii) Program and cutoff} With cutoff $x$ (continue iff $r\ge x$), the borrower maximises the NPV (lenders break even, so
she receives the whole surplus): date-1 profit, plus $\rho_1(r)$ on continuation states, plus $L$ on liquidation states (probability $F(x)$):
\begin{align}
  \max_{x}\;\text{NPV}(x) &= E[r]+\int_{x}^{\bar r}\rho_1(r)f(r)\,dr+F(x)L-I \label{t5e13-npv}\\
  \text{s.t.}\quad \mathcal P(x) &\equiv E[r]+\int_{x}^{\bar r}\rho_0(r)f(r)\,dr+F(x)L \;\ge\; I-A. \label{t5e13-be}
\end{align}
(A cutoff rule is without loss: continuing at $r$ gains $\rho_1(r)-L>0$ in value and costs $L-\rho_0(r)$ in pledgeable income; the gain/cost
ratio rises with $r$, so continuation goes to the highest $r$'s.) Differentiating by Leibniz's rule,
$\frac{d}{dx}\int_x^{\bar r}g(r)f(r)\,dr=-g(x)f(x)$ and $\frac{d}{dx}F(x)L=f(x)L$, so
\begin{align}
  \text{NPV}'(x) &= -\rho_1(x)f(x)+f(x)L \nonumber\\
                 &= -f(x)\big[\rho_1(x)-L\big] < 0, \label{t5e13-dnpv}\\
  \mathcal P'(x) &= f(x)\big[L-\rho_0(x)\big]. \label{t5e13-dP}
\end{align}
By \eqref{t5e13-dnpv} the borrower wants the lowest $x$ satisfying \eqref{t5e13-be}, so $r^*$ is the lowest root of
\begin{align}
  E[r]+\int_{r^*}^{\bar r}\rho_0(r)f(r)\,dr+F(r^*)L &= I-A. \label{t5e13-rstar}
\end{align}
By (2), $\mathcal P(0)=E[r+\rho_0(r)]<I-A<E[r]+L=\mathcal P(\bar r)$, so $r^*\in(0,\bar r)$; since $\rho_0$ is increasing, $\mathcal P$ rises while $\rho_0<L$ and falls after,
so at the lowest root $\rho_0(r^*)\le L$. (The statement's ``$r^*\in(0,r)$'' should read $(0,\bar r)$.)

\answer{(ii)} $\max_x\,E[r]+\int_x^{\bar r}\rho_1f\,dr+F(x)L-I$ s.t.\ $E[r]+\int_x^{\bar r}\rho_0f\,dr+F(x)L\ge I-A$;
$r^*$ is the lowest root of $E[r]+\int_{r^*}^{\bar r}\rho_0(r)f(r)\,dr+F(r^*)L=I-A$.

\step{3}{(iii) Short-term debt and equity} Let the firm owe short-term debt $d=r^*$ at date 1 (plus, inter alia, the residual $r-r^*$ to investors
and a long-term claim $R-B/\Delta p$ on success). If $r\ge r^*$, it pays $d$ out of $r$ and continues. If $r<r^*$, it defaults; creditors take control and
liquidate (``bankruptcy''), collecting $r+L$. They do not renegotiate, because continuing would give them at most
$\rho_0(r)<\rho_0(r^*)\le L$. Default cannot be strategic since $r$ is verifiable. Differentiating \eqref{t5e13-rstar} with respect to $A$ and using \eqref{t5e13-dP},
\begin{align}
  f(r^*)\big[L-\rho_0(r^*)\big]\,\frac{dr^*}{dA} &= -1 \nonumber\\
  \frac{dr^*}{dA} &= -\frac{1}{f(r^*)\,[L-\rho_0(r^*)]} < 0. \label{t5e13-drdA}
\end{align}
Less equity requires more pledgeable income, obtained by liquidating more often: short-term debt rises.

\answer{(iii)} Short-term debt $d=r^*$: the firm continues iff it can pay $d$; otherwise creditors take
control and liquidate. $dd/dA=dr^*/dA=-1/\{f(r^*)[L-\rho_0(r^*)]\}<0$: less equity, more short-term debt.

\step{4}{(iv) Court splits $L$, case $r^*>\hat r$} Investors now get $L/2$ on liquidation. Whatever the continuation policy, their income at state $r$
(besides $r$) is at most $\max\{L/2,\rho_0(r)\}$, so the best they can be promised is
\begin{align}
  \mathcal P^{c}_{\max} &= E[r]+\int_0^{\bar r}\max\{L/2,\rho_0(r)\}f(r)\,dr. \label{t5e13-Pmax}
\end{align}
Compare state by state with the contract of (ii), which gives $L$ on $[0,r^*)$ and $\rho_0(r)$ on $[r^*,\bar r]$. For $r<\hat r$: $\max=L/2<L$. For $\hat r\le r<r^*$:
$\max=\rho_0(r)<\rho_0(r^*)\le L$. For $r\ge r^*>\hat r$: $\max=\rho_0(r)$, the same. Hence $\mathcal P^{c}_{\max}<\mathcal P(r^*)=I-A$ by \eqref{t5e13-rstar}: no contract lets
investors break even, and the borrower loses access to financing.

\answer{(iv)} With $r^*>\hat r$, even the best policy under the court yields $\mathcal P^c_{\max}<I-A$:
no financing.

\step{5}{(v) Case $r^*<\hat r$} With cutoff $x$, investors get $\mathcal P^c(x)=E[r]+\int_x^{\bar r}\rho_0 f\,dr+F(x)L/2$, with
$\mathcal P^{c\prime}(x)=f(x)[L/2-\rho_0(x)]$, positive below $\hat r$ and negative above. At the old cutoff,
\begin{align}
  \mathcal P^c(r^*) &= \mathcal P(r^*)-F(r^*)\frac L2 = I-A-F(r^*)\frac L2 < I-A. \label{t5e13-Pc}
\end{align}
So if $\mathcal P^c(\hat r)<I-A$ there is no financing; otherwise the new cutoff $r^{**}\in(r^*,\hat r]$ solves $\mathcal P^c(r^{**})=I-A$: bankruptcy becomes
more likely. Lenders break even in every case (or do not lend), so they are not hurt. The borrower receives the NPV, and by \eqref{t5e13-dnpv}
$\text{NPV}(r^{**})<\text{NPV}(r^*)$: she is hurt (or loses financing altogether). Revealed preference says the same: under an enforcing court she
could have promised investors only $L/2$ in bankruptcy, and chose not to.

\answer{(v)} Either $\mathcal P^c(\hat r)<I-A$ (no financing) or the cutoff rises to $r^{**}\in(r^*,\hat r]$
(more bankruptcy); lenders always break even, the borrower's NPV falls.

\step{6}{Numerical illustration (ours)} $r\sim U[0,1]$ ($E[r]=0.5$), $p_H(r)=0.3+0.5r$, $R=2$, $B/\Delta p=1.6$ (e.g.\ $B=0.4$, $\Delta p=0.25$), so
$\rho_1(r)=(0.3+0.5r)(2)=0.6+r$, $\rho_0(r)=(0.3+0.5r)(2-1.6)=0.12+0.2r$, and $E[r+\rho_0]=0.5+0.12+0.2(0.5)=0.72$. With $f=1$ on $[0,1]$,
\begin{align}
  \mathcal P(x) &= 0.5+\int_x^1(0.12+0.2r)\,dr+Lx = 0.5+0.12(1-x)+0.1(1-x^2)+Lx \nonumber\\
                &= 0.72+(L-0.12)x-0.1x^2. \nonumber
\end{align}
\emph{Case (v):} $L=0.5$, $I-A=0.75$ (conditions (1)--(2) hold: $0.6>0.5$, $1>0.75>0.72$). \eqref{t5e13-rstar} is
$0.72+0.38x-0.1x^2=0.75$, i.e.\ $0.1x^2-0.38x+0.03=0$, with discriminant $0.38^2-4(0.1)(0.03)=0.1444-0.012=0.1324$; the lowest root is
\begin{align}
  r^* &= \frac{0.38-\sqrt{0.1324}}{0.2} = \frac{0.38-0.363868}{0.2} = 0.080659, \nonumber
\end{align}
with $\rho_0(r^*)=0.12+0.2(0.080659)=0.136132$ and, by \eqref{t5e13-drdA} with $f=1$, $dr^*/dA=-1/(0.5-0.136132)=-1/0.363868=-2.748249$.
Here $\hat r$ solves $0.12+0.2\hat r=0.25$, so $\hat r=(0.25-0.12)/0.2=0.65>r^*$.
Under the court $\mathcal P^c(x)=0.72+(0.25-0.12)x-0.1x^2=0.72+0.13x-0.1x^2$, so
$\mathcal P^c(\hat r)=0.72+0.0845-0.04225=0.76225\ge0.75$, and $\mathcal P^c(r^{**})=0.75$ gives $0.1x^2-0.13x+0.03=0$, discriminant
$0.0169-0.012=0.0049$, $r^{**}=(0.13-\sqrt{0.0049})/0.2=(0.13-0.07)/0.2=0.3$.
The borrower's loss is $\int_{r^*}^{r^{**}}(\rho_1-L)\,dr=\int_{r^*}^{r^{**}}(0.1+r)\,dr=0.1(0.3-0.080659)+(0.09-0.080659^2)/2=0.021934+0.041747=0.063681$.
\emph{Case (iv):} $L=0.3$, $I-A=0.78$: \eqref{t5e13-rstar} is $0.1x^2-0.18x+0.06=0$, discriminant $0.0324-0.024=0.0084$, so
$r^*=(0.18-0.091652)/0.2=0.441742$, with $\rho_0(r^*)=0.208348\le L$; $\hat r=(0.15-0.12)/0.2=0.15<r^*$. \eqref{t5e13-Pmax} gives
$0.5+\int_0^{0.15}0.15\,dr+\int_{0.15}^1(0.12+0.2r)\,dr=0.5+(0.15)(0.15)+0.12(0.85)+0.1(1-0.0225)=0.5+0.0225+0.102+0.09775=0.72225<0.78$: no financing.

\intuition Bankruptcy is the borrower's commitment device: handing $L$ to investors after bad news is what makes the loan pay. A court that
gives the borrower half of $L$ ex post destroys pledgeable income ex ante. Competitive lenders still break even, so the whole cost lands on the
borrower: more frequent (inefficient) bankruptcy, or no loan at all.

\pitfalls
\begin{itemize}
  \item Taking the highest root of the break-even equation: liquidation is inefficient by (1), so the
        borrower wants the \emph{lowest} cutoff that still lets investors break even.
  \item Concluding that the court helps borrowers or hurts lenders: lenders anticipate the split and
        break even (or do not lend); the whole cost is borne by the borrower.
  \item Thinking the default could be strategic here: $r$ is verifiable, so default happens only when
        $r<d$ and creditors liquidate because $\rho_0(r)<L$.
\end{itemize}

\bookans Matches the book's answer (pp.~586--587): same program, $r^*$ the lowest value satisfying breakeven, $d=r^*$, $dr^*/dA<0$, no financing when
$r^*>\hat r$, and when $r^*<\hat r$ either no financing or bankruptcy over a larger interval $[0,r^{**})$, $r^{**}>r^*$, hurting the borrower only.
(The book's wording ``no bankruptcy occurs on some interval $[0,r^{**}]$'' must mean bankruptcy occurs there.)

\refbib Tirole \S3.2.1 ``The Fixed-Investment Model'' and \S3.2.2 ``The Lenders' Credit Analysis'', pp.~115--117; Exercise 4.9, pp.~190--191; answer pp.~586--587.

% =============================================================================
% Original exercises 14--17 (sovereign debt; no book exercise exists; user agreed 2026-11-12)
% =============================================================================
\exsec{14. A two-period sovereign bond: default set, price and spread --- Original exercise}
%% BEGIN-STATEMENT 14
\stmtbegin
\textbf{Origin:} original exercise on sovereign default and bond pricing, built on Prof.~Ricca's
Topic 5 slides 13--18, 21 and 28--29 (Aula/Topic\_5.pdf) · Answer: in the solutions file

A small open economy lives for two periods, \(t = 0, 1\). Its government maximizes
\(u(c_0) + \beta E_0[u(c_1)]\), with \(\beta \in (0, 1)\), \(u' > 0\) and \(u'' < 0\). At \(t = 0\)
it has no assets (\(B = 0\)), observes output \(y_0 = 1\), and chooses the face value \(B'\) of the
one-period, non-state-contingent bonds it buys: a position \(B' < 0\) is a debt of face value
\(b \equiv -B' > 0\) due at \(t = 1\), and it raises \(-q(B')B'\) units of the good at \(t = 0\),
where \(q(B')\) is the unit price. At \(t = 1\) output is \(y_H = 1.2\) with probability \(0.8\) or
\(y_L = 0.8\) with probability \(0.2\). Period \(t = 1\) is the last one: nothing can be borrowed or
saved at \(t = 1\), and nothing happens afterwards. After observing \(y\), the government either
repays, consuming \(c_1 = y + B'\), or defaults, consuming \(h(y) = (1 - \lambda)y\) with
\(\lambda = 0.1\) (the output cost of default; exclusion from credit markets costs nothing more,
because the economy ends). Debt contracts cannot be enforced: the government defaults whenever
default gives strictly higher utility than repaying, and repays when indifferent. Foreign lenders
are risk-neutral and competitive, can invest at the risk-free rate \(r = 0.05\), and know the
distribution of \(y\).

\begin{enumerate}
\def\labelenumi{(\alph{enumi})}
\item
  Write the value of repaying, \(V^c(B', y)\), and the value of defaulting, \(V^d(y)\), at
  \(t = 1\). Show that the government defaults if and only if \(b > \lambda y\). Does the answer
  depend on the utility function \(u\)?
\item
  Find the default set \(D(B') = \{y \in \{y_L, y_H\} : V^d(y) > V^c(B', y)\}\) for every \(B'\),
  including \(B' \ge 0\) (the country lends).
\item
  Derive the default probability \(\delta(B')\) and the bond price \(q(B')\) implied by the lenders'
  zero-profit condition. Evaluate \(q\) at \(B' = 0.05\), \(-0.05\), \(-0.08\), \(-0.10\), \(-0.12\)
  and \(-0.15\).
\item
  For each \(B'\) in (c) with \(q(B') > 0\), compute the gross interest rate \(1 + r(B') = 1/q(B')\)
  and the spread \(r(B') - r\). At \(B' = -0.10\), check that a lender who buys the issue earns, in
  expectation, exactly the risk-free return.
\item
  The country raises \(-q(B')B'\) at \(t = 0\). (i) Find the bond position \(B^*\) that maximizes
  the amount raised, and that maximum amount. (ii) Show that any \(B'\) with
  \(-0.10 \le B' < -0.08\) is dominated: a risk-free bond raises the same amount at \(t = 0\) and
  costs the country no more at \(t = 1\) in either state (strictly less in at least one). (iii)
  Explain why the country never chooses \(B' < B^*\).
\end{enumerate}

\paragraph{Context needed to solve it}

\emph{From Prof.~Ricca's Topic 5 slides 13--18, 21 and 28--29 (Aula/Topic\_5.pdf), restated for this
exercise's two-period setting:}

\begin{itemize}
\tightlist
\item
  Small open economy with a non-storable, stochastic endowment \(\tilde y_t\); preferences
  \(E_0\sum_t \beta^t u(c_t)\) with \(\beta \in (0,1)\), \(u' > 0\), \(u'' < 0\) (slide 13). It
  smooths consumption by borrowing from or lending to competitive, risk-neutral foreign investors;
  bonds have a one-period maturity and are not state-contingent; debt contracts are not enforceable,
  so the country can renege even when it has the means to repay; default costs an output loss and
  exclusion from capital markets; the country defaults whenever this option is more valuable than
  repaying (slide 14).
\item
  A bond with face value \(B'\) costs \(q(B', y)B'\) today; the interest rate is
  \(1 + r(B', y) = B'/[q(B', y)B'] = 1/q(B', y)\) (slide 15). \(B' > 0\): the country pays \(qB'\)
  today and receives \(B'\) next period (it lends); \(B' < 0\): it receives \(-qB'\) today and pays
  \(-B'\) next period (it borrows) (slide 16).
\item
  When the country repays, its resource constraint is \(c = y + B - q(B', y)B'\); when it defaults
  it is excluded from foreign credit, re-enters with probability \(\theta\) each period, and
  consumes \(y^{def} = h(y) \le y\) while excluded (slide 17). In this exercise \(t = 1\) is the
  last period, so \(B' = 0\) is forced at \(t = 1\) and \(\theta\) plays no role.
\item
  Competition and risk neutrality: \(q(B', y)B'(1 + r) = B'[1 - \delta(B', y)]\),
  i.e.~\(q(B', y) = [1 - \delta(B', y)]/(1 + r)\); for \(B' > 0\) the default probability is zero
  (slide 18).
\item
  Default set \(D(B) = \{y \in Y : V^d(y) > V^c(B, y)\}\); the probability of default next period on
  debt \(B'\) is \(\delta(B', y) = \Pr[y' \in D(B') \mid y]\) (slide 21).
\item
  The amount raised is \(-q(B')B'\); the maximum increase in consumption is given by the \(B^*\)
  that solves \(\min_{B'} q(B')B'\), and it is never optimal to choose \(B' < B^*\) because the
  country can raise the same amount at a higher price, i.e.~with a smaller repayment next period
  (slides 28--29).
\end{itemize}
\stmtend
%% END-STATEMENT

\noindent\emph{Original exercise} (no course book has one on this model), built on Prof.\ Ricca's Topic~5 slides 13--18, 21 and 28--29.
Notation as on the slides: $B'$ are the \emph{assets} bought for $t=1$, so $B'<0$ is a debt of face value $b=-B'$.

\asked (a) Write the two values at the last date and find the rule that tells when the government defaults on a debt $b$.
(b) List, for every bond position, the output states in which it defaults. (c) Turn those sets into default probabilities and
the price a competitive lender pays per unit of face value, at six given positions. (d) Turn prices into interest rates and
spreads over the risk-free rate, and check the lender's break-even at $B'=-0.10$. (e) Find the position that raises the most
cash at $t=0$, show a band of risky positions is beaten by a risk-free bond, and explain why borrowing beyond $B^*$ never pays.

\approach At the last date there is no future, so $V^d$ and $V^c$ are single-period utilities and the comparison
reduces to consumption: default is a \emph{strategic} choice (nonverifiable repayment, slide 14), disciplined here only by the
output cost $h(y)\le y$. Lenders price the bond by zero profit, $q=(1-\delta)/(1+r)$ (slide 18), which turns the default set
(slide 21) into a step-shaped price schedule; the amount raised, $-q(B')B'$, then has a maximum $B^*$ (slides 28--29).

\notation
\begin{itemize}
  \item $B'$: assets bought at $t=0$ (face value due at $t=1$); $b=-B'$: debt; $y\in\{y_L,y_H\}=\{0.8,1.2\}$ with $\Pr(y_H)=0.8$, $\Pr(y_L)=0.2$.
  \item $h(y)=(1-\lambda)y$, $\lambda=0.1$: consumption after default; $\lambda y$: output lost by defaulting.
  \item $V^c(B',y)$, $V^d(y)$: values of repaying and of defaulting at $t=1$; $D(B')$: default set; $\delta(B')$: default probability.
  \item $r=0.05$: risk-free rate; $q(B')$: unit bond price; $1+r(B')=1/q(B')$: gross interest rate; $r(B')-r$: spread.
\end{itemize}

\step{1}{(a) The default rule} At $t=1$ nothing can be borrowed, so repaying leaves $c_1=y+B'$ and defaulting leaves $h(y)$:
\begin{align}
  V^c(B',y) &= u(y+B') = u(y-b), \label{t5e14-vc}\\
  V^d(y)    &= u\big(h(y)\big) = u\big((1-\lambda)y\big). \label{t5e14-vd}
\end{align}
Default happens iff $V^d(y)>V^c(B',y)$. Since $u$ is strictly increasing, $u(x_1)>u(x_2)\iff x_1>x_2$, so
\begin{align}
  V^d(y)>V^c(B',y) &\iff (1-\lambda)y > y-b \nonumber\\
                   &\iff b > y-(1-\lambda)y \nonumber\\
                   &\iff b > \lambda y. \label{t5e14-rule}
\end{align}

\answer{(a)} $V^c(B',y)=u(y+B')$, $V^d(y)=u((1-\lambda)y)$; default iff $b=-B'>\lambda y$ (the debt exceeds the output lost by
defaulting). The rule does not depend on $u$ (only on $u'>0$): at the last date default compares consumption levels.

\step{2}{(b) Default sets} The two thresholds of \eqref{t5e14-rule} are
\begin{align}
  \lambda y_L &= (0.1)(0.8) = 0.08, & \lambda y_H &= (0.1)(1.2) = 0.12. \nonumber
\end{align}
For $B'\ge0$ we have $b\le0<\lambda y$ in both states (the country is owed money and gains nothing by defaulting). Ties go to repayment:
at $b=0.08$, $y_L$, repaying leaves $0.8-0.08=0.72=(0.9)(0.8)$; at $b=0.12$, $y_H$, repaying leaves $1.2-0.12=1.08=(0.9)(1.2)$. Hence
\begin{align}
  D(B') = \begin{cases} \emptyset & B'\ge-0.08,\\ \{y_L\} & -0.12\le B'<-0.08,\\ \{y_L,y_H\} & B'<-0.12. \end{cases} \label{t5e14-D}
\end{align}

\answer{(b)} $D(B')=\emptyset$ for $B'\ge-0.08$ (all lending positions included); $D(B')=\{y_L\}$ for $-0.12\le B'<-0.08$; $D(B')=\{y_L,y_H\}$ for $B'<-0.12$.

\step{3}{(c) Default probability and price} $\delta(B')=\Pr[y\in D(B')]$ adds the probabilities of the default states in \eqref{t5e14-D}: $0$, $0.2$
or $0.2+0.8=1$. Zero profit for a lender buying face value $B'$ at $q(B')B'$ (slide 18) gives
\begin{align}
  q(B')B'(1+r) &= B'[1-\delta(B')] \nonumber\\
  q(B') &= \frac{1-\delta(B')}{1+r}
   = \begin{cases} \dfrac{1}{1.05} = 0.952381 & B'\ge-0.08,\\[6pt] \dfrac{1-0.2}{1.05}=\dfrac{0.8}{1.05} = 0.761905 & -0.12\le B'<-0.08,\\[6pt] \dfrac{1-1}{1.05}=0 & B'<-0.12. \end{cases} \label{t5e14-q}
\end{align}
At the six positions: $q(0.05)=q(-0.05)=q(-0.08)=0.952381$; $q(-0.10)=q(-0.12)=0.761905$; $q(-0.15)=0$.

\answer{(c)} $\delta(B')\in\{0,\,0.2,\,1\}$ on the three regions of \eqref{t5e14-D}; $q=0.952381$ at $0.05,-0.05,-0.08$; $q=0.761905$ at $-0.10,-0.12$; $q=0$ at $-0.15$.

\step{4}{(d) Interest rates and spreads} By slide 15, $1+r(B')=1/q(B')$. On the risk-free region,
\begin{align}
  1+r(B') &= \frac{1}{1/1.05} = 1.05, & r(B')-r &= 0.05-0.05 = 0. \nonumber
\end{align}
On the region where only $y_L$ defaults ($B'=-0.10,-0.12$),
\begin{align}
  1+r(B') &= \frac{1}{0.8/1.05} = \frac{1.05}{0.8} = 1.3125, & r(B')-r &= 0.3125-0.05 = 0.2625. \nonumber
\end{align}
At $B'=-0.15$, $q=0$: no lender pays anything for a bond that is never repaid, so no interest rate clears the market.
Lender check at $B'=-0.10$: the issue costs $q\,b=(0.761905)(0.10)=0.0761905$ and pays $0.10$ with probability $0.8$ and nothing
in default (the output loss $\lambda y_L$ is destroyed, not transferred), so
\begin{align}
  \frac{E[\text{payoff}]}{\text{price}} &= \frac{(0.8)(0.10)+(0.2)(0)}{0.0761905} = \frac{0.08}{0.0761905} = 1.05 = 1+r. \nonumber
\end{align}

\answer{(d)} Risk-free positions: $1+r(B')=1.05$, spread $0$. $B'=-0.10,-0.12$: $1+r(B')=1.3125$, spread $0.2625$ (26.25 percentage points).
$B'=-0.15$: no market. The lender earns exactly $1.05$ per unit at $B'=-0.10$.

\step{5}{(e)(i) Maximum amount raised} With $b=-B'$, the cash raised is $X(b)=-q(B')B'=q\,b$. By \eqref{t5e14-q}:
\begin{align}
  X(b) = \begin{cases} b/1.05, & 0<b\le0.08,\\ 0.8\,b/1.05, & 0.08<b\le0.12,\\ 0, & b>0.12. \end{cases} \nonumber
\end{align}
On the first step $X$ rises to $0.08/1.05=0.076190$; on the second it drops to just above $0.8(0.08)/1.05=0.060952$ and rises to $0.8(0.12)/1.05=0.091429$.
The largest value is at $b=0.12$: $B^*=-0.12$, raising $0.091429$.

\step{6}{(e)(ii) A dominated band} Take $b\in(0.08,0.10]$; it raises $0.8b/1.05$. A bond of face value $b_s=0.8b$ raises $b_s/1.05=0.8b/1.05$, the same
amount, and is risk-free because $b_s=0.8b\le0.8(0.10)=0.08$. What each costs at $t=1$:
\begin{align}
  y_H:&\quad \text{risky pays } b, \quad \text{safe pays } 0.8b < b; \nonumber\\
  y_L:&\quad \text{risky defaults and loses } \lambda y_L = 0.08, \quad \text{safe pays } 0.8b \le 0.08. \nonumber
\end{align}
So $t=1$ consumption is at least as high in both states with the safe bond, and strictly higher in $y_H$: every $B'\in[-0.10,-0.08)$ is dominated.
(For $b\in(0.10,0.12]$ the amount $0.8b/1.05>0.076190$ exceeds anything risk-free debt can raise, so no such comparison exists.)

\step{7}{(e)(iii) Never beyond $B^*$} For $B'<B^*=-0.12$ the price is $0$: the country raises nothing at $t=0$, and at $t=1$ defaults in both states, losing
$\lambda y$ of output. $B'=0$ gives the same $t=0$ consumption and no loss, so $B'<B^*$ is strictly dominated. Slide 29's general argument: any amount
raised with $B'<B^*$ can also be raised with some $B'\ge B^*$ at a higher price, hence with a smaller repayment.

\answer{(e)} (i) $B^*=-0.12$, maximum $0.091429$. (ii) For $B'\in[-0.10,-0.08)$ the risk-free bond $b_s=0.8b$ raises the same cash and costs $0.8b$ in both
states, against $b$ in $y_H$ and $0.08$ of output in $y_L$. (iii) Beyond $B^*$ the price is zero: nothing is raised and default costs output in every state.

\intuition A sovereign repays only because default destroys output: debt up to $\lambda y$ is credible in state $y$. Debt that exceeds the cost in
bad times but not in good times is repaid only in good times, so lenders charge a default premium. The price schedule is a staircase, and the cash
raised is a ``debt Laffer curve'': more face value raises more cash until the next default state is crossed, where cash drops. Default is pure
waste (lenders get nothing, the country loses $\lambda y$), which is why risky debt that a safe bond can replicate is never chosen.

\pitfalls
\begin{itemize}
  \item Sign of $B'$: on the slides $B'$ is \emph{assets}; debt is $B'<0$ and the cash raised is $-q(B')B'>0$.
  \item Writing the spread as $1/q-1=0.3125$: that is the risky rate $r(B')$; the spread subtracts $r=0.05$.
  \item Treating ties as defaults: $D(B)$ uses a strict inequality, so $b=\lambda y$ is repaid (this keeps $q(-0.12)>0$).
\end{itemize}

\refbib Prof.\ Ricca, Topic~5 slides 13--18 (setup, prices, interest rate, zero profit), 21 (default set), 28--29 (borrowing limit $B^*$). Original exercise; numbers in the check script.

% =============================================================================
\exsec{15. The price schedule $q(B',y)$ with persistent output --- Original exercise}
%% BEGIN-STATEMENT 15
\stmtbegin
\textbf{Origin:} original exercise on the sovereign bond price schedule, built on Prof.~Ricca's
Topic 5 slides 15, 17--18, 21, 23--24 and 27--29 (Aula/Topic\_5.pdf) · Answer: in the solutions file

A small open economy lives for two periods, \(t = 0, 1\), with preferences
\(u(c_0) + \beta E_0[u(c_1)]\), \(\beta \in (0,1)\), \(u' > 0\), \(u'' < 0\). Output takes three
values, \(y \in \{0.8, 1.0, 1.2\}\), and follows a Markov chain: given output \(y\) at \(t = 0\),
output \(y'\) at \(t = 1\) has the probabilities \(\pi(y' \mid y)\) below.

{\def\LTcaptype{none} % do not increment counter
\begin{longtable}[]{@{}llll@{}}
\toprule\noalign{}
current \(y\) ~next \(y'\) & \(0.8\) & \(1.0\) & \(1.2\) \\
\midrule\noalign{}
\endhead
\bottomrule\noalign{}
\endlastfoot
\(0.8\) & \(0.6\) & \(0.3\) & \(0.1\) \\
\(1.0\) & \(0.2\) & \(0.6\) & \(0.2\) \\
\(1.2\) & \(0.1\) & \(0.3\) & \(0.6\) \\
\end{longtable}
}

At \(t = 0\) the government has no assets, observes \(y\), and chooses the face value \(B'\) of
one-period, non-state-contingent bonds (\(B' < 0\) is a debt of face value \(b = -B'\) due at
\(t = 1\)), raising \(-q(B', y)B'\) at \(t = 0\). Period \(t = 1\) is the last one (no borrowing or
saving at \(t = 1\)). After observing \(y'\), the government either repays, consuming \(y' + B'\),
or defaults, consuming \(h(y') = \min\{y', \hat y\}\) with \(\hat y = 0.9\) (the output-cost form
used by Arellano 2008: default costs nothing when output is already below \(\hat y\) and costs
\(y' - \hat y\) when it is above). The government defaults whenever default gives strictly higher
utility, and repays when indifferent. Lenders are risk-neutral and competitive, can invest at
\(r = 0.04\), observe \(y\) and know the transition probabilities.

\begin{enumerate}
\def\labelenumi{(\alph{enumi})}
\item
  Let \(D(B') = \{y' : V^d(y') > V^c(B', y')\}\). Show that \(B'_2 \ge B'_1\) implies
  \(D(B'_2) \subseteq D(B'_1)\), and conclude that for every current \(y\) the price \(q(B', y)\) is
  nondecreasing in \(B'\) (nonincreasing in the debt \(b\)). Does \(D(B')\) depend on current output
  \(y\)? Does \(q(B', y)\)?
\item
  Find \(D(B')\) for every \(B'\).
\item
  Compute \(q(B', y)\) for \(B' \in \{0.10,\ 0,\ -0.05,\ -0.10,\ -0.20,\ -0.30,\ -0.40\}\) and
  \(y \in \{0.8, 1.0, 1.2\}\); present the results as a table.
\item
  At \(B' = -0.05\), compute the gross interest rate \(1 + r(B', y) = 1/q(B', y)\) and the spread
  \(r(B', y) - r\) for each current \(y\). Explain why the same bond pays a higher interest rate
  when current output is low. What would \(q(-0.05, y)\) be, for every current \(y\), if instead
  \(y'\) were i.i.d. with probabilities \((0.3, 0.4, 0.3)\) on \((0.8, 1.0, 1.2)\)?
\item
  For each current \(y\), find the bond position \(B^*(y) \in [-0.5, 0)\) that maximizes the amount
  raised, \(-q(B', y)B'\), and the maximum amount. Interpret the result as an endogenous borrowing
  limit.
\end{enumerate}

\paragraph{Context needed to solve it}

\emph{From Prof.~Ricca's Topic 5 slides 15, 17--18, 21, 23--24 and 27--29 (Aula/Topic\_5.pdf),
restated for this exercise's two-period setting:}

\begin{itemize}
\tightlist
\item
  A bond with face value \(B'\) costs \(q(B', y)B'\); the interest rate is
  \(1 + r(B', y) = 1/q(B', y)\) (slide 15). Repaying, the country consumes
  \(c = y + B - q(B', y)B'\); defaulting, it consumes \(y^{def} = h(y) \le y\) while excluded (slide
  17).
\item
  Lenders' zero profit: \(q(B', y) = [1 - \delta(B', y)]/(1 + r)\), with \(\delta(B', y) = 0\) for
  \(B' > 0\) (slide 18). Default set \(D(B) = \{y \in Y : V^d(y) > V^c(B, y)\}\);
  \(\delta(B', y) = \Pr[y' \in D(B') \mid y] = \sum_{y' \in D(B')} \pi(y' \mid y)\) (slide 21, with
  the integral over the density \(f(y', y)\) written as a sum for a finite state space).
\item
  Claim 1: for given \(y\), \(V^c(B, y)\) is increasing in \(B\). Claim 2: if \(B_2 \ge B_1\) then
  \(D(B_2) \subseteq D(B_1)\), since \(V^c\) increases with \(B\) while \(V^d(y)\) stays constant.
  Hence the equilibrium price function \(q(B', y)\) is increasing in \(B'\) and the interest rate
  decreasing (slides 23--24).
\item
  In the i.i.d. case the price function \(q(B') = [1 - F(y^*(B'))]/(1 + r)\) does not depend on
  current \(y\); equilibrium prices still depend on \(y\) because the chosen \(B'\) does (slide 27).
\item
  The maximum increase in consumption in a period is given by the \(B^*\) that solves
  \(\min_{B'} q(B')B'\); this endogenous borrowing limit means that, after a bad shock, the country
  would like to borrow more but such a contract is not available; choosing \(B' < B^*\) is never
  optimal (slides 28--29).
\end{itemize}
\stmtend
%% END-STATEMENT

\noindent\emph{Original exercise}, built on Prof.\ Ricca's Topic~5 slides 15, 17--18, 21, 23--24 and 27--29.

\asked (a) Prove, in this two-period economy, that more assets (less debt) shrink the default set, so the bond price falls with debt for any current
output; say what depends on current output. (b) Find the default set at every position. (c) Fill the price table for seven positions and three current
outputs. (d) Compute interest rates and spreads on a small debt for each current output, explain the recession premium, and price the same bond under
i.i.d.\ output. (e) For each current output, find the debt that raises the most cash.

\approach Claims 1--2 (slides 23--24) in their simplest form: at the last date $V^c=u(y'+B')$ rises with $B'$ while $V^d=u(h(y'))$ does not, so default
sets are nested. Prices follow from zero profit with Markov probabilities, $q=(1-\sum_{y'\in D(B')}\pi(y'\mid y))/(1+r)$ (slides 18, 21): the default
\emph{set} depends only on tomorrow's state, the default \emph{probability} on today's. The Arellano cost $h(y)=\min\{y,\hat y\}$ makes default cheap in
recessions.

\notation
\begin{itemize}
  \item $y$: output at $t=0$ (current); $y'$: output at $t=1$; $\pi(y'\mid y)$: transition probabilities (table in the statement).
  \item $h(y')=\min\{y',0.9\}$; $\kappa(y')\equiv y'-h(y')$: output lost by defaulting in state $y'$.
  \item $B'$, $b=-B'$, $D(B')$, $\delta(B',y)$, $q(B',y)$, $1+r(B',y)=1/q(B',y)$, $r=0.04$; $B^*(y)$: cash-maximising position at current $y$.
\end{itemize}

\step{1}{(a) Nested default sets and a monotone price} At $t=1$ (last date) $V^c(B',y')=u(y'+B')$ and $V^d(y')=u(h(y'))$. Let $B'_2\ge B'_1$ and
$y'\in D(B'_2)$. Then
\begin{align}
  u\big(h(y')\big) &> u(y'+B'_2) && (y'\in D(B'_2)) \nonumber\\
                   &\ge u(y'+B'_1) && (B'_2\ge B'_1,\ u'>0), \nonumber
\end{align}
so $y'\in D(B'_1)$: $D(B'_2)\subseteq D(B'_1)$. Since probabilities are nonnegative, summing over a smaller set gives a smaller number:
\begin{align}
  \delta(B'_2,y) = \sum_{y'\in D(B'_2)}\pi(y'\mid y) &\le \sum_{y'\in D(B'_1)}\pi(y'\mid y) = \delta(B'_1,y) \nonumber\\
  q(B'_2,y) = \frac{1-\delta(B'_2,y)}{1+r} &\ge \frac{1-\delta(B'_1,y)}{1+r} = q(B'_1,y). \label{t5e15-mono}
\end{align}

\answer{(a)} $D(B'_2)\subseteq D(B'_1)$ for $B'_2\ge B'_1$, hence $q(B',y)$ is nondecreasing in $B'$ (nonincreasing in debt) for every $y$. $D(B')$ depends only
on $B'$ and $y'$, not on current $y$; $\delta(B',y)$ and $q(B',y)$ depend on $y$ through $\pi(\cdot\mid y)$.

\step{2}{(b) Default sets} As in Exercise 14, default iff $h(y')>y'-b$, i.e.\ iff $b>\kappa(y')$. The costs are
\begin{align}
  \kappa(0.8) &= 0.8-\min\{0.8,0.9\} = 0, & \kappa(1.0) &= 1.0-0.9 = 0.1, & \kappa(1.2) &= 1.2-0.9 = 0.3. \nonumber
\end{align}
With ties going to repayment,
\begin{align}
  D(B') = \begin{cases} \emptyset & B'\ge0,\\ \{0.8\} & -0.1\le B'<0,\\ \{0.8,\,1.0\} & -0.3\le B'<-0.1,\\ \{0.8,\,1.0,\,1.2\} & B'<-0.3. \end{cases} \label{t5e15-D}
\end{align}

\answer{(b)} \eqref{t5e15-D}: any debt is defaulted on in the recession state $0.8$ (default costs nothing there); debt above $0.1$ also in state $1.0$; above $0.3$ always.

\step{3}{(c) The price table} $\delta(B',y)$ adds the row-$y$ probabilities of the states in \eqref{t5e15-D}, and $q=(1-\delta)/1.04$.
\begin{align}
  -0.1\le B'<0:&\quad \delta = \pi(0.8\mid y) = 0.6,\ 0.2,\ 0.1, \nonumber\\
               &\quad q = \tfrac{0.4}{1.04},\ \tfrac{0.8}{1.04},\ \tfrac{0.9}{1.04} = 0.384615,\ 0.769231,\ 0.865385; \nonumber\\
  -0.3\le B'<-0.1:&\quad \delta = \pi(0.8\mid y)+\pi(1.0\mid y) = 0.9,\ 0.8,\ 0.4, \nonumber\\
                  &\quad q = \tfrac{0.1}{1.04},\ \tfrac{0.2}{1.04},\ \tfrac{0.6}{1.04} = 0.096154,\ 0.192308,\ 0.576923; \nonumber
\end{align}
for $y=0.8,1.0,1.2$ respectively; $q=1/1.04=0.961538$ for $B'\ge0$ and $q=0$ for $B'<-0.3$.

\begin{center}
\begin{tabular}{lccccccc}
\toprule
$B'$ & $0.10$ & $0$ & $-0.05$ & $-0.10$ & $-0.20$ & $-0.30$ & $-0.40$\\
\midrule
$q(B',0.8)$ & $0.961538$ & $0.961538$ & $0.384615$ & $0.384615$ & $0.096154$ & $0.096154$ & $0$\\
$q(B',1.0)$ & $0.961538$ & $0.961538$ & $0.769231$ & $0.769231$ & $0.192308$ & $0.192308$ & $0$\\
$q(B',1.2)$ & $0.961538$ & $0.961538$ & $0.865385$ & $0.865385$ & $0.576923$ & $0.576923$ & $0$\\
\bottomrule
\end{tabular}
\end{center}

\answer{(c)} The table above: each row is a nonincreasing staircase in debt, and each column (for $B'<0$) rises with current $y$.

\step{4}{(d) Spreads and the recession premium} At $B'=-0.05$, $1+r(B',y)=1/q=1.04/(1-\delta)$:
\begin{align}
  y=0.8:&\quad 1+r(B',y)=\tfrac{1.04}{0.4}=2.6, & r(B',y)-r&=1.6-0.04=1.56; \nonumber\\
  y=1.0:&\quad 1+r(B',y)=\tfrac{1.04}{0.8}=1.3, & r(B',y)-r&=0.3-0.04=0.26; \nonumber\\
  y=1.2:&\quad 1+r(B',y)=\tfrac{1.04}{0.9}=1.155556, & r(B',y)-r&=0.155556-0.04=0.115556. \nonumber
\end{align}
The bond is defaulted on exactly when $y'=0.8$; with persistent output a recession today makes a recession tomorrow likelier ($0.6$ against $0.2$ or
$0.1$), so lenders demand more. With i.i.d.\ output $\delta(-0.05)=\Pr(y'=0.8)=0.3$ whatever $y$ is, and
\begin{align}
  q(-0.05,y) &= \frac{1-0.3}{1.04} = \frac{0.7}{1.04} = 0.673077 \quad\text{for every } y, \nonumber
\end{align}
as slide 27 says: the i.i.d.\ price function does not depend on current output.

\answer{(d)} Spreads at $B'=-0.05$: $1.56$ ($y=0.8$), $0.26$ ($y=1.0$), $0.115556$ ($y=1.2$). Persistence raises the default probability after a bad
shock. Under i.i.d.\ $(0.3,0.4,0.3)$: $q(-0.05,y)=0.673077$ for all $y$.

\step{5}{(e) Borrowing limit by current state} Within each price step, $q(B',y)b$ rises with $b$, so the maximum is at a right end of a step,
$b=0.1$ or $b=0.3$ (beyond $0.3$ the price is $0$):
\begin{align}
  y=0.8:&\quad (0.1)(0.384615)=0.038462 \;>\; (0.3)(0.096154)=0.028846 &&\Rightarrow B^*(0.8)=-0.1; \nonumber\\
  y=1.0:&\quad (0.1)(0.769231)=0.076923 \;>\; (0.3)(0.192308)=0.057692 &&\Rightarrow B^*(1.0)=-0.1; \nonumber\\
  y=1.2:&\quad (0.1)(0.865385)=0.086538 \;<\; (0.3)(0.576923)=0.173077 &&\Rightarrow B^*(1.2)=-0.3. \nonumber
\end{align}

\answer{(e)} $B^*(0.8)=-0.1$ (raises $0.038462$), $B^*(1.0)=-0.1$ ($0.076923$), $B^*(1.2)=-0.3$ ($0.173077$). The most the country can raise is
smallest in the recession, exactly when it would like to borrow to smooth consumption: an endogenous, procyclical borrowing limit (slide 28).

\intuition The default decision belongs to tomorrow's state, the price to today's information about it. Output costs that are small in recessions
(Arellano's $\min\{y,\hat y\}$) put the default states at the bottom of the distribution, and persistence makes a bad state today a signal of
default tomorrow. So credit is dearest and scarcest when the country needs it most --- the countercyclical spreads of emerging markets.

\pitfalls
\begin{itemize}
  \item Letting $D(B')$ depend on current $y$: the default decision is taken at $t=1$ and depends on $y'$ only; current $y$ enters through $\pi(\cdot\mid y)$.
  \item Using the stationary or unconditional probabilities with Markov output: lenders condition on current $y$.
  \item Looking for the maximum of $q\,b$ by calculus: $q$ jumps, so the candidates are the right ends of the steps.
\end{itemize}

\refbib Prof.\ Ricca, Topic~5 slides 15, 17--18, 21 (prices, default set), 23--24 (Claims 1--2, monotone $q$), 27 (i.i.d.\ price), 28--29 (borrowing limit). Original exercise.

% =============================================================================
\exsec{16. No capital inflows when default is possible (Claim 3) --- Original exercise}
%% BEGIN-STATEMENT 16
\stmtbegin
\textbf{Origin:} original exercise on the i.i.d. case with permanent exclusion of the
Eaton--Gersovitz model, built on Prof.~Ricca's Topic 5 slides 17--22 and 25 (Aula/Topic\_5.pdf) ·
Answer: in the solutions file

Consider the infinite-horizon small open economy with preferences
\(E_0\sum_{t=0}^{\infty}\beta^t u(c_t)\), \(\beta \in (0,1)\), \(u' > 0\), \(u'' < 0\). Output
\(y_t\) is i.i.d. over time. Each period the government enters with assets \(B\) (\(B < 0\) is
debt), observes \(y\), and either repays --- choosing next period's assets \(B' \ge Z\) (\(Z < 0\)
is a no-Ponzi bound) and consuming \(c = y + B - q(B')B'\) --- or defaults. After a default the
country is excluded from credit markets and re-enters each period with probability \(\theta\); while
excluded it consumes \(h(y) \le y\). Lenders are risk-neutral and competitive, with risk-free rate
\(r > 0\), so \(q(B') = [1 - \delta(B')]/(1 + r)\), where \(\delta(B') = \Pr[\tilde y \in D(B')]\)
and \(D(B) = \{y : V^d(y) > V^c(B, y)\}\). The value functions are
\[V^o(B, y) = \max\{V^d(y), V^c(B, y)\},\]
\[V^d(y) = u(h(y)) + \beta E\big[\theta V^o(0, \tilde y') + (1 - \theta)V^d(\tilde y')\big],\]
\[V^c(B, y) = \max_{B' \ge Z}\Big\{u\big(y + B - q(B')B'\big) + \beta E\big[V^o(B', \tilde y')\big]\Big\}.\]
In parts (a)--(d), \(\theta = 0\) (permanent exclusion) and \(h(y) = y\) (no output cost).

\begin{enumerate}
\def\labelenumi{(\alph{enumi})}
\item
  Show that \(V^d(y) = u(y) + \dfrac{\beta}{1 - \beta}E[u(\tilde y)]\). Evaluate \(V^d\) in both
  states when \(u(c) = \ln c\), \(y \in \{0.9, 1.1\}\) with probability \(1/2\) each, and
  \(\beta = 0.9\).
\item
  Prove: if, for some \(B\), the default set is non-empty, \(D(B) \ne \emptyset\), then there is no
  bond position \(\hat B\) available with \(B - q(\hat B)\hat B > 0\) (no capital inflow at \(B\)).
\item
  A commentator claims that, with the numbers in (a), a country holding \(B = -0.2\) defaults when
  \(y = 0.9\) even though lenders offer to buy bonds \(\hat B = -0.25\) at the risk-free price
  \(q(\hat B) = 1/(1 + r)\), with \(r = 0.05\). Compute the consumption and the net capital inflow
  that taking this offer would give, and show --- using only \(V^o \ge V^d\) --- that default at
  \(y = 0.9\) cannot then be optimal.
\item
  Show that \(D(0) = \emptyset\): a country with no debt never defaults. Does the result in (b) say
  that the country can never borrow?
\item
  Which assumptions does the proof in (b) use? Show that the result still holds (i) for any output
  cost \(h(y) \le y\), with \(\theta = 0\), and (ii) with Markov output, where prices are
  \(q(\hat B, y)\).
\item
  Now let \(\theta > 0\). Show at which step the proof in (b) breaks, and explain the economics.
\end{enumerate}

\paragraph{Context needed to solve it}

\emph{From Prof.~Ricca's Topic 5 slides 17--22 and 25 (Aula/Topic\_5.pdf):}

\begin{itemize}
\tightlist
\item
  Repaying, the country's resource constraint is \(c = y + B - q(B', y)B'\). Defaulting, it loses
  access to foreign credit for a stochastic number of periods, regains access next period with
  probability \(\theta\), and consumes \(y^{def} = h(y) \le y\) until it re-enters (slide 17).
  Lenders: \(q(B', y) = [1 - \delta(B', y)]/(1 + r)\), with zero default probability when \(B' > 0\)
  (slide 18).
\item
  \(V^o(B, y) = \max\{V^d(B, y), V^c(B, y)\}\) (slide 19);
  \(V^d(y) = u(y^{def}) + \beta\int_{y'}[\theta V^o(0, y') + (1 - \theta)V^d(y')]f(y', y)\,dy'\) and
  \(V^c(B, y) = \max_{B' \ge Z}\{u(y + B - q(B', y)B') + \beta\int_{y'}V^o(B', y')f(y', y)\,dy'\}\),
  where \(Z < 0\) is a lower bound on debt that prevents Ponzi schemes (slide 20). In the i.i.d.
  case the integrals are expectations over \(\tilde y'\) and \(q\) does not depend on current \(y\).
\item
  Default is optimal if \(V^d(y) > V^c(B, y)\); \(D(B) = \{y \in Y : V^d(y) > V^c(B, y)\}\);
  \(\delta(B', y) = \Pr[y' \in D(B') \mid y]\) (slide 21). An equilibrium is a set of policy
  functions \(c(s)\), \(B'(s)\), default sets \(D(B)\) and prices \(q(B', y)\) such that consumption
  satisfies the resource constraint, choices are optimal given prices, and prices reflect default
  probabilities with zero profits (slide 22).
\item
  Slide 25 states the result of part (b) for the i.i.d. case with \(\theta = 0\) and \(h(y) = y\)
  (``Claim 3'') and adds: if the country could roll over the debt, that is, borrow an amount large
  enough to repay the previous debt, it would choose to do so, given that it can still exercise the
  option to default in the next period.
\end{itemize}
\stmtend
%% END-STATEMENT

\noindent\emph{Original exercise}, built on Prof.\ Ricca's Topic~5 slides 17--22 and 25.

\asked (a) Solve the autarky value with permanent exclusion and no output cost, and evaluate it in a two-state log example. (b) Prove slide 25's Claim 3: if the
country would default in some state at debt $B$, no bond on offer can give it a net inflow at $B$. (c) Apply the argument to a concrete offer. (d) Show a debt-free
country never defaults, and say what Claim 3 does \emph{not} say. (e) Find which assumptions the proof needs and extend it to output costs and Markov output.
(f) Show where re-entry breaks the proof.

\approach Revealed preference on the Bellman equations of slides 19--20: $V^c$ is a maximum, so any feasible bond gives a lower bound, and $V^o=\max\{V^d,V^c\}\ge V^d$
everywhere. With permanent exclusion the continuation of default is $V^d$ itself, so ``roll over now, keep the option to default later'' is always at least as good
tomorrow as defaulting today (slide 25's remark); a positive inflow makes it strictly better today.

\notation
\begin{itemize}
  \item $B$: assets entering the period ($B<0$ debt); $\hat B$: a candidate bond position; $B-q(\hat B)\hat B$: net capital inflow (consumption minus output).
  \item $V^o,V^c,V^d$: value with the option, of repaying, of defaulting; $D(B)$: default set; $\theta$: re-entry probability; $h(y)\le y$: output after default.
  \item $E[\cdot]$: expectation over next period's i.i.d.\ output $\tilde y'$; numerical example: $u=\ln$, $y\in\{0.9,1.1\}$ equally likely, $\beta=0.9$, $r=0.05$.
\end{itemize}

\step{1}{(a) Autarky value} With $\theta=0$ and $h(y)=y$, $V^d(y)=u(y)+\beta E[V^d(\tilde y')]$. With i.i.d.\ output, $m\equiv E[V^d(\tilde y')]$ is the same number
whatever today's $y$; taking expectations of both sides,
\begin{align}
  m &= E[u(\tilde y)] + \beta m \nonumber\\
  m &= \frac{E[u(\tilde y)]}{1-\beta}, \nonumber\\
  V^d(y) &= u(y) + \beta m = u(y) + \frac{\beta}{1-\beta}E[u(\tilde y)]. \label{t5e16-vd}
\end{align}
Numbers: $\ln0.9=-0.1053605$, $\ln1.1=0.0953102$, so $E[\ln\tilde y]=\tfrac12(-0.1053605+0.0953102)=-0.0050252$ and, with $\beta/(1-\beta)=0.9/0.1=9$,
$9E[\ln\tilde y]=-0.0452265$ (products from unrounded logarithms):
\begin{align}
  V^d(0.9) &= -0.1053605-0.0452265 = -0.1505870, & V^d(1.1) &= 0.0953102-0.0452265 = 0.0500837. \nonumber
\end{align}

\answer{(a)} $V^d(y)=u(y)+\frac{\beta}{1-\beta}E[u(\tilde y)]$; $V^d(0.9)=-0.150587$, $V^d(1.1)=0.050084$.

\step{2}{(b) Proof of Claim 3} Suppose $D(B)\neq\emptyset$; pick $y\in D(B)$, and suppose, to reach a contradiction, that some $\hat B\ge Z$ has $B-q(\hat B)\hat B>0$. Three facts:
\begin{align}
  V^d(y) &> V^c(B,y) && (y\in D(B)) \label{t5e16-f1}\\
  V^c(B,y) &\ge u\big(y+B-q(\hat B)\hat B\big)+\beta E\big[V^o(\hat B,\tilde y')\big] && (\hat B \text{ feasible; } V^c \text{ is a max}) \label{t5e16-f2}\\
  V^d(y) &= u(y)+\beta E\big[V^d(\tilde y')\big] && (\theta=0,\ h(y)=y) \label{t5e16-f3}
\end{align}
Substituting \eqref{t5e16-f3} and \eqref{t5e16-f2} into \eqref{t5e16-f1},
\begin{align}
  u(y)+\beta E\big[V^d(\tilde y')\big] &> u\big(y+B-q(\hat B)\hat B\big)+\beta E\big[V^o(\hat B,\tilde y')\big] \nonumber\\
  u(y)-u\big(y+B-q(\hat B)\hat B\big) &> \beta E\big[V^o(\hat B,\tilde y')-V^d(\tilde y')\big] \;\ge\; 0, \label{t5e16-key}
\end{align}
the last step because $V^o(\hat B,y')=\max\{V^d(y'),V^c(\hat B,y')\}\ge V^d(y')$ in every state. So $u(y)>u(y+B-q(\hat B)\hat B)$, and since $u$ is increasing,
$y>y+B-q(\hat B)\hat B$, i.e.\ $B-q(\hat B)\hat B<0$ --- contradicting the inflow. With i.i.d.\ output the menu $\{q(\hat B),\hat B\}$ is the same in every state, so
no contract with an inflow is available at $B$ at all.

\answer{(b)} If $D(B)\neq\emptyset$, then $B-q(\hat B)\hat B\le0$ for every available $\hat B$: at debt levels where default is possible, the country can only make net repayments.

\step{3}{(c) The commentator's example} Taking the offer at $B=-0.2$, $y=0.9$:
\begin{align}
  c &= y+B-q(\hat B)\hat B = 0.9-0.2-\tfrac{1}{1.05}(-0.25) = 0.7+0.238095 = 0.938095, \nonumber\\
  B-q(\hat B)\hat B &= -0.2+0.238095 = 0.038095 > 0. \nonumber
\end{align}
By \eqref{t5e16-f2}, $V^o\ge V^d$ and \eqref{t5e16-f3},
\begin{align}
  V^c(-0.2,0.9) &\ge \ln0.938095+\beta E\big[V^o(-0.25,\tilde y')\big] \nonumber\\
                &\ge \ln0.938095+\beta E\big[V^d(\tilde y')\big] \nonumber\\
                &= V^d(0.9)+\ln\frac{0.938095}{0.9} = V^d(0.9)+0.041457 > V^d(0.9). \nonumber
\end{align}
So $0.9\notin D(-0.2)$; the same steps at $y=1.1$ give $1.1\notin D(-0.2)$. If such an offer exists, $D(-0.2)=\emptyset$: the claim is inconsistent.

\answer{(c)} $c=0.938095$, inflow $0.038095>0$; repaying and rolling over beats default by at least $0.041457$ utils, so default at $y=0.9$ is not optimal.

\step{4}{(d) A debt-free country never defaults} At $B=0$ the choice $B'=0$ is feasible and gives $c=y$. Then
\begin{align}
  V^c(0,y) &\ge u(y)+\beta E\big[V^o(0,\tilde y')\big] && (B'=0 \text{ feasible}) \nonumber\\
           &\ge u(y)+\beta E\big[V^d(\tilde y')\big] && (V^o\ge V^d) \nonumber\\
           &= V^d(y), && \text{by \eqref{t5e16-f3}} \nonumber
\end{align}
so $V^d(y)>V^c(0,y)$ never holds: $D(0)=\emptyset$. Claim 3 therefore does not apply at $B=0$, and a debt-free country can borrow ($-q(B')B'>0$ for any $B'<0$ with $q(B')>0$).
Claim 3 says something narrower: once debt is high enough that default is possible in some state, the market offers only net repayment --- default happens when the
country ``cannot roll over its debt'' (slides 25, 27). How much debt is sustained comes only from the value of keeping access to smoothing: with the numbers of (a),
a grid solution (check script; Exercise 17(e)) finds a risk-free limit of only $0.01$, one percent of mean output.

\answer{(d)} $V^c(0,y)\ge V^d(y)$, so $D(0)=\emptyset$. No: Claim 3 restricts net inflows only at debt levels where default is possible; positive (if small) borrowing is sustained.

\step{5}{(e) Which assumptions matter} The proof used: $\theta=0$ in \eqref{t5e16-f3} (the continuation of default is $V^d$, which is below $V^o$); $h(y)=y$ in \eqref{t5e16-f3}; and
i.i.d.\ output only to make the menu of contracts the same in every state.
(i) With $h(y)\le y$ and $\theta=0$, \eqref{t5e16-f3} becomes $V^d(y)=u(h(y))+\beta E[V^d(\tilde y')]$ and \eqref{t5e16-key} becomes
\begin{align}
  u\big(h(y)\big)-u\big(y+B-q(\hat B)\hat B\big) &> \beta E\big[V^o(\hat B,\tilde y')-V^d(\tilde y')\big] \ge 0 \nonumber\\
  \Rightarrow\quad h(y) &> y+B-q(\hat B)\hat B \nonumber\\
  \Rightarrow\quad B-q(\hat B)\hat B &< -\big[y-h(y)\big] \le 0, \nonumber
\end{align}
which is even stronger: in a default state, the net repayment demanded exceeds the output cost $y-h(y)$.
(ii) With Markov output, replace $E$ by $E[\cdot\mid y]$ and $q(\hat B)$ by $q(\hat B,y)$; every step holds state by state, so for each $y\in D(B)$ no contract offered
\emph{at that $y$} has an inflow (the menu may differ across states).

\answer{(e)} Only $\theta=0$ is essential. The result holds for any $h(y)\le y$ (with the stronger bound $B-q\hat B<-(y-h(y))$) and, state by state, with Markov output.

\step{6}{(f) Re-entry breaks the proof} With $\theta>0$, $V^d(y)=u(h(y))+\beta E[\theta V^o(0,\tilde y')+(1-\theta)V^d(\tilde y')]$, and the right side of \eqref{t5e16-key} becomes
\begin{align}
  &\beta E\big[V^o(\hat B,\tilde y')-\theta V^o(0,\tilde y')-(1-\theta)V^d(\tilde y')\big] \nonumber\\
  &\quad= \beta\theta E\big[V^o(\hat B,\tilde y')-V^o(0,\tilde y')\big] + \beta(1-\theta)E\big[V^o(\hat B,\tilde y')-V^d(\tilde y')\big]. \nonumber
\end{align}
The second term is still $\ge0$, but the first is $\le0$ for $\hat B<0$: $V^c$ rises with $B$ (Claim 1, slide 23) and $V^d$ does not depend on $B$, so $V^o=\max\{V^d,V^c\}$
is nondecreasing in $B$ and $V^o(\hat B,\cdot)\le V^o(0,\cdot)$. The right side can be negative, so a small inflow no longer contradicts default.

\answer{(f)} The step ``right side $\ge0$'' fails: with re-entry, defaulting can lead back to the market with $B=0$, worth more than returning with debt $\hat B$.

\intuition Rolling over debt is a way of postponing the default decision; with permanent exclusion, postponing is free (the option is kept) and an inflow makes it
strictly better today. So a country defaults only when markets demand net repayment. Re-entry turns default into a fresh start: it wipes out the debt, which
rolling over does not, and that can be worth more than a small inflow today.

\pitfalls
\begin{itemize}
  \item Reading Claim 3 as ``no borrowing'': it concerns debt levels with $D(B)\neq\emptyset$; at $B=0$ the default set is empty and the country can borrow.
  \item Writing \eqref{t5e16-f2} as an equality: $\hat B$ need not be optimal; the proof only needs the lower bound.
  \item Thinking output costs break the result: they strengthen it; only re-entry ($\theta>0$) does.
\end{itemize}

\refbib Prof.\ Ricca, Topic~5 slides 17--22 (setup, value functions, default set, equilibrium), 23 (Claim 1), 25 (Claim 3 and the roll-over remark). Original exercise.

% =============================================================================
\exsec{17. Infinite horizon with log utility: certifying risk-free debt --- Original exercise}
%% BEGIN-STATEMENT 17
\stmtbegin
\textbf{Origin:} original exercise on the risk-free region and the borrowing capacity of the
Eaton--Gersovitz model, built on Prof.~Ricca's Topic 5 slides 17--24 and 29 (Aula/Topic\_5.pdf) ·
Answer: in the solutions file

Consider the infinite-horizon small open economy of slides 17--22: preferences
\(E_0\sum_{t=0}^{\infty}\beta^t\ln c_t\) with \(\beta = 0.9\); output i.i.d.,
\(y \in \{y_L, y_H\} = \{0.9, 1.1\}\) with probability \(1/2\) each; one-period non-state-contingent
bonds priced by risk-neutral competitive lenders at \(q(B') = [1 - \delta(B')]/(1 + r)\) with
\(r = 0.05\); no enforcement. A defaulting country is excluded forever (\(\theta = 0\)) and consumes
\(h(y) = (1 - \lambda)y\) every period, with \(\lambda = 0.02\). The value functions are
\[V^o(B, y) = \max\{V^d(y), V^c(B, y)\},\qquad V^d(y) = \ln h(y) + \beta E[V^d(\tilde y')],\]
\[V^c(B, y) = \max_{B' \ge Z}\Big\{\ln\big(y + B - q(B')B'\big) + \beta E\big[V^o(B', \tilde y')\big]\Big\},\]
with \(Z < 0\) a no-Ponzi bound, and \(D(B) = \{y : V^d(y) > V^c(B, y)\}\). Fix a debt level
\(b > 0\) (assets \(B = -b\)) and consider the conjecture that bonds with face value \(b\) are
risk-free: \(q(-b) = 1/(1 + r)\).

\begin{enumerate}
\def\labelenumi{(\alph{enumi})}
\item
  Compute \(V^d(y_L)\) and \(V^d(y_H)\).
\item
  Under the conjecture, show that servicing the debt forever --- repaying and choosing \(B' = -b\)
  every period --- gives consumption \(y - s\) each period, with \(s \equiv rb/(1 + r)\), and that
  \[V^c(-b, y) \ge W(b, y) \equiv \ln(y - s) + \frac{\beta}{1 - \beta}E[\ln(\tilde y - s)].\]
  Conclude that if \(W(b, y) \ge V^d(y)\) in both states, then \(D(-b) = \emptyset\), so the
  conjecture \(q(-b) = 1/(1 + r)\) confirms itself.
\item
  Let \(\Delta(s, y) \equiv W(b, y) - V^d(y)\). Show that \(\Delta\) is increasing in \(y\), so the
  condition of (b) only needs to be checked at \(y_L\). Show also that with no output cost
  (\(\lambda = 0\)) \(\Delta(s, y) < 0\) for every \(s > 0\), and relate this to Exercise 16.
\item
  Find the largest \(\bar s\) with \(\Delta(\bar s, y_L) \ge 0\): write the equation that defines
  it, apply Newton's method from \(s_0 = 0\) until six decimals are stable, and verify the sign of
  \(\Delta(s, y_L)\) at \(s = 0.019\) and \(s = 0.020\). Convert \(\bar s\) into the certified debt
  level \(\bar b = \bar s(1 + r)/r\) and compare it with mean output.
\item
  A grid solution of the full model (slides 19--22: value iteration on an asset grid of step
  \(0.00375\) on \([-0.6, 0.3]\), with prices updated from the default sets until they converge)
  gives the lowest asset level with an empty default set, \(\bar B = \inf\{B : D(B) = \emptyset\}\):
  \(\bar B = -0.42\) for \((\lambda, \theta) = (0.02, 0)\); \(\bar B = -0.0375\) for
  \((\lambda, \theta) = (0.02, 0.5)\); and, on a grid of step \(0.0005\) on \([-0.05, 0.05]\),
  \(\bar B = -0.01\) for \((\lambda, \theta) = (0, 0)\). (i) Explain why \(\bar b\) from (d) is a
  lower bound on \(-\bar B\), and check it against the grid value. (ii) Compute \(d\bar s/d\lambda\)
  and \(d\bar b/d\lambda\) at \(\lambda = 0.02\) by implicit differentiation. (iii) Use the grid
  results to explain how the output cost \(\lambda\) and the re-entry probability \(\theta\) shape
  the country's borrowing capacity.
\end{enumerate}

\paragraph{Context needed to solve it}

\emph{From Prof.~Ricca's Topic 5 slides 17--24 and 29 (Aula/Topic\_5.pdf):}

\begin{itemize}
\tightlist
\item
  Repaying: \(c = y + B - q(B', y)B'\). Defaulting: exclusion from foreign credit, re-entry next
  period with probability \(\theta\), consumption \(y^{def} = h(y) \le y\) until re-entry (slide
  17). Lenders: \(q(B', y) = [1 - \delta(B', y)]/(1 + r)\) (slide 18).
\item
  \(V^o(B, y) = \max\{V^d(B, y), V^c(B, y)\}\) (slide 19);
  \(V^d(y) = u(y^{def}) + \beta\int_{y'}[\theta V^o(0, y') + (1 - \theta)V^d(y')]f(y', y)\,dy'\);
  \(V^c(B, y) = \max_{B' \ge Z}\{u(y + B - q(B', y)B') + \beta\int_{y'}V^o(B', y')f(y', y)\,dy'\}\)
  with \(Z < 0\) a no-Ponzi bound (slide 20); \(D(B) = \{y \in Y : V^d(y) > V^c(B, y)\}\),
  \(\delta(B', y) = \Pr[y' \in D(B') \mid y]\) (slide 21); equilibrium: resource constraint, optimal
  choices given prices, zero-profit prices (slide 22).
\item
  Claim 1: \(V^c(y, B)\) is increasing in \(B\); Claim 2:
  \(B_2 \ge B_1 \Rightarrow D(B_2) \subseteq D(B_1)\) (slide 23).
  \(\underline B = \sup\{B : D(B) = Y\}\) and \(\bar B = \inf\{B : D(B) = \emptyset\}\), with
  \(\underline B \le \bar B \le 0\) (slide 24). When \(B > \bar B\), debt is risk-free and
  \(q(B) = 1/(1 + r)\) (slide 29).
\item
  Newton's method for \(G(s) = 0\): \(s_{k+1} = s_k - G(s_k)/G'(s_k)\). Implicit differentiation of
  \(G(\bar s, \lambda) = 0\): \(d\bar s/d\lambda = -G_\lambda/G_s\).
\end{itemize}

\emph{From Tirole §3.8 ``Nonverifiable Income'', ``Relation to the sovereign debt literature'',
pp.~143--144 (PDF 156--157), for part (e)(iii):}

\begin{quote}
Repayment of debt by the sovereign responds to two incentives: international sanctions and the
future cost of being shut down from the international capital market after default. A subliterature,
starting with Bulow and Rogoff (1989a,b), assumes away sanctions and focuses on the incentives
provided by exclusion. {[}\ldots{]} Bulow and Rogoff consider an infinite-horizon,
symmetric-information model in which (a) the sovereign can decide not to reimburse and (b) the
sovereign can save, and (c) the rate of growth of the economy is smaller than the rate of interest.
They show that no lending is feasible as the borrower always prefers to default (and save some of
the concomitant extra income). {[}\ldots{]} Second, an outright exclusion, in which the defaulting
sovereign cannot even save, makes it particularly costly for the sovereign to repudiate its debt.
Again, some sovereign debt may then be issued in equilibrium (Kehoe and Levine 1993; Kocherlakota
1996).
\end{quote}
\stmtend
%% END-STATEMENT

\noindent\emph{Original exercise}, built on Prof.\ Ricca's Topic~5 slides 17--24 and 29, with Tirole \S3.8, pp.~143--144, for part (e)(iii).

\asked (a) Compute the autarky values with a 2\% output cost. (b) Show that, if a debt is priced as risk-free, the value of repaying is at least the value of
servicing the debt forever, and that this makes a risk-free price self-confirming. (c) Show the test binds in the low state and fails for every debt without
output costs. (d) Solve the test's threshold by Newton's method and convert it into a debt level. (e) Compare with a grid solution of the full model, derive
the threshold's sensitivity to the cost, and explain how the cost and re-entry drive borrowing capacity.

\approach A full equilibrium (slides 19--22) needs a computer, but a lower bound on $V^c$ is available by hand: $V^c$ is a maximum, so any feasible plan
--- here ``repay and keep the same debt forever'' --- bounds it from below. If that plan beats autarky in every state, default never happens at that debt,
the debt is risk-free (slide 29: $q=1/(1+r)$ above $\bar B$), and the risk-free price used to compute the plan's value is confirmed. The threshold is a
one-dimensional equation in $s$, solved by Newton's method; implicit differentiation gives its comparative statics.

\notation
\begin{itemize}
  \item $y\in\{y_L,y_H\}=\{0.9,1.1\}$ i.i.d., probability $\tfrac12$ each; $\beta=0.9$, $K\equiv\beta/(1-\beta)=9$; $r=0.05$; $h(y)=(1-\lambda)y$, $\lambda=0.02$; $\theta=0$.
  \item $b>0$: debt (assets $B=-b$); $s\equiv rb/(1+r)$: per-period net debt service when $b$ is rolled over at the risk-free price.
  \item $W(b,y)$: value of servicing $b$ forever; $\Delta(s,y)=W-V^d$; $G(s)\equiv\Delta(s,y_L)$; $\bar s$, $\bar b=\bar s(1+r)/r$: certified threshold.
  \item $\bar B=\inf\{B:D(B)=\emptyset\}$: true risk-free limit (slide 24).
\end{itemize}

\step{1}{(a) Autarky values} Exactly as in \eqref{t5e16-vd}, with $h(y)$ in place of $y$: $V^d(y)=\ln[(1-\lambda)y]+K\,E\{\ln[(1-\lambda)\tilde y]\}$. With $(0.98)(0.9)=0.882$ and
$(0.98)(1.1)=1.078$: $\ln0.882=-0.1255632$, $\ln1.078=0.0751075$, $E=\tfrac12(-0.1255632+0.0751075)=-0.0252279$, $9E=-0.2270509$ (from unrounded logarithms), so
\begin{align}
  V^d(0.9) &= -0.1255632-0.2270509 = -0.3526141, \nonumber\\
  V^d(1.1) &= 0.0751075-0.2270509 = -0.1519434. \nonumber
\end{align}

\answer{(a)} $V^d(y_L)=-0.352614$, $V^d(y_H)=-0.151943$.

\step{2}{(b) Servicing the debt forever} Enter with $B=-b$, repay, and buy $B'=-b$ at $q(-b)=1/(1+r)$:
\begin{align}
  c &= y+B-q(B')B' = y-b+\frac{b}{1+r} = y-\frac{(1+r)b-b}{1+r} = y-\frac{rb}{1+r} = y-s. \nonumber
\end{align}
Since this choice is feasible and $V^o\ge V^c$,
\begin{align}
  V^c(-b,y) &\ge \ln(y-s)+\beta E\big[V^o(-b,\tilde y')\big] \;\ge\; \ln(y-s)+\beta E\big[V^c(-b,\tilde y')\big]. \label{t5e17-iter}
\end{align}
Apply \eqref{t5e17-iter} to $V^c(-b,\tilde y')$ inside the expectation, and again $T$ times (i.i.d.: each future period has the same distribution):
\begin{align}
  V^c(-b,y) &\ge \ln(y-s)+\sum_{t=1}^{T}\beta^t E[\ln(\tilde y-s)] + \beta^{T+1}E\big[V^c(-b,\tilde y_{T+1})\big]. \nonumber
\end{align}
Values are bounded (consumption lies in a bounded set), so the last term vanishes as $T\to\infty$, and $\sum_{t\ge1}\beta^t=\beta/(1-\beta)$:
\begin{align}
  V^c(-b,y) &\ge \ln(y-s)+\frac{\beta}{1-\beta}E[\ln(\tilde y-s)] = W(b,y). \label{t5e17-W}
\end{align}
If $W(b,y)\ge V^d(y)$ for $y=y_L,y_H$, then $V^c(-b,y)\ge V^d(y)$, so $V^d>V^c$ never holds: $D(-b)=\emptyset$, $\delta(-b)=0$, and slide 18 gives $q(-b)=1/(1+r)$ --- the price assumed.

\answer{(b)} $c=y-rb/(1+r)$; $V^c(-b,y)\ge W(b,y)$; $W\ge V^d$ in both states $\Rightarrow D(-b)=\emptyset\Rightarrow q(-b)=1/(1+r)$: the conjecture confirms itself.

\step{3}{(c) Only $y_L$ matters; nothing is certified without output costs} Subtract $V^d$ term by term, writing $\ln(y-s)-\ln[(1-\lambda)y]=\ln(1-s/y)-\ln(1-\lambda)$:
\begin{align}
  \Delta(s,y) &= \ln\!\Big(1-\frac sy\Big)-\ln(1-\lambda) + K\,E\Big[\ln\!\Big(1-\frac s{\tilde y}\Big)-\ln(1-\lambda)\Big]. \label{t5e17-Delta}
\end{align}
Only the first term depends on current $y$, and
\begin{align}
  \frac{\partial}{\partial y}\ln\!\Big(1-\frac sy\Big) &= \frac{s/y^2}{1-s/y} = \frac{s}{y(y-s)} > 0 \quad (0<s<y), \nonumber
\end{align}
so $\Delta(s,y_H)>\Delta(s,y_L)$: if the test passes at $y_L$ it passes at $y_H$. With $\lambda=0$ the $\ln(1-\lambda)$ terms vanish and, for $s>0$, $\ln(1-s/y)<0$ in every
state, so $\Delta(s,y)<0$: servicing any debt forever is worse than autarky. Link to Exercise 16: with $h(y)=y$, a constant net outflow $s$ buys nothing that autarky
lacks; whatever debt is sustained must rest on the value of smoothing (varying $B'$ with $y$), which this bound ignores.

\answer{(c)} $\partial\Delta/\partial y=s/[y(y-s)]>0$, so test only at $y_L$. With $\lambda=0$, $\Delta(s,y)<0$ for all $s>0$: the bound certifies no debt.

\step{4}{(d) The threshold by Newton's method} At $y_L$, \eqref{t5e17-Delta} with $K=9$ and probabilities $\tfrac12$ collects to $(1+\tfrac K2)\ln(1-s/0.9)+\tfrac K2\ln(1-s/1.1)-(1+K)\ln(0.98)$:
\begin{align}
  G(s) &= 5.5\ln\!\Big(1-\frac{s}{0.9}\Big)+4.5\ln\!\Big(1-\frac{s}{1.1}\Big)-10\ln0.98, \label{t5e17-G}\\
  G'(s) &= -\frac{5.5}{0.9-s}-\frac{4.5}{1.1-s} < 0. \nonumber
\end{align}
$G$ is decreasing, so $\{s:G(s)\ge0\}=[0,\bar s]$ with $G(\bar s)=0$. From $s_0=0$: $G(0)=-10\ln0.98=-10(-0.0202027)=0.202027$ and
$G'(0)=-5.5/0.9-4.5/1.1=-6.111111-4.090909=-10.202020$, so
\begin{align}
  s_1 &= 0-\frac{0.202027}{-10.202020} = 0.019803, \nonumber\\
  s_2 &= s_1-\frac{G(s_1)}{G'(s_1)} = 0.019803-\frac{-0.002089}{-10.414504} = 0.019803-0.000201 = 0.019602, \nonumber\\
  s_3 &= s_2-\frac{G(s_2)}{G'(s_2)} = 0.019602-\frac{-0.0000002}{-10.412307} = 0.019602. \nonumber
\end{align}
Six decimals are stable: $\bar s=0.019602$. Signs: $G(0.019)=5.5(-0.0213371)+4.5(-0.0174236)+0.202027=0.006266>0$ and $G(0.020)=-0.004145<0$, bracketing $\bar s$. Then
\begin{align}
  \bar b &= \bar s\,\frac{1+r}{r} = 0.019602\times\frac{1.05}{0.05} = 0.019602\times21 = 0.411642, \nonumber
\end{align}
about 41\% of mean output ($E[y]=1$).

\answer{(d)} $\bar s=0.019602$ (Newton: $0.019803\to0.019602\to0.019602$); $G(0.019)>0>G(0.020)$; $\bar b=0.411642\approx41\%$ of mean output.

\step{5}{(e)(i) A lower bound} \eqref{t5e17-W} is an inequality: the country can do better than constant service (borrow more after $y_L$, repay after $y_H$), so
$V^c(-b,y)$ may exceed $V^d(y)$ even when $W<V^d$. Every $b\le\bar b$ passes the test, so $D(-b)=\emptyset$ for all $b\le\bar b$ and $\bar B\le-\bar b$, i.e.\ $-\bar B\ge\bar b$.
Grid: $-\bar B=0.42\ge0.411642$, a gap of $0.42-0.411642=0.008358$: the bound is within 2\% of the grid value.

\step{6}{(e)(ii) Sensitivity to the output cost} Write \eqref{t5e17-G} with general $\lambda$: $G(s,\lambda)=5.5\ln(1-s/0.9)+4.5\ln(1-s/1.1)-10\ln(1-\lambda)$. Then
$G_\lambda=10/(1-\lambda)=10/0.98=10.204082$ and $G_s=G'(\bar s)=-10.412307$, so
\begin{align}
  \frac{d\bar s}{d\lambda} &= -\frac{G_\lambda}{G_s} = \frac{10.204082}{10.412307} = 0.980002, &
  \frac{d\bar b}{d\lambda} &= 21\times0.980002 = 20.580043. \nonumber
\end{align}
One more percentage point of output cost raises the certified debt by about $0.21$, a fifth of mean output.

\step{7}{(e)(iii) Output costs and re-entry} Output cost: moving $\lambda$ from $0$ to $0.02$ moves the grid limit $-\bar B$ from $0.01$ to $0.42$. Exclusion alone
sustains almost nothing here, because losing access to one-period non-contingent bonds costs little when shocks are $\pm10\%$; a permanent output loss is a large,
certain punishment and sustains most of the debt. Re-entry: moving $\theta$ from $0$ to $0.5$ (with $\lambda=0.02$) moves $-\bar B$ from $0.42$ to $0.0375$. With
$\theta>0$, $V^d$ puts weight $\theta$ on $V^o(0,\cdot)\ge V^d(\cdot)$, so (holding $V^o(0,\cdot)$ fixed) $V^d$ rises with $\theta$; the punishment lasts $1/\theta=2$ periods on
average instead of forever, default sets grow (Exercise 16(f)), and risk-free capacity collapses. This matches Tirole's summary (\S3.8, pp.~143--144): if a defaulter
could also save, Bulow--Rogoff show no lending is feasible; outright exclusion, under which the defaulter cannot save (autarky here), is what makes some debt
possible --- and costs of default beyond exclusion make it larger.

\answer{(e)} (i) $\bar b=0.411642\le-\bar B=0.42$: the bound holds and is nearly tight. (ii) $d\bar s/d\lambda=0.980002$, $d\bar b/d\lambda=20.580043$. (iii) Output costs
raise capacity ($0.01\to0.42$ as $\lambda$ goes $0\to0.02$); re-entry lowers it ($0.42\to0.0375$ as $\theta$ goes $0\to0.5$), because default becomes cheaper.

\intuition A sovereign's debt is worth only what it would lose by walking away. Here that loss has two parts: the output cost, paid every period of autarky, and the
lost ability to smooth. ``Service the debt forever'' prices the first part exactly and ignores the second, which is why it is a lower bound --- and a close one, since
smoothing is worth little with small i.i.d.\ shocks. Anything that shortens or softens the punishment (re-entry, saving abroad) shrinks the debt the country can credibly owe.

\pitfalls
\begin{itemize}
  \item Using $s=rb$: rolling over raises only $b/(1+r)$ against $b$ repaid, so the net service is $b-b/(1+r)=rb/(1+r)$.
  \item Testing only $y_H$, or averaging: the default incentive is strongest at $y_L$ (concavity), so that is where the test binds.
  \item Reading $\bar b$ as the exact risk-free limit: it is a sufficient condition; the full model's limit is weakly larger.
\end{itemize}

\refbib Prof.\ Ricca, Topic~5 slides 17--22 (model), 23--24 (Claims 1--2, $\underline B$, $\bar B$), 29 (risk-free region). Tirole \S3.8 ``Nonverifiable Income'', pp.~143--144 (PDF 156--157). Original exercise; grid solution in the check script (\texttt{e17solve}).

\end{document}
